\documentclass[11pt,a4paper]{article}
\pdfinfoomitdate=1
\usepackage[T1]{fontenc}
\usepackage[utf8]{inputenc}
\usepackage[margin=28mm]{geometry}
\usepackage{amsmath,amssymb,amsthm}
\usepackage{mathtools}
% amsthm supplies theorem styles; pf2 supplies every proof environment.
% Release the three conflicting amsthm definitions before loading pf2.
\let\proof\relax
\let\endproof\relax
\let\qed\relax
% Obtain pf2.sty from https://lamport.azurewebsites.net/latex/latex.html
% and put it beside this file or on the TeX input path.
\IfFileExists{pf2.sty}{\usepackage{pf2}}{%
  \PackageError{bennett-model}{pf2.sty is required}{%
    Download pf2.sty from Leslie Lamport's LaTeX packages page:\MessageBreak
    https://lamport.azurewebsites.net/latex/latex.html\MessageBreak
    Put the file beside bennett-model.tex or on the TeX input path.}%
}
\pfshortnumbers
\pfshortindent
\usepackage{microtype}
\usepackage[hidelinks,unicode]{hyperref}
\urlstyle{same}

\newtheorem{theorem}{Theorem}[section]
\newtheorem{lemma}[theorem]{Lemma}
\newtheorem{proposition}[theorem]{Proposition}
\newtheorem{corollary}[theorem]{Corollary}
\theoremstyle{definition}
\newtheorem{definition}[theorem]{Definition}
\newtheorem{example}[theorem]{Example}
\newtheorem{counterexample}[theorem]{Counterexample}
\newtheorem{assumption}[theorem]{Assumption}
\theoremstyle{remark}
\newtheorem{remark}[theorem]{Remark}
\newtheorem{convention}[theorem]{Convention}
\newtheorem{problem}[theorem]{Open problem}

\newcommand{\ket}[1]{\lvert #1\rangle}
\newcommand{\bra}[1]{\langle #1\rvert}
\newcommand{\proj}[1]{\ket{#1}\bra{#1}}
\DeclareMathOperator{\tr}{Tr}
\newcommand{\id}{\mathbb{I}}
\newcommand{\reg}[1]{\mathsf{#1}}
\newcommand{\cH}{\mathcal{H}}
\newcommand{\cL}{\mathcal{L}}
\newcommand{\C}{\mathbb{C}}
\newcommand{\N}{\mathbb{N}}
\newcommand{\eps}{\varepsilon}

\newcommand{\notedsection}[3]{%
  \section[#1]{#1\protect\footnotemark}\label{#2}\footnotetext{#3}}
\newcommand{\earlierforms}{the earlier \texttt{mark} and \texttt{rewind} forms}

\title{Mark, rewind and keep in the quantum circuit model:\\
  definitions, equivalence, soundness, cleanliness and resources}
\author{Working draft}
\date{}

\begin{document}
\maketitle
\begin{abstract}
We define classical control (a deterministic Turing machine), finite
reversible-Turing-machine encodings, and quantum circuits with fresh
$\ket0$ allocation and drop by partial trace, stating every result first
for the retained isometry on all qubits. For devices presented as explicit
total classical circuits with finite capacities and an input-independent
time bound, the nine basic classical and quantum instructions are
equivalent to the circuit model in the sense of equality of final
instruments (not of output or stopping times), with explicit overheads;
exact equivalence fails for arbitrary rotation angles, and approximation
is proved under stated effective-data hypotheses. Mark, keep and rewind
with pairings fixed in the syntax then denote total unitaries
$U^\dagger KU$ at every finite nesting depth, with fixed total completions
and no allocation, drop, observed measurement, external console observation
or actual stop inside a marked body; whether
the auxiliary register is thereby restored is a separate question, for
which we prove exact tests, a sufficient condition, and a four-qubit
counterexample to universal restoration. The prescribed recursive
expansion doubles work per nesting level, tightly for that expansion
(these are not lower bounds for optimal equivalent circuits); a
restricted delayed-erasure construction and an exact test for deleting
inner inverses are proved. For a bounded dynamic interpreter with all
registers retained, a supplied certificate gives unitary soundness and
exact undo under whole-sector matching encodings and routing obligations.
We also prove a static specialization, an exact collective pure-factorization
test for disposal from a fixed pure initial interpreter state, a timing-record
counterexample and conditional per-tick circuit costs. Compiling general
data-dependent mark/rewind into such a certified interpreter remains open.
\end{abstract}

\noindent\textbf{Status of this paper.}
Sections~\ref{sec:circuit-model}, \ref{sec:machines} and \ref{sec:primitives}
describe the foundations of today's device language: ordinary code,
allocation, gates, measurement and release by partial trace.
Today's \texttt{apply} computes an ordinary function for each definite
value of its arguments, copies the result into a target register by XOR
and undoes the computation.
Sections~\ref{sec:intro} and \ref{sec:instructions} introduce that
compute, copy, undo idea in the earlier design.
Sections~\ref{sec:instructions} and \ref{sec:soundness}--\ref{sec:tierB}
study \earlierforms, and Sections~\ref{sec:intro}, \ref{sec:open-problems}
and \ref{sec:summary} also discuss them.
The instructions were removed on 17 November 2026. The theorems for the
earlier forms still hold, and the forms can be recovered from commit
\texttt{f81a70a}. The paper's open problems in Section~\ref{sec:open-problems}
stand.

\tableofcontents

\notedsection{Introduction}{sec:intro}
  {Part of this section introduces \earlierforms.}
We study a computation described using exactly three machines: a deterministic Turing machine that supplies classical control, a reversible Turing machine whose finite encoded states can be used coherently, and a quantum circuit with fresh zero qubits and partial trace. The second machine is a convenient way to describe reversible classical operations on quantum registers; it adds no primitive physical operation. Every finite instance specifies its register sizes and running-time bound before receiving the unknown quantum input (the bound is part of the supplied data; no termination bound is inferred from the instruction list). We first retain all auxiliary qubits and describe an isometry, and only afterward take partial traces to obtain the observable instrument.

The central instruction pair has a direct operator meaning. If a delimited computation implements a total unitary $U$, and a reversible copying operation $K$ writes a specified Boolean function of its machine register into designated quantum targets by XOR, the complete operation is $U^\dagger KU$. A mark identifies the start of the delimited computation; rewind applies the inverse circuit after this copying operation, called keep. These terms do not imply that the machine register is restored. In particular, a gate changing a previously read qubit can prevent restoration even when the keep target was untouched by the forward computation.

Primitive~10 has the controlled-reversible meaning of
Lemma~\ref{lem:read-product} when the subsequent operations preserve the
read register's computational-basis sectors. A gate changing that register
removes this guarantee; the full joint unitary and its retained records
then specify the semantics.

Theorems~\ref{eqv:forward} and \ref{eqv:reverse} make the equivalence of the nine basic classical and quantum instructions precise for bounded uniform families and final instruments, with explicit overheads. Theorem~\ref{thm:sA-sound} proves soundness for any finite nesting of statically paired marks, allowing quantum gates between each mark and rewind. Theorem~\ref{thm:sA-clean} gives necessary and sufficient tests for
restoration on the whole initialised input space and characterises when
the reduced operation is nevertheless unitary; Theorem~\ref{thm:sA-recursive}
gives the corresponding tests on specified reached subspaces and at nested
reset points. Theorem~\ref{thm:sA-classical} gives a sufficient condition---product
output with computational-basis machine factor and a forward operation
independent of the keep target---which in particular covers reversible
classical computations; Theorem~\ref{thm:sA-four} gives an explicit four-qubit failure of universal restoration. Finite-alphabet exact equivalence excludes arbitrary ideal angles; Theorems~\ref{eqv:approx} and \ref{eqv:rotations} state the separate approximation guarantees and their effective-data assumptions.

Theorems~\ref{thm:resA-resources}--\ref{thm:resA-deletion} prove prescribed-expansion resource counts, a restricted delayed-erasure comparison and the exact inner-inverse deletion criterion; Counterexample~\ref{ex:resA-quantum} gives its five-qubit limitation. The tight exponential count concerns the prescribed expansion alone, with no lower bound for optimal equivalent circuits. Allocation, drop, observed measurement, external console observation and actual stop inside marked computations are excluded from the structured soundness theorem. No pebble-game trade-off continuum is proved. The full list of hypotheses and limitations is collected in Section~\ref{sec:summary}; Section~\ref{sec:open-problems} records the unresolved questions. These are mathematical statements about the model defined here, rather than a correctness claim for a particular software executor.

For dynamic control, Theorem~\ref{thm:tB-sound} proves a conditional result. A supplied certificate must verify the interpreter's local rules, orthogonal rule domains and orthogonal rule ranges, retained records, matching forward and undo encodings, and every reached state, including after keep. Under these hypotheses a fixed, bounded run is unitary on all retained logical registers, and undo with its specified routing is the exact adjoint of the recorded completed operation. A supplied clean implementation of the whole tick restores its workspace at each tick; this does not imply reset of the retained logical records. Theorem~\ref{thm:tB-A} includes Tier~A as a static clock schedule with its original machine registers retained; this specialization constructs no dynamic certificate.

Theorem~\ref{thm:tB-clean} gives the exact synchronization and disposal criterion: from a fixed pure initial machine state and for every data input, all retained interpreter information must finish in one fixed pure state independent of the data if tracing it is to leave a unitary data channel. Equal completion times alone do not ensure this. In Counterexample~\ref{ex:tB-delay}, undo restores the target at two different completion ticks, but the remaining timing record distinguishes the data values; disposing of it decoheres a superposition. Theorem~\ref{thm:tB-cost} charges each tick for scanning the instruction table, journal and mark stack, plus the supplied Boolean, quantum-action and clock circuits; a bounded run uses exactly $T$ tick circuits.

Here and below, \emph{source} refers to a hypothetical program text using
the instruction names listed in Section~\ref{sec:primitives}, before translation
into the terms defined here; no such text is given a general translation or
semantics in this paper.

Compiling general dynamic source mark/rewind into an interpreter satisfying that certificate remains open. In particular, no general translation for data-dependent repeated visits or an inverse stopped at a coherent mark identifier is proved (Open problem~\ref{prob:tB-compiler}). The conditional interpreter results do not establish such a source compiler or correctness of the existing executor.

Every proof uses Lamport's hierarchical notation~\cite[Section~3, PDF pp.~9--11]{Lamport2012}.
To typeset the proofs, obtain the original \texttt{pf2.sty} from
\href{https://lamport.azurewebsites.net/latex/latex.html}{Lamport's LaTeX packages page}.

\section{Quantum circuit model}\label{sec:circuit-model}
\begin{definition}[Types and observable equality]\label{def:registers}
A register list $L$ is a finite ordered list of distinct wire names, with
$\cH_L=(\C^2)^{\otimes |L|}$. Subregisters are ordered subsets; juxtaposition
of disjoint lists denotes tensor product. Products of operators act from
right to left. A width-$w$ value is a binary string, also interpreted as
an unsigned integer in $\{0,\ldots,2^w-1\}$; $\oplus$ is bitwise XOR.
A density operator is a positive operator of trace one. A channel is a
completely positive trace-preserving linear map. An instrument is a finite
family of completely positive trace-nonincreasing maps summing to a channel.
For channels with the same input and output, define
\[
 d_{\rm out}(\Phi,\Psi)=
 \sup_{R,\rho}\|((\Phi-\Psi)\otimes\id_R)(\rho)\|_1,
\]
where $R$ ranges over finite reference registers and $\rho$ over density
operators on the input and $R$; $\|A\|_1=\tr\sqrt{A^\dagger A}$.
An instrument is compared as a channel retaining its orthogonal outcome
labels. The operator norm is $\|A\|=\sup_{\|v\|=1}\|Av\|$.
All bounds and circuit descriptions are independent of the unknown input.
\end{definition}

\begin{definition}[Machine 3 and fresh names]\label{def:circuit}
Machine 3 is the finite quantum circuit model with labelled gates
\[
 H=2^{-1/2}\begin{pmatrix}1&1\\1&-1\end{pmatrix},\quad
 T=\operatorname{diag}(1,e^{i\pi/4}),\quad
 \operatorname{CNOT}\ket{a,b}=\ket{a,b\oplus a}.
\]
Write $\mathcal G=\{H,T,\operatorname{CNOT}\}$,
$\mathcal G^{\pm}=\mathcal G\cup\{T^\dagger\}$,
$Z=T^4$, $X=HZH$, and $S=T^2$.
Exactly two additional register transformations are allowed:
\[
 A_B:\cH_L\longrightarrow\cH_{LB},\quad
 A_B\ket\psi=\ket\psi\ket{0^{|B|}},\qquad
 D_B(\rho)=\tr_B\rho\quad(B\subseteq L).
\]
The names in $B$ in $A_B$ have never occurred earlier in the description.
Drop removes names from the accessible list permanently; later allocations
use new names. The circuit data designate an accessible output list $O$.
In the retained interpretation, a dropped factor is moved to an inaccessible
list $E$ and kept, with no subsequent operation on it. Thus the denotation
is first an isometry $V:\cH_{L_{\rm in}}\to\cH_{OE}$ and then the channel
$\rho\mapsto\tr_E(V\rho V^\dagger)$.
An optional ideal alphabet adds the actual matrices
$R_z(\theta)=\operatorname{diag}(e^{-i\theta/2},e^{i\theta/2})$.
An angle is supplied by a finite description together with a procedure
that, on input $b$, outputs dyadic approximations to the four complex
entries of $R_z(\theta)$ within $2^{-b}$ (the entry procedure of
Theorem~\ref{eqv:rotations}, whose cost is charged there). A description
without such a procedure does not define an ideal gate in this paper.
Their controlled versions use these matrices including their phases;
$T$ is not identified with $R_z(\pi/4)$ under a control.

Terminal computational-basis observation is part of the circuit model,
not an additional register transformation. Before choosing $O$, let $A$
be the final accessible list. The data may designate an ordered list
$(i_1,\ldots,i_b)$ of wires in $A$ to measure (repetitions allowed).
For $y\in\{0,1\}^b$ set
$P_y=\prod_{j=1}^b(\proj{y_j}_{i_j}\otimes\id)$.
Let $E_0$ be the list of previously dropped wires, disjoint from $A$.
If $V_0:\cH_{L_{\rm in}}\to\cH_{AE_0}$ is the retained circuit
isometry before observation, the
instrument on $O\subseteq A$ is
\[
 \rho\longmapsto\tr_{(AE_0)\setminus O}
       (P_yV_0\rho V_0^\dagger P_y),\qquad y\in\{0,1\}^b,
\]
where projectors act as identity on $E_0$. Its retained form first
copies each $i_j$ into
a fresh zero wire $Y_j$ by CNOT, then copies $Y$ into a fresh zero
$E_Y$. Selecting $OY$ and tracing its complement gives the displayed
instrument with orthogonal classical labels $y$. Without observation
the earlier output isometry notation is unchanged; with observation
it has output $OY$ and the complementary retained environment.
\end{definition}


\section{Classical and reversible Turing machines}\label{sec:machines}
\subsection{The classical control computer and reversible computation}

\begin{definition}[Machines 1 and 2; uniform descriptions]\label{def:rtm}
Machine 1 is a deterministic multitape Turing machine with finite control
and alphabet and partial transition function, with designated stopping
states. It generates descriptions and handles genuinely observed classical
data. A family is uniform when one such machine on $1^n$ and classical
parameters outputs the complete finite instance and all its bounds.
For polynomial uniformity its output time and the numerical bounds are
polynomial in $n$ and the lengths of classical inputs. Resource parameters
used in compilation estimates are encoded in unary; wire addresses and
Boolean circuit descriptions are binary. A time bound supplied in binary (a horizon, in the sense of
Definitions~\ref{def:read} and \ref{def:adaptive}) still has the same
numerical execution cost, which may be exponential in its encoding length. Approximation requests are positive rational numbers.

Machine 2 is a deterministic reversible multitape Turing machine:
commands on each tape either read and replace a symbol or move the head
by $-1,0,1$ without reading or writing; distinct commands have disjoint
domains and disjoint ranges. Its global transition $\delta$ is partial
and injective. This is Bennett's quadruple convention
\cite[pp.~526--527, definition preceding the main theorem]{Bennett1973}.
For a finite encoding choose, on each bilateral tape, the window
$\{0,\ldots,S_0-1\}$, with $S_0\geq1$, and heads in that window.
For $k$ tapes, alphabet $\Gamma$, and control set $F$, use
\[
 m_{\rm RTM}=kS_0\lceil\log_2|\Gamma|\rceil+
 k\lceil\log_2 S_0\rceil+\lceil\log_2|F|\rceil
\]
bits, with padding if needed. Invalid codes and transitions leaving the
window are excluded from the domain $D$ of the restricted $\delta$.
On $\{0,1\}^{m_{\rm RTM}}\setminus D$ fix the lexicographic bijection
to the complement of $\delta(D)$; together with $\delta$ this defines
$\widehat\delta$. Its permutation matrix is the specified completion.
On excluded codes this is only a completion, not a fault-reporting TM
simulation. Any assertion of agreement with the original TM requires
every applied forward transition to remain in $D$. No such promise is
needed for unitarity or for applying the completed inverse.
The simulation in Definition~\ref{def:device} below uses explicit circuits
instead of this potentially expensive completion. There is no autonomous
cyclic scheduler in the present definition.
\end{definition}

\begin{lemma}[Finite completion]\label{lem:completion}
The completion $\widehat\delta$ is a permutation and hence defines a
unitary on its entire finite register.
\end{lemma}
\begin{proof}
 \step{r2comp-size}{$|D|=|\delta(D)|$.}
 \begin{proof}\pf\ The restriction is injective with image $\delta(D)$.
 \end{proof}
 \step{r2comp-complement}{The ordered complements have equal size and
 their prescribed correspondence is bijective.}
 \begin{proof}\pf\ Subtract the equality in \stepref{r2comp-size} from
 the size of the common ambient finite set.\end{proof}
 \qedstep
 \begin{proof}\pf\ The domains partition the set, and their images
 partition it into two disjoint sets, each mapped bijectively. A
 permutation takes an orthonormal basis to an orthonormal basis.\end{proof}
\end{proof}

\begin{definition}[Clean implementation workspace]\label{def:scratch}
An implementation of a logical unitary $A$ on $L$ is a total circuit
unitary $\widetilde A$ on $LZ$ satisfying
\[
 \widetilde A(\ket\psi\ket{0^{|Z|}})=
 A\ket\psi\ket{0^{|Z|}}\quad\hbox{for every }\ket\psi\in\cH_L.
\]
$Z$ is called clean workspace. This is a zero-input, zero-output
condition, not the assertion $\widetilde A=A\otimes\id_Z$ on other
workspace inputs. Logical recording cells are part of $L$, not $Z$.
For disjoint additional wires $Z'$, the padding of an operator $A$ is
$A\otimes\id_{Z'}$, with factors placed by their names. Padding a
clean implementation in this way preserves its defining equality.
\end{definition}

\begin{lemma}[Workspace and inverses]\label{lem:scratch}
Clean implementations compose, admit clean implementations of their
adjoints by circuit reversal, and may reuse the same zero workspace after
padding its width to the maximum required width.
\end{lemma}
\begin{proof}
 \step{r2scratch-invariant}{The subspace $\cH_L\otimes\ket0_Z$ is mapped
 onto itself by $\widetilde A$.}
 \begin{proof}\pf\ Its restriction is the unitary $A$ under the
 isometric identification $\psi\mapsto\psi\otimes\ket0$.\end{proof}
 \step{r2scratch-inverse}{$\widetilde A^\dagger(A\psi\otimes\ket0)
 =\psi\otimes\ket0$.}
 \begin{proof}\pf\ Apply $\widetilde A^\dagger$ to the defining equality;
 surjectivity of $A$ covers every logical input.\end{proof}
 \qedstep
 \begin{proof}\pf\ After each forward or inverse implementation the
 workspace is zero by the definition and \stepref{r2scratch-inverse}.
 Pad each workspace by tensor identity on the unused part of a common
 maximum-width pool. Its unused wires stay zero. Induction over the
 finite list proves composition and reuse.\end{proof}
\end{proof}

\begin{lemma}[Exact finite control and Boolean evaluation]\label{lem:synthesis}
NOT, CNOT, Toffoli, controlled-$H$, controlled-$T^{\pm1}$, and
controlled-CNOT have constant-size exact-phase implementations over
$\mathcal G^{\pm}$ with constant clean workspace. If a total Boolean
function $f:\{0,1\}^u\to\{0,1\}^v$ is supplied by an acyclic circuit
of $\beta$ AND, XOR, and NOT gates, then
$\ket{x,y}\mapsto\ket{x,y\oplus f(x)}$ has a clean implementation of
size and workspace $O(\beta+u+v)$. For a total Boolean predicate $p$ on
a register $X$ and a unitary $U$ on a disjoint target $R$, suppose
controlled-$U$ has a specified clean implementation. Computing $p(x)$
into a fresh flag, applying controlled-$U$, and undoing the predicate
implements $\sum_x\proj x_X\otimes U^{p(x)}$.
The controlled implementation acts on neither $X$ nor the predicate
evaluation workspace, apart from using the flag as an unchanged control.
\end{lemma}
\begin{proof}
 \step{r2synth-toffoli}{Toffoli has an exact implementation.}
 \begin{proof}
  \step{r2synth-phase}{For bits $a,b,c$,
  $4abc=a+b+c-(a\oplus b)-(a\oplus c)-(b\oplus c)
  +(a\oplus b\oplus c)$ as integers.}
  \begin{proof}\pf\ Substitute $a\oplus b=a+b-2ab$ and
  $a\oplus b\oplus c=a+b+c-2ab-2ac-2bc+4abc$.\end{proof}
  \step{r2synth-ccz}{There is a diagonal circuit with multiplier
  $(-1)^{abc}$ and zero-output workspace.}
  \begin{proof}\pf\ Compute the seven indicated parities into zero
  wires with CNOTs. Apply $T$ for positive coefficients and $T^\dagger$
  for negative coefficients, then erase the parities by reversing their
  CNOTs. The phase is $e^{i\pi(4abc)/4}$ by
  \stepref{r2synth-phase}.\end{proof}
  \qedstep
  \begin{proof}\pf\ Conjugate the third wire by $H$ before and after
  \stepref{r2synth-ccz}; $HZH=X$.\end{proof}
 \end{proof}
 \step{r2synth-controls}{The asserted controlled gates have exact phases.}
 \begin{proof}\pf\ Controlled-CNOT is Toffoli. For controlled-$T$,
 compute the AND of control and target into a zero wire, apply $T$ there,
 then uncompute; replace $T$ by $T^\dagger$ for its inverse. Put
 $V=SHTHS^\dagger$. Direct matrix multiplication gives $VZV^\dagger=H$.
 Thus $(\id\otimes V)\operatorname{CZ}(\id\otimes V^\dagger)$ is
 controlled-$H$, with $\operatorname{CZ}=(\id\otimes H)
 \operatorname{CNOT}(\id\otimes H)$. Scalar phases of $V$ cancel in
 this displayed expression. NOT is $HT^4H$.\end{proof}
 \step{r2synth-evaluate}{The Boolean evaluator has the asserted bounds.}
 \begin{proof}\pf\ Store each circuit gate's value in a fresh zero
 bit in topological order: AND by Toffoli, XOR by two CNOTs, NOT by
 CNOT followed by NOT on the new bit. Copy output bits into $y$, then
 undo the stored evaluations in reverse topological order. This uses
 $O(\beta+u+v)$ gates and bits
 and preserves $x$ on every logical input.\end{proof}
 \qedstep
 \begin{proof}\pf\ The construction for a predicate is the evaluator
 of \stepref{r2synth-evaluate}, the supplied controlled implementation,
 and the inverse evaluator. On $\ket x\ket\psi\ket0$ their successive
 logical values are $(x,\psi,p(x))$, $(x,U^{p(x)}\psi,p(x))$, and
 $(x,U^{p(x)}\psi,0)$. Linearity and Lemma~\ref{lem:scratch} give the
 asserted operator and zero workspace.\end{proof}
\end{proof}

\subsection{A total joint device tick}

A \emph{device} below is not a fourth machine. It is finite classical data,
supplied by Machine~1, that couples that control to Machine~3: at each
scheduled step (a \emph{tick}) the data select at most one quantum operation
and then update the classical configuration, with the previous configuration
retained so the step is a total unitary. Machine~2 enters only through its
completed permutation steps, admitted as atoms in Section~\ref{sec:instructions}.

\begin{definition}[Finite device data, not an implicit instruction table]
\label{def:device}
A register in the computational-basis state $\ket{0\cdots0}$ is called
\emph{blank}. A device instance supplies the following finite data. The quantum pool
$Q$ has $q$ named wires. Its classical configuration $C$ has $s$ bits,
including program position, word slots, memory, stack, quantum identifiers
and live flags, console tape and length, status, and outcome. Parameters
$w\geq8,B_{\rm mem},L_{\rm slot},h,o_{\max},p$ bound word width, memory bytes, word slots, call depth,
console bytes, and instruction positions respectively. Their field layout
and a validity predicate on all $s$-bit strings are part of the data.
A canonical valid fault configuration $c_{\rm fault}$ and an initial
configuration $c_{\rm in}$ are supplied.
The data also designate one active status code
$\sigma_{\mathrm{act}}$ and one timeout status code
$\sigma_{\mathrm{to}}$, each of the status-field width.

Crucially the data include total Boolean circuits for a decoder $d(c)$,
an update $f(c,z)$, and an output extraction $o(c)$. The word $z$ has a
fixed width $k\geq1$, the maximum read width. Write $[q]=\{0,\ldots,q-1\}$.
The decoder's semantic codomain is the disjoint union
\[
 \begin{split}
 \mathcal D={}&\{\mathsf{ordinary},\mathsf{idle},\mathsf{fault}\}
 \sqcup\{\mathsf{gate}(g,i,j):g\in\{H,T,\operatorname{CNOT}\}\}\\
 &\sqcup\{\mathsf{read}(r;i_0,\ldots,i_{r-1}):0\leq r\leq k\}
 \sqcup\{\mathsf{measure}(i),\mathsf{allocate}(i):i\in[q]\}.
 \end{split}
\]
In gate and read modes all used addresses belong to $[q]$; CNOT has
$i\ne j$; read lists may repeat addresses. A one-wire gate uses
only $i$. The fixed binary layout, in order, is a 3-bit mode tag (the
seven displayed modes in order, starting at zero), a 2-bit gate tag
($H,T,\operatorname{CNOT}$ numbered $0,1,2$), two address fields, a
read-length field, and $k$ read-address fields. Measurement and allocation
use the first address field; one-wire gates do likewise. Address width is
$a_q=\max(1,\lceil\log_2(q+1)\rceil)$; length width is
$a_k=\lceil\log_2(k+1)\rceil$. Unused fields are zero. Thus decoder
width is $5+(k+2)a_q+a_k$. No alternative redundant encoding is admitted.
The supplied Boolean decoder always outputs a canonical code in
$\mathcal D$; a malformed intended action must instead encode fault.
A read XORs wire $i_j$ into sample bit $j$, starting at bit zero;
measurement XORs wire $i$ into sample bit zero only. The sample's other
bits are untouched, hence zero on a blank sample. Sample destinations
in $C$ are determined by $f$ alone, not by additional decoder fields.

For every configuration code, a gate, read, or measurement may address
only wires whose live flags are one. Allocation selects a zero live
flag and sets just that flag to one; its quantum action is identity on
that preallocated factor. Except when sending to $c_{\rm fault}$, all
other updates preserve live flags for every $z$; allocation preserves
the other flags. There is no deallocation instruction. Invalid codes,
addresses, repeated gate operands, exhausted capacities, mode violations,
or illegal instructions have decoder fault, identity quantum action,
and $f(c,z)=c_{\rm fault}$ for every $z$. The fault configuration is valid
and terminal. These are requirements on the total supplied circuits,
not an implicit run-time repair. In particular no gate may precede
allocation of its operand in a valid forward execution.
Valid terminal configurations have idle decoder and $f(c,z)=c$.
For ordinary, gate, allocation and idle modes, $f$ is independent of $z$.
For read and measurement it writes the sampled value into its specified
destination and performs the specified program-position update.

All remaining behavior is supplied by the circuits, including which
operands are immediate or indirect, signed comparison, shift directions,
division, call-frame layout, console overflow, and every output bit on
invalid encodings. No theorem identifies unspecified choices.
For example modular addition means addition modulo $2^w$ only if that
equation is imposed on the supplied circuit. This parametrization fixes
one operator for every complete instance. Instruction names 1--6 mean
the corresponding classical functions when the instance implements them;
no efficiency of an arbitrary supplied function is presumed.
The circuit for $f$ includes instruction selection and has size $\beta$.
Let $\Delta$ bound the sum of sizes of decoder and output circuits,
including their input and output ports. The encoding has
$s\geq a_q$; in particular the live-flag field contains $q$ bits.
\end{definition}

For a first reading, keep one small schematic instance in mind. Take a
pool of two wires, a configuration $C$ holding a two-bit program position,
two live flags and a status bit, with the other required fields padded to
the device layout. Total circuits $d,f$ decode position $0$ as $H$ on
wire $0$ and advance to position $1$, decode position $1$ as halted/idle,
and send every other position code to the fault configuration, subject to
the validity and live-flag checks above. Nothing in what follows depends on
the size of the fields: a device is simply a finite classical control
state $C$ together with total Boolean circuits saying, for \emph{every}
value of $C$, which single quantum operation (if any) happens next and
what $C$ becomes. The binary layout and the every-code conditions exist
so that this statement has no undefined cases when $C$ itself is later
held in qubits.

One classical update $c\mapsto f(c,z)$ is in general many-to-one, so it
has no inverse and no unitary. The construction below makes it reversible
by retaining the old configuration in a fresh register $J_t$ before
overwriting $C$. The map
$(c,j,z)\mapsto(j\oplus f(c,z),c,z)$ is then a permutation of \emph{all}
basis states, blank or not, with the displayed inverse; that totality is
what later lets a rewind act after a keep has moved the state outside the
image of the forward steps on blank records. The price is the distinction
used throughout: the operator is a unitary unconditionally, but it means
``perform the classical update'' only on blank $J_t,B_t$, where the full
tick acts as $(c,0,0)\mapsto(f(c,z),c,z)$ on the machine factors.
The order in $S_t=F_tD_t$ is essential: the quantum operation $D_t$ is
selected by the configuration \emph{before} the tick, and only afterwards
does $F_t$ change that configuration.

\begin{definition}[Snapshot, sampling, and ordering]\label{def:joint-tick}
For each scheduled tick $t$ reserve distinct registers $J_t$ of $s$ bits
and $B_t$ of $k$ bits, initially zero. They are retained. Define
\[
 F_t:\ket{c,j,z}\longmapsto
       \ket{j\oplus f(c,z),c,z}_{C J_t B_t},\qquad
 F_t^{-1}:\ket{c',j',z}\longmapsto
       \ket{j',c'\oplus f(j',z),z}.
\]
Define $G_c$ on $QB_t$ using the \emph{pre-tick} decoder: the selected
gate in gate mode; CNOTs from selected pool wires to the designated
positions of $B_t$ in read or measurement mode; identity in all other
modes. In particular coherent read and measurement have the same retained
copying operation, but only the latter has an observational interpretation.
The joint tick on the full retained list is
\[
 D_t=\sum_{c\in\{0,1\}^s}\proj c_C\otimes G_c,
 \qquad S_t=F_tD_t.                         \tag{1}\label{eq:tick}
\]
No decoder field changes during $D_t$. The update follows the quantum
action, so an address changed by $f$ cannot change that action's operand.
The logical machine register $M$ includes $C$, every $J_t,B_t$, and any
other explicit retained records; implementation workspace is separate.
For an ordinary tick with blank records,
$(c,0,0)\mapsto(f(c,0),c,0)$. For a read of basis value $x$ it is
$(c,0,0)\mapsto(f(c,x),c,x)$, with the read pool wires unchanged.

All syntactic tick occurrences of a complete term receive globally
distinct $J_t,B_t$ names, including occurrences in sequential fragments
and different marked bodies. Expanding a term's adjoint uses the same
names as the term, in inverse order; it allocates no new record.
The ordinary interpretation of a forward tick requires its records blank
at that application. If a nonblank cell is encountered the completed
unitary still applies, but it is not asserted to implement the ordinary
update. Fresh cells may be reused only with a proved zero-restoration
condition. Every name possibly read by an enclosing inverse is retained.
\end{definition}

\begin{lemma}[Joint tick well-posedness]\label{lem:joint}
Equation~\eqref{eq:tick} is a total unitary with inverse
$D_t^\dagger F_t^{-1}$, independent of blank-record promises. It has a
clean implementation over $\mathcal G^{\pm}$. On blank records it has
the stated ordinary and coherent-read actions.
\end{lemma}
\begin{proof}
 \step{r2joint-permutation}{$F_t^{-1}F_t=F_tF_t^{-1}=\id$.}
 \begin{proof}\pf\ Substitute the displayed binary formulas and use
 $a\oplus a=0$, while $z$ is preserved.\end{proof}
 \step{r2joint-controlled}{$D_t$ is unitary.}
 \begin{proof}\pf\ Each $G_c$ is a product of unitary gates and the
 projectors $\proj c$ are orthogonal and sum to identity. Hence both
 $D_t^\dagger D_t$ and $D_tD_t^\dagger$ equal identity.\end{proof}
 \step{r2joint-implementation}{Both factors have clean implementations.}
 \begin{proof}\pf\ Evaluate $f(c,z)$ by Lemma~\ref{lem:synthesis}, XOR
 it into $J_t$, restore workspace, and swap $C,J_t$ using three CNOTs
 per pair. For $D_t$, compute the decoder into workspace, enumerate
 gate/measurement actions and read position/address pairs, compute
 their selection predicates, apply each controlled gate, and uncompute.
 For reads, the predicate also tests that the position is below the
 decoded length. The decoder depends only on the unchanged $C$;
 selected read operations are controlled-CNOTs. Uncompute the decoder.
 Lemmas~\ref{lem:synthesis} and \ref{lem:scratch} apply.\end{proof}
 \qedstep
 \begin{proof}\pf\ Multiply the inverses from
 \stepref{r2joint-permutation} and \stepref{r2joint-controlled} in reverse
 order. Substitution of $j=z=0$ gives the stated forward actions.
 \stepref{r2joint-implementation} supplies the realization.\end{proof}
\end{proof}


\notedsection{Coherent reads, mark, rewind and keep}{sec:instructions}
  {This section defines \earlierforms.}
This section defines the new instructions and their elementary
well-posedness properties. The first subsection makes a coherent read one
joint tick whose sampled value is stored in $C$ and used by later ticks,
with stopping replaced by a fixed schedule. The second is needed for the
equivalence in Section~\ref{sec:primitives}: it shows an observed, adaptive
protocol is already described by retaining its measurement copies and
tracing at the end. The third is the heart of the paper: keep is a
controlled XOR $K$ from the machine register to quantum targets, a marked
term denotes $U^\dagger KU$, and a \emph{menu} is the one data-dependent
construction admitted by these structured terms---a fixed finite direct
sum $\sum_i\proj i_D\otimes\mathcal U_{P_i}$, with the selector register
$D$ retained, not a data-dependent mark. Soundness, restoration and
resources for these definitions are Sections~\ref{sec:soundness}--%
\ref{sec:resources}, respectively; equivalence of the nine basic
instructions, which excludes coherent reads and marks, is
Section~\ref{sec:primitives} and can be read before those three sections.
The separate conditional bounded dynamic interpreter is
Section~\ref{sec:tierB}.

\subsection{Coherent reads and a fixed stopping schedule}

\begin{definition}[Read instruction, return, and continuation]\label{def:read}
Primitive 10 is exactly a read-mode tick, at its decoded program position.
Before such a tick, its input value is not made available to $f$ by an
external parameter. Later ordinary ticks use the stored value in $C$.
Bounded fragments share that same $C$ and obey the global cell discipline
of Definition~\ref{def:joint-tick}. Intervening gates are composed in their
actual order, including gates changing an earlier read input.
Run a fragment for its supplied number of ticks. Its ordinary return,
stop or fault status selects idle updates for subsequent ticks, but even
an idle update uses the completed $F_t$. The schedule never stops by
observing a coherent status bit. A separate final timeout tick, included
in a declared timeout horizon $T\geq1$, has its own total update
$f_\tau(c,z)=h_\tau(c)$, where $h_\tau(c)$ is obtained from $c$ by
replacing the status field by $\sigma_{\mathrm{to}}$ if it equals
$\sigma_{\mathrm{act}}$, and otherwise preserving every bit of $c$,
on every configuration code. Its decoder
is identically idle. Thus $S_{T-1}=F_{T-1}^{\tau}$ uses $h_\tau$ in
the snapshot formula and $D_{T-1}=\id$; only $S_0,\ldots,S_{T-2}$
use the ordinary device pair $(d,f)$. The final tick has its own blank
$J_{T-1},B_{T-1}$, counted among the $T$ ticks. It is a total unitary
by the same snapshot inverse. It is an explicitly scheduled snapshot
operation, not an assertion that the device's time-independent $f$
can infer the end of its horizon. Without a declared timeout horizon a
fixed term has only its explicitly listed atoms.

For a specified continuation address $\ell$ whose binary code fits
the program-position field, define $h_\ell(c)$ by
replacing just the program-position bits with the binary code of $\ell$
and preserving every other bit, on \emph{all} codes. A continuation atom
uses a fresh $s$-bit $J$ and the same snapshot formula with $f=h_\ell$.
On blank $J$ it installs $\ell$ and retains the previous configuration.
It is a total permutation, including on nonblank $J$.
A structured source block with that continuation denotes the marked
operation below followed by this explicit atom. Its continuation record
is part of $M$ until all enclosing inverses finish. Cleanliness of the
marked operation does not imply cleanliness of this subsequent update.
\end{definition}

In particular, ordinary code after a read acts as a family of reversible
operations controlled by the read register, in the sense of
Lemma~\ref{lem:read-product}, while the operations preserve its
computational-basis sectors. A gate changing that register removes this
guarantee: the controlled description applies to the preceding segment,
and the read register can be entangled with the machine register.
Theorem~\ref{thm:sA-four} is an explicit four-qubit calculation of the
failure, using copying atoms rather than full-snapshot read ticks.

\begin{lemma}[Preserved read register]\label{lem:read-product}
Consider a fixed tick sequence whose only actions on a factor $X\subseteq Q$
are uses as CNOT controls for reads, and whose other operations preserve
each computational-basis sector of $X$. Then on $X L$, where $L$ is all
remaining retained wires, its unitary has the form
\[
 U=\sum_x\proj x_X\otimes V_x,\qquad
 V_x=V_{x,T-1}\cdots V_{x,0}.
\]
Each $V_{x,t}$ is obtained by substituting $x$ in the actual tick circuit;
it is not an independently supplied oracle. If no gates on other data
occur and all other actions are Boolean, these are permutation unitaries.
\end{lemma}
\begin{proof}
 \step{r2read-sector}{Every tick restricts to a unitary $V_{x,t}$ in
 each fixed-$x$ sector.}
 \begin{proof}\pf\ The hypothesis gives an invariant sector. Its
 finite dimension and the unitary tick imply the restriction maps it
 onto itself. Substitution of $x$ computes this restriction.\end{proof}
 \qedstep
 \begin{proof}\pf\ Multiply the sector sums using
 $\proj x\proj y=0$ for $x\ne y$. Products of permutations are
 permutations. A gate changing $X$ violates the hypothesis and gives
 no such formula for the entire sequence.\end{proof}
\end{proof}

\subsection{Bounded adaptivity and its retained interpretation}

\begin{definition}[Bounded adaptive protocol and outputs]\label{def:adaptive}
A protocol is the device data of Definition~\ref{def:device}, an initial
valid classical configuration $c_{\rm in}$ whose live flags identify
exactly the input list, and a horizon of $T\geq1$ ticks including timeout.
A measurement-outcome sequence is \emph{reachable} from $c_{\rm in}$ if
the successive configurations obtained by starting at $c_{\rm in}$ and
applying $f$ with those sampled bits (zero at non-measurement ticks) are
the configurations at which those samples are taken; the final tick uses
$h_\tau$ as specified below.
For every $c$, the decoder range excludes read mode: an intended read
or any other mode violation must select fault with update $c_{\rm fault}$.
All live-flag requirements of Definition~\ref{def:device} hold for every
code, not just for a selected input. There are no marks. At the first
$T-1$ ticks the classical configuration selects the quantum operation,
then $f$ updates it. The last tick uses exactly the idle decoder and
$h_\tau$ of Definition~\ref{def:read}, on both interpretations.
Measurement of wire $i$ has outcomes $z=0,1$ and
unnormalized update $\rho\mapsto P_{i,z}\rho P_{i,z}$, where
$P_{i,z}=\proj z_i\otimes\id$ on other wires. Its outcome is supplied
to $f$ in sample bit zero, with other sample bits zero.
Allocation selects a never-used zero factor in a pool reserved
in advance; the live flags forbid operations on unallocated factors.
At most $T$ allocations are needed beyond the input list. Terminal data
are padded as above. Designated quantum output wires $O\subseteq Q$ are
fixed, and for every measurement-outcome sequence reachable from
$c_{\rm in}$ every wire in $O$ is live in the final configuration $c_T$;
the classical output is
the total specified bit string $o(c_T)$, which includes any requested
outcome, memory and console fields. All pool wires outside $O$, including
live non-output wires, are traced in this output instrument. These rules
define it by finite repeated composition and summation over outcomes.

To retain the protocol use exactly $S_{T-1}\cdots S_0$, with the last
factor defined by $h_\tau$, on the input state and the reserved zero
non-input pool, records and samples, with $C$ initialized to $c_{\rm in}$.
Copy $o(C)$ into a fresh output register
$Y$ by its total XOR evaluator, and copy $Y$ into a retained register $E_Y$.
The environment is exactly the complement of $OY$ in the entire retained
list: in particular it contains $Q\setminus O$, $C$, all $J_t,B_t$,
$E_Y$, and evaluation workspace. The second copy makes $Y$
classical after that environment is traced. This specifies the output
extraction, rather than merely observing measurement samples.
\end{definition}

\begin{lemma}[Bounded deferred observation, proved here]\label{lem:deferred}
For a protocol satisfying all the global decoder and live-flag restrictions
and the $T\geq1$ timeout convention in Definition~\ref{def:adaptive},
the retained construction there is an isometry.
Tracing its designated environment gives exactly the protocol's output
instrument, including classical outputs, for inputs with any reference.
\end{lemma}
\begin{proof}
 \step{r2defer-expand}{After $t$ ticks the retained vector can be written
 $\sum_{\zeta} A_{t,\zeta}\ket\psi\otimes\ket{e_{t,\zeta}}$, where
 $\zeta$ is the length-$t$ list of sampled bits, padded by zeros at
 nonmeasurement ticks. Distinct possible lists have orthogonal retained
 sample registers. $A_{t,\zeta}$ is the ordered product of the protocol's
 selected gates and measurement projectors on the reserved pool, composed
 with the initial zero embedding of its non-input wires; its configuration is the
 deterministic configuration prescribed by those samples.}
 \begin{proof}
  \step{r2defer-base}{The assertion holds at $t=0$.}
  \begin{proof}\pf\ There is one empty list, the initial zero embedding,
  and the specified initial configuration and blank records.\end{proof}
  \step{r2defer-next}{The assertion at $t$ implies it at $t+1$.}
\begin{proof}
 \step{r2defer-next-hyp}{\assume{The asserted expansion holds at
 $t<T$.}\prove{It holds after the next scheduled tick.}}
 \begin{proof}
  \step{r2defer-next-timeout}{\case{$t=T-1$.}}
  \begin{proof}\pf\ The scheduled quantum action is identity and
  the new sample is zero. The snapshot applies exactly $h_\tau$
  and stores the old configuration, giving the prescribed next
  configuration without altering the quantum product.\end{proof}
  \step{r2defer-next-gate}{\case{$t<T-1$ and the decoder selects
  gate mode.}}
  \begin{proof}\pf\ Multiply each summand by the selected gate.
  The sample is zero; the snapshot computes $f(c,0)$ and stores $c$.
  This is the corresponding protocol step.\end{proof}
  \step{r2defer-next-measure}{\case{$t<T-1$ and the decoder selects
  measurement mode on wire $i$.}}
  \begin{proof}\pf\ The copying isometry is
  $v\mapsto\sum_{z=0}^1P_{i,z}v\otimes\ket z_{B_t}$,
  with unused sample bits zero. The snapshot stores the old
  configuration and computes precisely the next one from $z$.
  These are the two projector terms of the protocol update.\end{proof}
  \step{r2defer-next-allocate}{\case{$t<T-1$ and the decoder selects
  allocation mode.}}
  \begin{proof}\pf\ Its zero live flag implies the factor is a
  non-input factor initially zero and never acted on: ordinary
  updates preserve flags, allocation only changes zero to one,
  and every quantum action requires a one flag. A fault is terminal
  and cannot be followed by allocation. The factor is therefore
  still zero. The identity quantum action and prescribed flag update,
  with zero sample, implement allocation from the reserved pool.\end{proof}
  \step{r2defer-next-other}{\case{$t<T-1$ and the decoder selects
  fault, ordinary or idle mode.}}
  \begin{proof}\pf\ These modes have identity quantum action,
  zero new sample and their specified updates. Their snapshots
  store the old configuration. Read mode is excluded on every code.\end{proof}
  \step{r2defer-next-orthogonal}{Distinct extended sample lists
  have orthogonal retained sample registers.}
  \begin{proof}\pf\ Distinct old lists differ in an old retained
  sample bit. Equal old lists with different new samples differ
  in $B_t$. No case erases these bits.\end{proof}
  \qedstep
  \begin{proof}\pf\ The cases exhaust the schedule and permitted
  modes. Each appends the correct gate or projector and deterministic
  configuration update; \stepref{r2defer-next-orthogonal} establishes
  the required orthogonality. Sum over the old sample lists.\end{proof}
 \end{proof}
 \qedstep
 \begin{proof}\pf\ By \stepref{r2defer-next-hyp}.\end{proof}
\end{proof}
  \qedstep
  \begin{proof}\pf\ Finite induction using
  \stepref{r2defer-base} and \stepref{r2defer-next}.\end{proof}
 \end{proof}
 \step{r2defer-output}{Tracing the environment gives the required instrument.}
 \begin{proof}\pf\ Orthogonality of the samples removes cross terms
 for different lists in \stepref{r2defer-expand}. The extra copy of
 $o(c_T)$ removes off-diagonal classical output labels. The remaining
 terms are exactly $A_{T,\zeta}\rho A_{T,\zeta}^\dagger$ grouped by
 $o(c_T)$, with all $Q\setminus O$ wires traced, as in the definition.\end{proof}
 \qedstep
 \begin{proof}\pf\ The initial embeddings and subsequent unitaries are
 isometries. The calculation in \stepref{r2defer-output} applies to
 matrix units by linearity, and every operator acts as identity on a
 reference. It therefore applies to all density operators with references.
 No unsupplied measurement theorem is used.\end{proof}
\end{proof}

\subsection{Structured marks, allocation types, and a finite menu}

\begin{definition}[Keep]\label{def:keep}
For the specified current machine list $M$ of $m_M$ bits, choose a total
Boolean circuit $v:\{0,1\}^{m_M}\to\{0,1\}^{|r|}$ and a fixed ordered
list $r$ of distinct quantum wires disjoint from $M$. Set
\[
 K_{v,r}=\sum_{u\in\{0,1\}^{m_M}}X_r^{v(u)}\otimes\proj u_M,
 \qquad X_r^z=\bigotimes_i X_{r_i}^{z_i},                 \tag{2}
\]
with identity on all other wires. A keep list is a finite product of
these maps. Different entries may share target wires; repeated XORs
have their usual meaning. Every target is in the register context of
its enclosing mark. There is no computed-address exception in this
definition. Targets used by the body are allowed, without a cleanliness
promise. A standalone use is a circuit atom with the same typing.
\end{definition}

\begin{lemma}[Keep well-posedness]\label{lem:keep}
Every keep list is a unitary involution and has a clean finite-alphabet
implementation.
\end{lemma}
\begin{proof}
 \step{r2keep-sectors}{In a fixed $u$ sector each entry is a product
 of Pauli $X$ operators; all entries commute and square to identity.}
 \begin{proof}\pf\ Every entry preserves $M$, and $X^2=\id$ on each
 target. Overlaps merely add their exponents modulo two.\end{proof}
 \qedstep
 \begin{proof}\pf\ Sum over the orthogonal $u$ sectors for unitarity
 and the involution identity. Evaluate $v$, XOR its output into $r$,
 and uncompute by Lemma~\ref{lem:synthesis}; $M$ is preserved.\end{proof}
\end{proof}

We use \emph{Tier~A} for the class defined next: marks and their pairings are
fixed in the syntax, hence known to the classical control computer before any
quantum input. \emph{Tier~B} is the data-dependent analogue, in which
whether a rewind is reached and which mark it names can depend on values read
from qubits. Section~\ref{sec:tierB} defines a bounded interpreter and proves
conditional results under a supplied certificate; constructing such a
certificate for general dynamic programs remains open
(Sections~\ref{sec:tierB} and \ref{sec:open-problems}).

\begin{definition}[Tier~A, fixed-space terms and typed outer programs]
\label{def:tiera}
Fix a retained list $L=QM$. An atom is an occurrence-labelled elementary
quantum gate from $\mathcal G^{\pm}$ (or the explicitly declared ideal
alphabet), a joint tick of Definition~\ref{def:joint-tick}, a continuation atom, a specified
completed RTM step, or a finite composition of these. Each is a specified
logical unitary, with an exact clean implementation where finite-alphabet
resources are claimed. Boolean elementary gates on machine bits are
also atoms by Lemma~\ref{lem:synthesis}.
Every atom specifies a support list contained in $L$ and acts as identity
on its complement. An RTM atom's support is its specified
$m_{\rm RTM}$-bit encoding register. For a joint-tick atom the supplied
decoder, on \emph{every} $c$, excludes measurement and allocation modes;
mode violations must instead select fault with update $c_{\rm fault}$.
Read mode is allowed, with all the live-flag conditions of
Definition~\ref{def:device}. The explicitly scheduled timeout snapshot
of Definition~\ref{def:read} is also an allowed joint-tick atom.
Atom labels include syntactic occurrence identifiers: two written copies
of a tick receive distinct recording cells even if their instruction
data agree. Adjoint expansion retains the original occurrence identifiers.
The grammar, interpreted on the \emph{same} $L$, is
\[
 P::=I\mid a\mid P;P\mid \mathsf{mark}\{P\}\mathsf{rewind}\ K.
\]
Denote the last constructor by $\mathsf M(P,K)$ and its term denotation
by $\mathcal U_P$. Set
\[
 \mathcal U_I=\id,
 \quad\mathcal U_a=A_a,
 \quad\mathcal U_{P;R}=\mathcal U_R\mathcal U_P,
 \quad\mathcal U_{\mathsf M(P,K)}
   =\mathcal U_P^\dagger K\mathcal U_P.       \tag{3}
\]
The sans-serif $\mathsf M(P,K)$ is a term constructor, not the register $M$.
The mark is a static delimiter in machine 1. Inverting a completed inner
term takes its complete adjoint, including its keep. A source block with
post-rewind address $\ell$ is, by definition,
$\mathsf M(P,K);\mathsf{cont}_\ell$, with the atom of
Definition~\ref{def:read}. The first operation has exactly the denotation
$\mathcal U_P^\dagger K\mathcal U_P$; it need not restore the machine-held
position. Position restoration requires a separately proved condition
on the allowed input space. The second operation installs $\ell$ when
its own fresh record is blank, retaining the previous configuration,
regardless of the position left by the first operation.

For changing register lists the outer grammar is
$E::=I_L\mid A_B\mid\mathsf{run}_L(P)\mid E;E$,
with types $I_L:L\to L$, $A_B:L\to LB$ for globally fresh $B$,
$\mathsf{run}_L(P):L\to L$, and composition only when the middle
lists coincide. Its denotations are respectively identity, zero embedding,
$\mathcal U_P$, and ordered composition. Outputs are a
designated subset; dropping is performed afterward on its complement.
All logical records and implementation workspace are allocated before
their enclosing run. Disjoint recording cells are assigned to syntactic
occurrences, not to the number of gates in the expanded circuit.
There is a padding rule: a term $P$ on $L$ may be used on $LZ$ for any
disjoint list $Z$, by lifting every atom and keep with tensor identity;
its denotation is $\mathcal U_P\otimes\id_Z$. In particular
$\mathsf{run}_{LZ}(P\otimes I_Z)$ is the logical run after allocating
workspace. A physical implementation may use $Z$, and agrees with this
logical run on zero $Z$ by Definition~\ref{def:scratch}; equality on
other $Z$ states is not required.

\emph{Restrictions of this tier:} no allocation, drop, observed measurement,
external console observation, or actual stop inside a marked body. Bounded
idle statuses and updates of an unobserved console tape are allowed.
Every mark instance and pairing is static. A loop or call containing marks
is covered only if it has a supplied fixed finite syntactic expansion;
data-dependent numbers of mark visits and general recursive calls are
excluded. Bounded loops without marks can be joint-tick fragments.
Source allocation inside a mark is not automatically admitted by this
grammar. Pre-reserving its wires is a correction requiring a separate
zero-input justification, not equality on arbitrary wire states.
\end{definition}

\begin{lemma}[Algebraic typing of structured terms]\label{lem:term-types}
Every finite fixed-space term of Definition~\ref{def:tiera} denotes a
total unitary. Every well-typed outer term denotes an isometry into its
entire retained output space. These conclusions concern the algebraic
denotation, without an assertion about compilation cost or source coverage.
\end{lemma}
\begin{proof}
 \step{r2types-induction}{Every fixed-space term denotes a unitary.}
 \begin{proof}
  \step{r2types-base}{\case{The term is identity or an atom.}}
  \begin{proof}\pf\ Identity is unitary. Elementary gates, completed
  permutations and joint ticks are unitary by their matrices,
  Lemmas~\ref{lem:completion} and \ref{lem:joint}. The continuation
  has the snapshot inverse, and finite compositions of atoms multiply
  unitaries.\end{proof}
  \step{r2types-sequence}{\case{The term is $P;R$, and the assertion
  holds for its proper subterms.}}
  \begin{proof}\pf\ The product of the two unitary denotations has
  inverse the reversed product of their adjoints.\end{proof}
  \step{r2types-mark}{\case{The term is $\mathsf M(P,K)$, and the
  assertion holds for $P$.}}
  \begin{proof}\pf\ Put $U=\mathcal U_P$. By
  Lemma~\ref{lem:keep}, $K$ is unitary; multiplication gives
  \[
   (U^\dagger KU)^\dagger(U^\dagger KU)=\id.
  \]
  The opposite product also equals identity.\end{proof}
  \qedstep
  \begin{proof}\pf\ Structural induction on the finite syntax, by
  \stepref{r2types-base}, \stepref{r2types-sequence}, and
  \stepref{r2types-mark}.\end{proof}
 \end{proof}
 \step{r2types-outer}{Every well-typed outer term is an isometry.}
 \begin{proof}\pf\ Identity and run are isometries by
 \stepref{r2types-induction}. Allocation satisfies $A_B^\dagger A_B=\id$
 because $\ket0_B$ has unit norm. If $V^\dagger V=\id$ and
 $W^\dagger W=\id$ with matching middle lists, then
 $(WV)^\dagger WV=V^\dagger V=\id$. Structural induction gives
 the assertion.\end{proof}
 \qedstep
 \begin{proof}\pf\ By \stepref{r2types-induction} and
 \stepref{r2types-outer}.\end{proof}
\end{proof}

\begin{lemma}[Moving an unused allocation earlier]\label{lem:reserve}
Let $B$ be fresh, and let a finite circuit segment $V$ before its
allocation act on $L$ alone. Then
$(V\otimes\id_B)A_B=A_BV$. Consequently finitely many allocations can
be moved to the start of a forward circuit if none of their names is used
before its original allocation. The equality is of embeddings with
$B$ initialized to zero, not of an inverse allocation with a unitary.
\end{lemma}
\begin{proof}
 \step{r2reserve-tensor}{Both sides send $\psi$ to $V\psi\otimes\ket0_B$.}
 \begin{proof}\pf\ Substitute the definition of $A_B$ and the
 tensor-product action on distinct factors.\end{proof}
 \qedstep
 \begin{proof}\pf\ Move each fresh allocation across each earlier
 operation by \stepref{r2reserve-tensor}. There are finitely many such
 pairs. This proves the equality on the stated input space only.
 It does not define an adjoint deletion inside a marked body.\end{proof}
\end{proof}

\begin{definition}[Restricted menu; dynamic source semantics remains open]
\label{def:tierb}
Let $D$ be a $d$-qubit selector disjoint from a common retained payload
$L$. Choose $0\leq j\leq2^d$ and Tier~A terms $P_0,\ldots,P_{j-1}$
on $L$. For the implementation claim, each term must supply an explicit
exact clean circuit over $\mathcal G^{\pm}$, with its actual matrix
phase, on the same zero-initialized workspace list after padding.
An ideal-alphabet term without such a supplied circuit is excluded from
that claim; the algebraic direct sum still exists for any total payload
unitaries. No general exact controlled-rotation synthesis is presumed.
All $D$ wires are preserved throughout this operation and retained for
every enclosing inverse. Define
\[
 \mathsf{SELECT}=
 \sum_{i<j}\proj i_D\otimes\mathcal U_{P_i}+
 \sum_{j\leq i<2^d}\proj i_D\otimes\id_L.                  \tag{4}
\]
Implement this by enumerating the $j$ circuit lists, controlling each
elementary gate on $D=i$. The schedule is fixed; shorter lists may be
padded with identity. No selected circuit acts on $D$. If $D$ was
computed earlier, that computation is an explicit reversible operation
with retained records; no erasure of $D$ is implied by selection.

The finite-menu extension of the Tier~A syntax admits this construction
as a derived atom $\mathsf{sel}_\nu(P_0,\ldots,P_{j-1})$ on $DL$,
with identity on any other named factors. Its support contains $D,L$;
the global machine list includes $D$ and the payload's machine factors.
The entries here are base Tier~A terms, not menus, so this rule is not
circular. A derived atom may occur in a sequence or a marked body.
Every payload tick and continuation has its own globally fresh occurrence
identifier across all entries and menu occurrences. Reversing this atom
reverses its entire controlled circuit with the same identifiers and
selector wires; it introduces no new cells. Its cost as an atom is the
length of that supplied controlled circuit. Its internal payload ticks
and continuations are counted once for logical record allocation.

More generally a supplied list of $N\geq0$ total \emph{unitaries}
$S_t:\cH_L\to\cH_L$ is a finite unitary schedule. Calling a member
``rewind'' supplies no source semantics. Such arbitrary schedules are
not a definition of dynamic mark/rewind and carry no equivalence claim.
The construction admitted by this definition is (4), with its stated
static menu. Theorems~\ref{thm:tB-sound} and \ref{thm:tB-cost} separately
prove bounded dynamic-interpreter soundness and costs under the supplied
certificate of Definition~\ref{def:tB-cert}. A translation of source-level
mutable mark stacks, repeated visits, and data-dependent inverse stopping
into that certified interpreter remains open
(Open problem~\ref{prob:tB-compiler}).
\end{definition}

\begin{lemma}[Menu well-posedness]\label{lem:menu}
Under Definition~\ref{def:tierb}, if the payload terms denote unitaries,
(4) is a unitary on $DL$. For the supplied exact finite-alphabet lists
required there, its specified controlled implementation has the same
logical action, with zero workspace restored. Adding such derived atoms
to Tier~A preserves the conclusions of Lemma~\ref{lem:term-types}.
\end{lemma}
\begin{proof}
 \step{r2menu-blocks}{The summands have disjoint selector sectors
 covering $\cH_D$, and unitary payloads.}
 \begin{proof}\pf\ Capacity $j\leq2^d$ makes every projector defined;
 the binary basis resolves identity. Multiply (4) by its adjoint in
 either order to obtain identity.\end{proof}
 \step{r2menu-implementation}{In sector $i<j$ precisely list $i$ acts,
 with its specified phase; in other sectors every gate is identity.}
 \begin{proof}\pf\ The equality predicates depend only on the
 unchanged $D$. Lemma~\ref{lem:synthesis} controls each elementary
 gate exactly, and Lemma~\ref{lem:scratch} restores workspace. No
 independent scalar phase may be discarded from a selected list.\end{proof}
 \qedstep
 \begin{proof}\pf\ By \stepref{r2menu-blocks} and
 \stepref{r2menu-implementation}, followed by linearity across sectors.
 For the extension, the new atom case is \stepref{r2menu-blocks};
 the product, marked-term and outer-isometry induction steps in
 Lemma~\ref{lem:term-types} are unchanged. For implementations use
 \stepref{r2menu-implementation} and Lemma~\ref{lem:scratch}.
 \end{proof}
\end{proof}

\subsection{Costs, certificates and cleanliness: definitions used later}
The definitions in this subsection are used in Sections~\ref{sec:primitives}--%
\ref{sec:resources}. They are collected here, next to the marked terms they
describe; no result is asserted in this subsection.

\begin{definition}[Tick selection cost]\label{def:selection-cost}
For the data of Definitions~\ref{def:device} and \ref{def:adaptive}, write
\[
\mathcal R=\beta+\Delta+s+k+(q^2+kq+1)(s+k+1).
\]
The supplied sizes include Boolean-circuit input and output ports. Uniform families and unary numerical bounds have the meaning of Definition~\ref{def:rtm}.
\end{definition}

\begin{remark}[Explicit corrections to the primitive list]\label{rem:corrections}
The corrections to the nine instructions that the equivalence requires
are listed, with their justifications, in Section~\ref{sec:summary}; the
additional restrictions for marked code are those of
Definition~\ref{def:tiera} for Tier~A and
Definition~\ref{def:tB-cert} for the conditional Tier~B interpretation.
\end{remark}

\begin{definition}[Expansion and cost model]\label{def:costs}
Assign every atom and keep an explicit clean circuit over $\mathcal G^{\pm}$
of length $c_a,c_K$. One gate in this alphabet costs one sequential gate
tick (a gate tick, the time unit, is not a joint tick $S_t$; one joint
tick costs many gate ticks), with no parallel execution assumed. Define the unsimplified expansion
by concatenation and by body, keep, reversed adjoint body for a mark.
Let $C(P)$ be its gate count, $L(P)$ the sum of all $c_a,c_K$ over
syntactic occurrences counted once, and $d(P)$ the maximum number of
enclosing marks: zero on atoms and identity, maximum under sequence,
and one plus body depth at a mark. A zero-initialized scratch pool of
width $a_*$ is the maximum local implementation workspace.

Let $n_{\rm tick}(P)$ count distinct syntactic joint-tick occurrences,
$n_{\rm cont}(P)$ continuation occurrences, and $q_{\rm new}(P)$ declared
fresh data wires in the enclosing typed outer program. With one shared
$s$-bit configuration and $m_{\rm extra}$ explicitly declared other
logical wires (including any RTM or Boolean-value registers), the derived
logical machine size for the snapshot construction is
\[
 m_M=s+(s+k)n_{\rm tick}(P)+s n_{\rm cont}(P)+m_{\rm extra}.
\]
For the menu extension, each derived menu atom is one leaf of cost
equal to its compiled controlled-list length in $C(P),L(P),d(P)$.
Its payload occurrences contribute to $n_{\rm tick},n_{\rm cont}$ for
record allocation, but their gate costs are not counted again outside
the atom cost. For selectors of widths $\delta_\nu$, add
$\sum_\nu\delta_\nu$ fresh disjoint selector wires to
the displayed machine size unless already included in $m_{\rm extra}$,
and count the records of their explicit computations once. The same
selector names are retained under adjoints. Selectors computed temporarily during
one tick are clean scratch and are not retained logical selectors.
Compiler time means steps of a multitape classical TM, with explicit
output gate list and binary wire labels; it is distinct from gate ticks.
\end{definition}


\begin{definition}[Approximation certificates]\label{def:certificates}
For a fixed expanded list of $N_g\ge1$ unitary factors $A_i$, a certificate procedure supplies, for every rational $0<\eta<1$, a clean $\mathcal G^{\pm}$ implementation of a logical unitary $\widetilde A_i$ with
\[
\|\widetilde A_i-A_i\|\le\eta,\qquad
\text{length}\le c_*\log^\alpha(2+1/\eta),
\]
where $c_*>0$, $\alpha\ge0$ are fixed independently of $i,N_g,\eta$.
The error concerns actual matrices, including the entire controlled matrix when present. Inverses use adjoints of the same chosen approximating circuits. Occurrences inside inverses and controls are counted in $N_g$.
\end{definition}

\begin{remark}[Scope of the approximation hypothesis]\label{rem:sk}
The certificate hypothesis above is an explicit restriction: no
unconditional approximation theorem for arbitrary additional gates is
proved here. Its removal for arbitrary effectively specified bounded-arity
gates over $\mathcal G$, with all control phases accounted for, remains
an open synthesis task.
Dawson--Nielsen's instruction-set hypothesis is a finite inverse-closed
subset of $SU(d)$ whose generated subgroup is dense in $SU(d)$; their Theorem~1 gives
operator-norm length $O(\log^c(1/\eta))$ for fixed $d$
\cite[Definition~1, PDF p.~2; Theorem~1, PDF p.~4]{DawsonNielsen2006}.
This statement is not applied literally to $\{H,T,\operatorname{CNOT}\}$:
its determinants and inverse closure do not satisfy those hypotheses.
The exact controlled finite gates needed above were instead constructed
in Lemma~\ref{lem:synthesis}. No approximation step relies on an
unread source or on discarding a control-dependent phase.
\end{remark}

\begin{definition}[Cleanliness, reached subspaces, and restoration]
\label{def:clean}
For a fixed unit vector $\ket{m_0}_M$, define
$E\psi=\psi\otimes\ket{m_0}$ and $P_0=EE^\dagger$.
A unitary $W$ is clean on this input space if $WE=EV$ for a unitary
$V$ on $Q$. For a particular input density operator the weaker condition
is support of its final state in $\operatorname{ran}E$.
For nested use let $\mathcal S\subseteq\cH_L$ be a specified reached
subspace with orthogonal projection $P_{\mathcal S}$. It is computed
mathematically as the image of the initial allowed subspace under the
fully specified preceding isometry, or it can be a supplied containing
subspace. It need not factor across data and machine registers.
Given a designated reusable auxiliary factor $A\subseteq L$, put
$P_A=\id_{L\setminus A}\otimes\proj0_A$. Reset on $\mathcal S$ means
$(\id-P_A)WP_{\mathcal S}=0$; for reuse the incoming subspace must
also satisfy $\mathcal S\subseteq\operatorname{ran}P_A$.
For an initially entangled machine factor, preservation of its reduced
state is a different condition, called marginal restoration, and does
not authorize replacing it by fresh zeros or tracing it before an inverse.

Cleanliness here concerns the logical machine register $M$. It is a
different property from a \emph{clean implementation}
(Definition~\ref{def:scratch}), which concerns only implementation
workspace $Z$: a marked operation can have a clean implementation while
leaving $M$ entangled with $Q$ (Theorem~\ref{thm:sA-four}), and
Theorem~\ref{thm:sA-sound} asserts the former kind of statement, not
cleanliness of $M$. In phrases such as ``clean XOR'' or ``clean CNOT'',
clean means that the auxiliary factors of the stated initialized
computation return to their initial vectors; it is not merely a
workspace assertion.
\end{definition}


\begin{definition}[Restricted delayed-erasure comparison]\label{def:delayed}
For $1\leq j\leq d$ let
$f_j:\{0,1\}^{b_{j-1}}\to\{0,1\}^{b_j}$ be deterministic.
A specified total unitary evaluator $A_j$ of cost $a_j$ preserves its
input and, with output and workspace zero, writes
$(f_j(x),g_j(x))$ in computational basis into an output of $b_j$ bits
and private workspace of $u_j$ bits. It touches no other factors.
Copying that output into a disjoint external target costs $k_j$;
compute, copy, inverse gives a composite XOR unitary of cost
$o_j=2a_j+k_j$. The eager composite at level $j$ computes the previous
composite into a zero intermediate, applies the level-$j$ XOR unitary,
and inverts the previous composite. The delayed construction evaluates $A_1,\ldots,A_d$ on
distinct retained intermediates and workspaces, copies the last value,
then inverts $A_d,\ldots,A_1$. Targets are never evaluator inputs and no
other gate changes the retained values during these sequences.
Theorem~\ref{thm:resA-delayed} uses these compositions under a more general
evaluator hypothesis, (RA8) there, of which the basis-valued $g_j(x)$ above
is the special case $\ket{\gamma_{j,x}}=\ket{g_j(x)}$.
\end{definition}


\begin{remark}[Open point: remaining scope]\label{problem:remaining}
Theorems~\ref{thm:tB-sound}--\ref{thm:tB-cost} give conditional bounded
interpreter results, a static specialization, and the exact disposal
criterion. They do not construct a certified interpreter for general
dynamic source mark/rewind, data-dependent repeated mark visits or an
inverse stopped at a coherent mark identifier
(Open problem~\ref{prob:tB-compiler}). The full-snapshot form of the
four-qubit example also remains open.
They also do not settle general phase-correct approximation certificates
for arbitrary additional gates or an optimal recursive resource bound.
Source allocation within a marked body requires a proved reservation
translation and remains excluded from Tier~A as stated.
The proved deferred-observation lemma concerns bounded adaptive protocols
without marks. It does not authorize observed measurement inside a
reversible marked body.
These open points carry no proof step in the settled results.
An optimal time--space trade-off for general intervening quantum gates,
including a transformation to delayed erasure, also remains open.
\end{remark}

\begin{remark}[Collected restrictions]\label{rem:r3-restrictions}
The complete restrictions are collected in Section~\ref{sec:summary}.
They include the supplied-circuit and final-instrument scope of the
nine-instruction equivalence, the static Tier~A rules, and the conditional
Tier~B certificate and implementation hypotheses; none is inferred from
an instruction name alone.
\end{remark}

\section{Equivalence of the basic primitives}\label{sec:primitives}
\subsection{The nine basic instructions}
The instruction list being compared with the circuit model is the following.
\begin{enumerate}
\item Integer arithmetic and comparison on named values.
\item Reading and writing classical memory.
\item Conditional jump.
\item Function call and return.
\item Stop with an outcome.
\item Write one byte to a console output tape.
\item Allocate fresh zero qubits.
\item Apply $H$, $T$, and controlled $X$ (CNOT), with optional ideal rotations treated separately.
\item Computational-basis measurement with an ordinary recorded result available to later classical control.
\end{enumerate}
Instructions 1--6 refer to the complete bounded classical functions and capacities supplied in Definition~\ref{def:device}. The quantum instructions and their observed or retained interpretations are fixed in Definitions~\ref{def:circuit} and \ref{def:adaptive}. Coherent reads and marked computations are additional instructions, outside this nine-instruction equivalence.

\subsection{The meaning and cost of bounded equivalence}

\begin{definition}[The equivalence asserted here]\label{eqv:def}
Two bounded descriptions are observationally equivalent if their channels
to the same ordered quantum output $O$ and classical outcome register $Y$
agree, including after tensoring with an arbitrary reference identity.
Both descriptions have the same ordered input list. Observational equivalence is equality of the final instruments on $OY$; it does not assert that a console byte is produced at the same tick, that a run stops at the same tick, or that any intermediate classical value is available to an external reader before the final output.
The classical register is diagonal in the computational basis. Thus this
is equality of instruments with their outcome labels retained, not merely
equality of outcome probabilities. Before taking any partial trace we
specify an isometry on every input, including all inaccessible factors.
Different implementations may have different environment registers;
equality of their entire retained vectors is not part of observational
equivalence. When implementing a particular specified retained construction,
we do assert equality of that construction's isometry, with additional
zero-output workspace.

Every allocation counted below is a one-qubit allocation. A larger
allocation is expanded into these. A circuit description includes ordered
input, output, and inaccessible lists and any classical output labels.
Gate cost counts elementary gates, not zero embeddings or partial traces.
Classical time is multitape Turing-machine time with an explicit output
description, not unit-cost arithmetic on words of arbitrary length.
\end{definition}

\begin{lemma}[Cost of the specified joint tick]\label{eqv:tick-cost}
For the data of Definition~\ref{def:device}, define $\mathcal R$ as in
Definition~\ref{def:selection-cost}. One ordinary joint tick and its adjoint have exact
clean implementations of $O(\mathcal R)$ gates and $O(\mathcal R)$
workspace qubits over $\mathcal G^{\pm}$. The logical cells $J_t,B_t$
are additional to that workspace. The decoder always emits the canonical
layout of Definition~\ref{def:device}; a malformed intended action emits
canonical fault, not malformed output requiring a subsequent repair.
The separately scheduled timeout is treated in Theorem~\ref{eqv:forward}.
\end{lemma}
\begin{proof}
 \step{eqv-tick-update}{The snapshot factor costs
 $O(\beta+s+k)$ gates and workspace.}
 \begin{proof}\pf\ Compute $f(C,B_t)$ into $J_t$ by the XOR evaluator
 of Lemma~\ref{lem:synthesis}, restoring its workspace. Swap $C,J_t$
 with $3s$ CNOTs. This is exactly $F_t$, on all inputs, by its formula
 in Definition~\ref{def:joint-tick}.\end{proof}
 \step{eqv-tick-decode}{The selected quantum action has a clean circuit
 of size $O(\Delta+(q^2+kq+1)(s+k+1))$.}
 \begin{proof}
  \step{eqv-tick-flags}{One may form all relevant action flags using
  $O(\Delta+(q^2+kq+1)(s+k+1))$ Boolean gates and temporary bits.}
  \begin{proof}\pf\ Evaluate the canonical decoder. Enumerate its
  $O(q^2+1)$ gate/measurement actions and its $kq$ read
  position/address pairs. Each test compares mode and gate tags and
  at most two $a_q$-bit addresses; a read test also compares its
  fixed position with the $a_k$-bit read length. Its cost is
  $O(a_q+a_k+1)=O(s+k+1)$, using $s\geq a_q$ and $a_k\leq k$.
  Unused canonical fields are already zero and need no repeated test.
  Measurement selects sample bit zero; a read selects its listed bit.
  Decoder ports are counted in $\Delta$. For $q=0$ every
  quantum-action flag is false. All flags depend only on pre-tick $C$.
  \end{proof}
  \step{eqv-tick-actions}{Apply the flagged operations and erase all
 flags and decoded data.}
\begin{proof}
 \step{eqv-tick-actions-cost}{Each of the $O(q^2+kq+1)$ enumerated
 operations has an exact constant-cost controlled implementation.}
 \begin{proof}\pf\ Each is controlled-$H$, controlled-$T$ or
 controlled-CNOT, including the CNOT from a sampled pool wire to
 $B_t$. Apply Lemma~\ref{lem:synthesis}, with constant workspace.
 \end{proof}
 \step{eqv-tick-actions-gate}{\assume{On a computational-basis
 sector of $C$, the decoder selects gate or measurement mode.}
 \prove{The flagged product applies the selected operation exactly once.}}
 \begin{proof}\pf\ The decoder enables at most one operation.
 Enumerate a one-wire gate once for each wire $i$, using its first
 address field. Its second field is canonically zero and is not
 tested again. Thus the selected operation is applied exactly once.
 \end{proof}
 \step{eqv-tick-actions-read}{\assume{On a computational-basis
 sector of $C$, the decoder selects read mode.}\prove{The flagged
 product is precisely the prescribed read-copying operation.}}
 \begin{proof}\pf\ Precisely the listed read copies are enabled.
 Their targets are distinct bits of $B_t$, disjoint from their pool
 controls. These CNOTs commute even for repeated pool addresses,
 so their enumeration order does not change the specified read.
 \end{proof}
 \step{eqv-tick-actions-idle}{\assume{On a computational-basis
 sector of $C$, the decoder selects any remaining mode.}
 \prove{The flagged product is identity on $QB_t$.}}
 \begin{proof}\pf\ No gate, measurement or read flag is enabled;
 the selected quantum action is identity.\end{proof}
 \step{eqv-tick-actions-clean}{All flags and decoded data can be
 erased with the asserted resource bound.}
 \begin{proof}\pf\ Their input $C$ is unchanged. Reverse their
 evaluations to restore zero workspace. Reversible Boolean evaluation
 changes the Boolean counts by a constant factor; input/output-port
 cost is absorbed by $\Delta+s+k$ and the displayed bound in
 \stepref{eqv-tick-flags}.\end{proof}
 \qedstep
 \begin{proof}\pf\ The three conditional action statements cover
 every decoder mode on every basis sector of $C$; linearity gives
 the selected action on arbitrary inputs. Steps
 \stepref{eqv-tick-actions-cost} and \stepref{eqv-tick-actions-clean}
 give its implementation and erasure.\end{proof}
\end{proof}
  \qedstep
  \begin{proof}\pf\ By \stepref{eqv-tick-flags} and
  \stepref{eqv-tick-actions}.\end{proof}
 \end{proof}
 \qedstep
 \begin{proof}\pf\ Add \stepref{eqv-tick-update} and
 \stepref{eqv-tick-decode}. Lemma~\ref{lem:scratch} permits shared
 workspace and circuit reversal. No quantum random-access operation
 occurs in this construction.\end{proof}
\end{proof}

\begin{theorem}[Device to circuit, with explicit classical cost]
\label{eqv:forward}
Let an instance satisfy Definitions~\ref{def:device},
\ref{def:joint-tick}, and \ref{def:adaptive}, with no primitives 10 or 11.
Assume $T\geq1$ consists of $T-1$ ordinary ticks and the separately
scheduled final timeout with designated codes $\sigma_{\mathrm{act}}$,
$\sigma_{\mathrm{to}}$. Every forward tick uses distinct blank records.
The initial $c_{\rm in}$ is valid with exactly the input wires live;
read mode is excluded on every configuration code. The canonical total
decoder and update obey all-code live checks and flag preservation.
The fixed output list $O$ is live in the final configuration for every
reachable measurement-outcome sequence; earlier allocation of its wires
is allowed. The entire complement of $OY$ is traced at output.
All bounds, total circuits and initial
classical bits are supplied explicitly. Let $y=|Y|$ be the number of output bits of $o$, and let $d$ be the
bit length of those supplied descriptions, including unary bounds.
There is a finite retained circuit implementing exactly the specified
isometry, with at most
\[
 G=O(T\mathcal R),\qquad
 Q_*=q+s+T(s+k)+2y+O(\mathcal R)
\]
gates and qubits respectively. Preparing $c_{\rm in}$ is included in
$G$. The gate alphabet is $\mathcal G^{\pm}$; using $\mathcal G$ instead
gives $G_{\mathcal G}\leq7G_{\mathcal G^{\pm}}$. After partial trace its instrument is
observationally equivalent to the device's, on arbitrary quantum inputs.

For some absolute constant $c$, the construction, after receipt of the
data, takes at most
\[
 c(d+G+Q_*+1)^4
\]
classical steps, where $G,Q_*$ can be replaced by the explicit upper
bounds from this proof. A uniform data generator's runtime must be
added to this bound. A direct classical controller evaluating the
supplied Boolean circuits during an observed execution takes at most
\[
 O\bigl(T(d+\beta+\Delta+s+k+q+y+1)^4\bigr)
\]
classical steps, apart from the physical quantum operations. These are
deliberately conservative bounds, including binary address processing.
Uniform families therefore remain uniform; polynomially bounded data
generated in polynomial time give polynomially uniform circuits.
\end{theorem}
\begin{proof}
 \step{eqv-fwd-reserve}{There is an isometric initialization of the
 input and all required ancillary factors.}
 \begin{proof}\pf\ Keep the specified live input factors, append
 the other $q$-pool factors, $C,J_t,B_t,Y,E_Y$ and workspace in zero,
 and apply $X$ to the nonzero bits of $c_{\rm in}$. This costs $O(s)$
 gates by $X=HT^4H$. The initialization is an isometry on every input.
 A pool factor with live flag zero has not been acted on before its allocation: it is initially a non-input zero factor, every quantum action requires live flag one, every update other than allocation and sending to $c_{\mathrm{fault}}$ preserves all live flags, allocation changes one zero flag to one, and $c_{\mathrm{fault}}$ is terminal so no allocation follows it. Hence the factor is still zero at allocation, and Lemma~\ref{lem:reserve} justifies moving that allocation to this initialisation.\end{proof}
 \step{eqv-fwd-circuit}{The tick product and output extraction have
 the asserted gate and qubit bounds.}
\begin{proof}
 \step{eqv-fwd-circuit-ordinary}{The $T-1$ ordinary ticks use
 $O((T-1)\mathcal R)$ gates and shared clean workspace.}
 \begin{proof}\pf\ Apply Lemma~\ref{eqv:tick-cost} with separate
 $J_t,B_t$ for every tick. Lemma~\ref{lem:scratch} permits reuse
 of the clean implementation workspace.\end{proof}
 \step{eqv-fwd-circuit-timeout}{The final tick implements
 $F^\tau_{T-1}$ with $O(s+k)$ gates and workspace.}
 \begin{proof}\pf\ Compare status with $\sigma_{\mathrm{act}}$ and
 conditionally replace it by $\sigma_{\mathrm{to}}$, preserving
 all other bits. This is $h_\tau$ on every code, including invalid
 ones, and has Boolean size $O(s)$. XOR-evaluate it into $J_{T-1}$
 and swap $C,J_{T-1}$ by Lemma~\ref{lem:synthesis}. Leave
 $QB_{T-1}$ unchanged. This also covers $T=1$, with no ordinary
 tick. The last records are among the $T$ reserved pairs.\end{proof}
 \step{eqv-fwd-circuit-output}{Output extraction costs
 $O(\Delta+s+y)$ gates and workspace.}
 \begin{proof}\pf\ Compute $o(C)$ into $Y$ with its clean XOR
 evaluator and copy $Y$ into $E_Y$. Lemma~\ref{lem:synthesis}
 bounds the evaluator, and the copy uses $y$ CNOTs. Output ports
 are counted in $\Delta$, so $y\leq\Delta$.\end{proof}
 \step{eqv-fwd-circuit-count}{The retained logical list has exactly
 $q+s+T(s+k)+2y$ factors.}
 \begin{proof}\pf\ Count $Q,C$, the $T$ disjoint pairs $J_t,B_t$,
 and $Y,E_Y$. Implementation workspace is additional and shared
 cleanly by the preceding constructions.\end{proof}
 \qedstep
 \begin{proof}\pf\ Since $T\geq1$, the initialization from
 \stepref{eqv-fwd-reserve} and extraction costs are absorbed in
 $O(T\mathcal R)$. Add the $O(\mathcal R)$ shared workspace
 to \stepref{eqv-fwd-circuit-count}.\end{proof}
\end{proof}
 \step{eqv-fwd-defer}{Tracing the specified environment yields exactly
 the device instrument.}
 \begin{proof}
  \pf\ This is Lemma~\ref{lem:deferred} with the tick costs of
  Lemma~\ref{eqv:tick-cost}; we spell out the expansion once more because
  the gate count is read off it.
  \step{eqv-fwd-expand}{For a pure input $\psi$, the final retained
  vector, before output copying, has form
  $\sum_z A_z\psi\otimes\ket{e_z}$, with orthonormal $e_z$ for
  different sampled-bit lists $z$, padded by zeros at nonmeasurement
  ticks, including the final tick. Here $A_z=B_zE_{\rm in}$, where
  $E_{\rm in}\psi=\psi\otimes\ket0_{Q\setminus L_{\rm in}}$ and
  $B_z$ is the ordered product on $Q$ of selected gates and projectors.}
\begin{proof}
 \step{eqv-fwd-expand-initial}{At $t=0$ the quantum map is
 $E_{\rm in}$, with configuration $c_{\rm in}$ and blank records.}
 \begin{proof}\pf\ This is the initialization in
 \stepref{eqv-fwd-reserve}.\end{proof}
 \step{eqv-fwd-expand-suffices}{\suffices{Verify that every scheduled
 tick has the corresponding action in the finite induction of
 Lemma~\ref{lem:deferred}.}}
 \begin{proof}\pf\ That induction starts with the embedding and
 configuration in \stepref{eqv-fwd-expand-initial} and retains
 precisely the sample and snapshot registers used here.\end{proof}
 \step{eqv-fwd-expand-measure}{\case{An ordinary tick measures
 wire $i$.}}
 \begin{proof}\pf\ The copying CNOT maps $v$ to
 $P_{i,0}v\otimes\ket0+P_{i,1}v\otimes\ket1$.
 Subsequent choices are controlled by deterministically updated $C$;
 old samples and snapshots remain retained.\end{proof}
 \step{eqv-fwd-expand-allocate}{\case{An ordinary tick allocates
 a wire.}}
 \begin{proof}\pf\ Its zero live flag implies this non-input
 factor has never been acted on: every quantum action requires a
 one flag, flags cannot be reset except by fault, and fault is
 terminal. It remains zero, as required for the reserved-pool
 allocation.\end{proof}
 \step{eqv-fwd-expand-gate}{\case{An ordinary tick selects a gate.}}
 \begin{proof}\pf\ Its selected matrix is applied on $Q$ and
 multiplies the ordered product; its sample stays zero and the
 snapshot implements the supplied deterministic update.\end{proof}
 \step{eqv-fwd-expand-other}{\case{An ordinary tick has any remaining
 permitted mode.}}
 \begin{proof}\pf\ The quantum action is identity and the sample
 is zero. Its snapshot performs the specified update. Read mode
 is excluded by hypothesis.\end{proof}
 \step{eqv-fwd-expand-timeout}{\case{The final tick.}}
 \begin{proof}\pf\ Its quantum action is identity, its sample is
 zero, and $c_T=h_\tau(c_{T-1})$, also when $T=1$.\end{proof}
 \qedstep
 \begin{proof}\pf\ These are exactly the finite-induction cases,
 yielding $B_zE_{\rm in}$ and deterministic configuration and record
 bits $e_z$. Distinct lists differ in a retained sample bit, hence
 their basis record vectors are orthogonal and have unit norm.
 \end{proof}
\end{proof}
  \step{eqv-fwd-trace}{The reduced output is
  \[
   \sum_z \tr_{Q\setminus O}(A_z\rho A_z^\dagger)
       \otimes\proj{o(c_T(z))}_Y.
  \]}
  \begin{proof}\pf\ The retained samples erase the cross terms
  between unequal strings under partial trace. Copying the output to
  $E_Y$ gives its orthogonal classical label. Equal output labels are
  summed, not identified coherently. The formula holds for matrix
  units and hence for all density operators; tensoring every step
  with a reference identity changes none of the calculations.
  This also proves the bounded deferred-observation claim here
  without an external measurement theorem.\end{proof}
  \qedstep
  \begin{proof}\pf\ The displayed sum is the definition of sequential
  composition and grouping of the device's instrument.\end{proof}
 \end{proof}
 \step{eqv-fwd-time}{The claimed classical bounds and uniformity hold.}
\begin{proof}
 \step{eqv-fwd-time-emission}{Emitting $G$ gates with scans and
 binary counter arithmetic fits within $O((d+G+Q_*+1)^3)$ steps.}
 \begin{proof}\pf\ For each gate scan the supplied description and
 output so far once. Each emitted gate has
 $O(\log(Q_*+2))$ address bits; the output length is
 $O(G\log(Q_*+2))$. The scans and arithmetic on counters of at most
 $d+G+Q_*+1$ bits cost at most a constant times
 $(d+G+Q_*+1)^2$ per gate. Sum over $G$ gates and include the
 initial input scan.\end{proof}
 \step{eqv-fwd-time-construction}{The complete construction costs
 at most $c(d+G+Q_*+1)^4$ for an absolute $c$.}
 \begin{proof}\pf\ Evaluate or translate a topologically ordered
 Boolean circuit by scanning its description and stored values for
 each gate. Binary counters and comparisons have at most the
 description length. Per tick enumerate at most $O(\mathcal R)$
 flags and gates; locate operands by full scans and emit the gates
 as in \stepref{eqv-fwd-time-emission}. The unary bounds are part
 of $d$ and the explicit gate and wire bounds may be used for
 $G,Q_*$. Thus this fixed work per tick and its repeated scans are
 absorbed by enlarging the absolute constant in the displayed
 fourth power.\end{proof}
 \step{eqv-fwd-time-controller}{The observed controller costs
 $O(T(d+\beta+\Delta+s+k+q+y+1)^4)$ steps.}
 \begin{proof}\pf\ At each observed tick use the same full-scan
 Boolean evaluation and address comparisons on the supplied
 descriptions and current stored values. The last tick uses the
 $O(s)$ circuit for $h_\tau$. The displayed fourth power bounds
 one tick, including binary address processing; sum over $T$ ticks.
 Physical quantum operations are excluded from this classical count.
 \end{proof}
 \step{eqv-fwd-time-uniform}{These constructions preserve uniformity
 and polynomial uniformity with the stated qualifications.}
 \begin{proof}\pf\ Every scan and enumeration has an explicitly
 bounded length and is specified by a single machine. Compose it
 with the data generator and add that generator's runtime. If data
 and their generation are polynomially bounded, the displayed
 construction bound is polynomial as well.\end{proof}
 \step{eqv-fwd-time-adjoint}{Using $\mathcal G$ instead of
 $\mathcal G^{\pm}$ costs at most a factor of seven.}
 \begin{proof}\pf\ $T^\dagger=T^7$ exactly, including phase;
 substitute seven $T$ gates for each $T^\dagger$ and leave every
 other gate unchanged.\end{proof}
 \qedstep
 \begin{proof}\pf\ By \stepref{eqv-fwd-time-construction}--\stepref{eqv-fwd-time-adjoint}.
 \end{proof}
\end{proof}
 \qedstep
 \begin{proof}\pf\ By \stepref{eqv-fwd-reserve},
 \stepref{eqv-fwd-circuit}, \stepref{eqv-fwd-defer}, and
 \stepref{eqv-fwd-time}.\end{proof}
\end{proof}

\begin{theorem}[Circuit to device]\label{eqv:reverse}
Let a finite well-typed circuit over $\mathcal G$ have $g$ gates,
$a$ one-qubit fresh allocations, $n_{\rm in}$ inputs and $b$ designated
terminal computational-basis measurements. Names are never reused;
no gate acts on an inaccessible name. Quantum outputs and measurement
labels are specified. Put
\[
 q=n_{\rm in}+a,\quad N=g+a+b+1,\quad
 w=8+\lceil\log_2(N+q+b+2)\rceil.
\]
There is an observationally equivalent device using primitives 1--9,
with $p=N,T=N+1,h=o_{\max}=B_{\rm mem}=0,L_{\rm slot}=b+1$,
$k=1$, and width $w$. Its $N$ source
instructions are gates, allocations, measurements and stop; the extra
tick performs timeout or terminal padding. Its quantum execution uses
exactly $g$ original gates, $a$ allocations, $b$ measurements and a
pool of $q$ wires. Dropped factors are retained and omitted from the
output, so this is not a simulation of storage reclamation.

For an input description of $d_0$ bits, set
${\mathfrak m}=d_0+N+q+(b+1)w+1$. One can choose
$s=(b+2)w+q+2\leq(b+2)w+q+8$ and
$\beta+\Delta=O(N(s+w+1))\subseteq O({\mathfrak m}^3)$,
construct the device description in $O({\mathfrak m}^6)$
classical steps, and execute its classical controller in
$O((N+1){\mathfrak m}^6)$ steps. Applying Theorem~\ref{eqv:forward} also gives a
fully retained circuit of $O((N+1){\mathfrak m}^3)$ gates and
$O(q+(N+1){\mathfrak m}+{\mathfrak m}^3)$ qubits. All constants are independent of the
circuit instance. The translation preserves uniformity and polynomial
uniformity.
\end{theorem}
\begin{proof}
 \step{eqv-rev-drop}{Replace every drop by retaining an inaccessible
 factor, with no later operation on that factor.}
 \begin{proof}\pf\ For an operator $X$ on $LE$ and a later map
 $\Phi$ on $L$, expansion in matrix units on $E$ gives
 $\tr_E((\Phi\otimes\id_E)(X))=\Phi(\tr_E X)$.
 No later operation acts on a dropped factor: drop removes a name permanently and no gate acts on an inaccessible name.
Iterating this identity moves all partial traces to the output.
 On all retained factors the circuit is an isometry before any
 terminal observation, with measurement dilation given by copying
 into fresh outcome factors.\end{proof}
 \step{eqv-rev-device}{The literal instruction list has the stated
 bounded implementation with a total decoder and update.}
\begin{proof}
 \step{eqv-rev-device-layout}{\pick\ the following instruction list,
 configuration fields, validity predicate and initial code.}
 \begin{proof}
  \step{eqv-rev-device-list}{Number the gates, allocations and measurements
  in their given order as $0,\ldots,N-2$, and append stop at $N-1$.}
  \begin{proof}\pf\ There are $g+a+b=N-1$ operations before stop,
  so this gives exactly $p=N$ instruction positions.\end{proof}
  \step{eqv-rev-device-fields}{Use $b+1$ word slots, one $w$-bit
  position, $q$ live flags and two status bits, with $00$ active,
  $01$ stopped, $10$ fault and $11$ timeout. All other fields have
  width zero; literal quantum identifiers reside in the instructions.}
  \begin{proof}\pf\ The first $b$ words hold measurement results;
  the last word is the outcome, with no duplicated outcome field.
  These fields have total width $(b+2)w+q+2$. The specified $w$
  accommodates the instruction positions and literal identifiers.
  Memory, stack and console capacities are zero, as asserted.\end{proof}
  \step{eqv-rev-device-valid}{Let $c_{\rm fault}$ have status $10$
  and every other bit zero. A code is valid precisely when its
  position is below $N$ and its status is not $10$, or it equals
  $c_{\rm fault}$. Slots and flags have no further range restrictions.}
  \begin{proof}\pf\ This is a predicate on every configuration
  string and includes $c_{\rm fault}$. Among valid codes, status
  $00$ is active; statuses $01$, $11$, and the single code
  $c_{\rm fault}$ are terminal. These alternatives partition the
  valid codes.\end{proof}
  \step{eqv-rev-device-initial}{Initially position, slots and status
  are zero, and exactly the input live flags are one.}
  \begin{proof}\pf\ Since $N\geq1$, position zero is below $N$;
  status $00$ makes this a valid active configuration with the
  prescribed input live list.\end{proof}
  \qedstep
  \begin{proof}\pf\ The preceding explicit choices give the
  required finite data.\end{proof}
 \end{proof}
 \step{eqv-rev-device-rules}{\assume{$c$ is any configuration code
 and $z$ is any one-bit sample.}\prove{The following decoder and
 update rules are defined and meet the all-code conditions.}}
 \begin{proof}
  \step{eqv-rev-device-invalid}{\case{$c$ is invalid.}}
  \begin{proof}\pf\ Emit canonical fault and set $f(c,z)=c_{\rm fault}$
  for every $z$. This gives the required fault behavior without
  a quantum action.\end{proof}
  \step{eqv-rev-device-terminal}{\case{$c$ is valid and terminal.}}
  \begin{proof}\pf\ Emit canonical idle and preserve the entire
  code, independently of $z$, as required for terminal padding.\end{proof}
  \step{eqv-rev-device-active}{\case{$c$ is active.}}
  \begin{proof}
   \step{eqv-rev-device-select}{Compare position with the instruction
   literals and check the selected instruction's operands.}
   \begin{proof}\pf\ Validity makes exactly one position literal
   match. Allocation requires its flag zero; gates and measurements
   require all operand flags one. CNOT operands are distinct by the
   supplied circuit's typing. A failed check emits canonical fault
   and sets $f(c,z)=c_{\rm fault}$ for every $z$.\end{proof}
   \step{eqv-rev-device-allocate}{\case{The selected allocation passes
   its check.}}
   \begin{proof}\pf\ Emit its allocation code, set just its zero
   live flag to one and advance position. The quantum action is
   identity on its reserved factor; all other flags are preserved
   and the update ignores $z$.\end{proof}
   \step{eqv-rev-device-gate}{\case{The selected gate passes its check.}}
   \begin{proof}\pf\ Emit the gate and its live operands and
   advance position. Preserve all flags and ignore $z$.\end{proof}
   \step{eqv-rev-device-measure}{\case{The selected measurement passes
   its check.}}
   \begin{proof}\pf\ Emit its first address only and sample into
   $B_t[0]$. Write sample bit zero into the low bit of its designated
   word, zero that word's other bits and advance position. Preserve
   all live flags. The destination is the slot for the measurement's
   designated output order and belongs solely to $f$.\end{proof}
   \step{eqv-rev-device-stop}{\case{The selected instruction is stop.}}
   \begin{proof}\pf\ Emit ordinary mode, set status $01$, set the
   outcome word to zero and retain position and all live flags.
   The update ignores $z$.\end{proof}
   \qedstep
   \begin{proof}\pf\ A failed check and the four successful
   instruction cases exhaust the active codes. Advancing position
   occurs only before the final stop position, hence stays below $N$.
   Each nonallocation update preserves all flags unless it produces
   $c_{\rm fault}$.\end{proof}
  \end{proof}
  \qedstep
  \begin{proof}\pf\ Invalid, valid terminal and active codes exhaust
  all strings, by \stepref{eqv-rev-device-layout}.\end{proof}
 \end{proof}
 \step{eqv-rev-device-canonical}{The decoder outputs are canonical
 and read mode never occurs.}
 \begin{proof}\pf\ Use the prescribed 3-bit mode, 2-bit gate tag
 ($0,1,2$ for $H,T,\operatorname{CNOT}$), $a_q$-bit addresses
 and $a_k$-bit length, with $k=1$. Set every unused address, gate,
 length and read field to zero. The cases in
 \stepref{eqv-rev-device-rules} use only the listed canonical modes
 with legal live operands; none selects read.\end{proof}
 \step{eqv-rev-device-output}{The output circuit and timeout have
 the specified form.}
 \begin{proof}\pf\ Set $o(c)$ to the low bits of the first $b$
 word slots in their designated order, on every code. The outcome
 word is internal and not part of $Y$. Follow the $N$ ordinary
 ticks by the separate $h_\tau$ snapshot with
 $\sigma_{\mathrm{act}}=00$, $\sigma_{\mathrm{to}}=11$;
 this gives $T=N+1$ exactly.\end{proof}
 \qedstep
 \begin{proof}\pf\ The explicit tests, bit selections and constant
 assignments above are total Boolean circuits. Their cases verify
 all global conditions of Definition~\ref{def:device}.\end{proof}
\end{proof}
 \step{eqv-rev-size}{The advertised finite costs hold.}
\begin{proof}
 \step{eqv-rev-size-width}{The layout has $s=(b+2)w+q+2=O({\mathfrak m})$
 bits, with $s\geq a_q$ and $N,w=O({\mathfrak m})$.}
 \begin{proof}\pf\ Count the fields in \stepref{eqv-rev-device}.
 The definition of ${\mathfrak m}$ includes $(b+1)w$, $q$ and $N$; also
 $w\leq(b+1)w$. The specified logarithmic width bounds $a_q$.\end{proof}
 \step{eqv-rev-size-flags}{The validity test costs $O(s+w)$ and
 the literal selection tests cost $O(Nw)$.}
 \begin{proof}\pf\ Test position below $N$, status and equality
 with $c_{\rm fault}$ by linear-size bit comparisons. Each position
 equality costs $O(w)$; that literal's live tests cost $O(1)$.
 Compute the disjoint selection flags once, including the invalid
 and terminal cases.\end{proof}
 \step{eqv-rev-size-circuits}{Including ports,
 $\beta+\Delta=O(N(s+w+1))\subseteq O({\mathfrak m}^3)$.}
 \begin{proof}\pf\ For each of the $s$ update output bits, select
 the preserved bit, specified literal or sampled bit by ANDs with
 the flags and their XOR sum. This costs $O(Ns)$; invalid cases
 select $c_{\rm fault}$ and terminal cases copy $c$. Select the
 decoder's literal canonical codes of width $6+3a_q=O(w+1)$
 in the same manner. Output extraction copies $b$ low sample bits.
 Add the tests in \stepref{eqv-rev-size-flags} and the ports.
 In fact $N,s,w=O({\mathfrak m})$ gives $O({\mathfrak m}^2)$ Boolean gates.\end{proof}
 \step{eqv-rev-size-time}{Emission costs $O({\mathfrak m}^6)$ classical steps;
 evaluation costs $O((N+1){\mathfrak m}^6)$.}
 \begin{proof}\pf\ Emit one Boolean gate at a time, scanning the
 supplied description and output so far and doing binary counter
 arithmetic. There are $O({\mathfrak m}^2)$ gates with
 $O(\log({\mathfrak m}+2))$-bit addresses. Each full scan and its counter
 arithmetic cost $O({\mathfrak m}^2\log({\mathfrak m}+2)+{\mathfrak m}^2)$, so complete emission
 fits in $O({\mathfrak m}^6)$. The same repeated-scan procedure evaluates the
 circuits at each of $N+1$ ticks within $O((N+1){\mathfrak m}^6)$ steps.
 No enumeration of $2^s$ configurations is used.\end{proof}
 \step{eqv-rev-size-retained}{The fully retained bounds in the
 theorem follow from Theorem~\ref{eqv:forward}.}
 \begin{proof}\pf\ Here $k=1$, $y=b$, $q,s=O({\mathfrak m})$ and
 $\beta+\Delta=O({\mathfrak m}^3)$, so
 $\mathcal R=\beta+\Delta+s+k+(q^2+kq+1)(s+k+1)=O({\mathfrak m}^3)$.
 Substitute $T=N+1$ into the forward theorem: gates are
 $O((N+1){\mathfrak m}^3)$ and qubits are $O(q+(N+1){\mathfrak m}+{\mathfrak m}^3)$.\end{proof}
 \qedstep
 \begin{proof}\pf\ By \stepref{eqv-rev-size-width}--\stepref{eqv-rev-size-retained}.
 \end{proof}
\end{proof}
 \step{eqv-rev-equality}{The output instruments agree.}
 \begin{proof}\pf\ Induction over the literal instructions shows
 that each allocated factor is zero, each gate has the same ordered
 operands and matrix, and each sampled bit is recorded in its
 designated slot. Induction on the instructions maintains the invariant that the live flags are one exactly on the input names together with the names allocated so far: fresh names are never reused, so an allocation flag is zero, and well-typedness gives every gate and measured wire among those names. Hence no check fails and fault is unreachable from the specified initial configuration. Stop precedes timeout. Classical
 output extraction selects exactly the requested labels. The
 retained snapshots may differ from the circuit's own environment,
 but Lemma~\ref{lem:deferred} and \stepref{eqv-rev-drop} give equality
 of the designated instruments, also with references.\end{proof}
 \qedstep
 \begin{proof}\pf\ By \stepref{eqv-rev-device},
 \stepref{eqv-rev-size}, and \stepref{eqv-rev-equality}.
 The explicitly bounded transformations are single classical
 algorithms, so composing them with a uniform generator proves
 the uniformity assertions.\end{proof}
\end{proof}

\begin{corollary}[Coverage of bounded adaptive circuits]\label{eqv:adaptive-circuits}
Let a bounded adaptive controller be presented directly as device data satisfying Definitions~\ref{def:device} and \ref{def:adaptive}: explicit total decoder, update and output circuits, a fixed horizon $T\ge1$ including the separate final timeout, and a fixed output list live at final time for every reachable outcome list. Assume
a valid specified initial state with exactly the input wires live,
and total controller data obeying the canonical decoder, all-code
live-flag discipline, terminal padding, and all-code exclusion of read
mode in Definitions~\ref{def:device} and \ref{def:adaptive}.
Use fresh blank records and trace the entire complement of $OY$.
These circuit instruments and the
bounded devices are equivalent in both directions, with polynomial
overhead in the horizon, pool size, classical capacity, and supplied
Boolean-circuit sizes. Purifications are retained before taking outputs.
Gates in $\mathcal G^{\pm}$ are admitted after exact $T^\dagger=T^7$
expansion and adjustment of the horizon. Ideal rotations instead require
the separate approximation theorem and its certificate-generation costs.
This corollary concerns formal device data; it does not identify an
unspecified source-language executor with those data.
\end{corollary}
\begin{proof}
 \step{eqv-ac-encode}{The supplied controller has finite device data.}
 \begin{proof}\pf\ These are the supplied total decoder, update and output circuits in the hypothesis, with their canonical encoding, initial live set, horizon and separate final timeout. Definitions~\ref{def:device} and \ref{def:adaptive} therefore apply directly.\end{proof}
 \step{eqv-ac-translate}{Both required simulations follow.}
 \begin{proof}\pf\ Theorem~\ref{eqv:forward} gives the circuit simulation of these data. Theorem~\ref{eqv:reverse} gives the device simulation of every finite circuit over $\mathcal G$ with terminal observations, and composing it with the device presentation just given yields the stated two-direction equivalence, with inaccessible factors retained until the final partial trace.\end{proof}
 \qedstep
 \begin{proof}\pf\ Compose the polynomial bounds in
 Theorems~\ref{eqv:forward} and \ref{eqv:reverse}. This statement does
 not give a succinct encoding to an exponentially large explicitly
 supplied schedule. 
 \end{proof}
\end{proof}

\subsection{Why the restrictions cannot be suppressed}

\begin{proposition}[Finite capacity and stopping restrictions]
\label{eqv:restrictions}
An instance with fixed console capacity $o_{\max}$ cannot implement the family
whose output tape is $n$ zero bytes for every $n$. Width-$w$ modular
addition is not unbounded integer addition. Primitives 1--9 alone
do not ensure termination, and no horizon obtained by truncating a given unbounded execution and reporting timeout need preserve that execution's instrument. Defining a timeout is therefore an
additional specification, unless a sufficient bound has been proved.
\end{proposition}
\begin{proof}
 \step{eqv-rest-capacity}{The output-tape family exceeds fixed capacity.}
 \begin{proof}\pf\ At $n=o_{\max}+1$ its required length exceeds $o_{\max}$;
 Definition~\ref{def:device} requires overflow behavior instead. Independently, a fixed instance has a fixed output circuit $o$ of fixed output width, so the same conclusion holds for $n$ larger than that width in bits.
 Growing instances can reserve $n$ bytes.\end{proof}
 \step{eqv-rest-width}{The two addition operations differ.}
 \begin{proof}\pf\ On $2^w-1$ and $1$, modular addition returns
 $0$, whereas integer addition returns $2^w$. The theorem preserves
 whichever total function is supplied, not an unspecified alternative.
 Multiword arithmetic is permitted only with its memory and Boolean
 cost included.\end{proof}
 \step{eqv-rest-time}{There need not be a sufficient finite horizon.}
 \begin{proof}\pf\ A conditional jump with an always true condition
 back to itself never reaches stop. For a terminating-with-probability-one
 example, initialize a qubit to zero, apply $H$, measure, stop on $1$,
 and on $0$ repeat $H$ and measurement on the now-zero qubit. After
 $r$ trials the probability of no stop is $2^{-r}>0$, by multiplication
 of the conditional probabilities $1/2$. A cutoff reporting timeout
 has that positive extra outcome, while the unbounded run stops
 with probability one since $2^{-r}\to0$. The conclusion concerns
 the cutoff rule; a different circuit may reproduce a particular
 unbounded program's final output by other means.\end{proof}
 \qedstep
 \begin{proof}\pf\ By \stepref{eqv-rest-capacity},
 \stepref{eqv-rest-width}, and \stepref{eqv-rest-time}.\end{proof}
\end{proof}

\begin{counterexample}[Exact finite-alphabet equivalence cannot include
arbitrary angles]\label{eqv:exact-obstruction}
Start with $\ket0$, apply $H R_z(\theta)H$, and measure, where
$\theta=2\arccos(1/\sqrt3)$ is a computable real with this finite
description. The probability of outcome zero is $1/3$. No finite
bounded adaptive circuit over $\mathcal G$, with every quantum wire
initially in a computational-basis state, zero-initialized allocations,
deterministic classical data, and computational-basis measurements and with output probabilities given by the unconditional instrument (no postselection or conditional renormalisation),
has this same output distribution, even allowing auxiliary qubits and
partial traces. No additional quantum state or random classical seed
is supplied as input.
\end{counterexample}
\begin{proof}
 \step{eqv-ring-amplitudes}{Every amplitude of an unnormalized vector
 conditional on a finite measurement string belongs to
 $\mathcal A=\mathbb Z[1/\sqrt2,i]$.}
 \begin{proof}\pf\ Initial basis amplitudes are integers. Entries
 of $H,T,\operatorname{CNOT}$, basis projectors, and zero embeddings
 belong to this ring. Finite matrix multiplication and addition
 preserve the ring. For each fixed measurement string adaptivity
 selects one such finite product, without conditional normalization.
\end{proof}
 \step{eqv-ring-probabilities}{Every rational output probability of
 such a circuit is dyadic.}
 \begin{proof}\pf\ An element of $\mathcal A$ has form
 $(a+b\sqrt2+i(c+d\sqrt2))/2^m$ with integer $a,b,c,d$ and $m\geq0$.
 Its squared modulus belongs to $\mathbb Z[1/2,\sqrt2]$. Sum these
 squared moduli over unobserved basis values and the finite set of
 measurement strings assigned to the output. This still gives
 $(u+v\sqrt2)/2^j$ for integers $u,v,j$, $j\geq0$. If it is rational,
 irrationality of $\sqrt2$ implies $v=0$. The rational number is
 therefore dyadic.\end{proof}
 \step{eqv-ring-ideal}{The stated ideal circuit gives a nondyadic
 probability.}
 \begin{proof}\pf\ Its zero amplitude is $\cos(\theta/2)$, so its
 zero probability is $\cos^2(\theta/2)=1/3$. If $1/3=u/2^j$ then
 $3u=2^j$, impossible since $3$ does not divide a power of $2$.
\end{proof}
 \qedstep
 \begin{proof}\pf\ The exact equality of instruments would imply
 equality of this probability, contradicting
 \stepref{eqv-ring-probabilities} and \stepref{eqv-ring-ideal}.
 This refutes the unrestricted exact claim, not
 the finite-alphabet equivalence of Theorems~\ref{eqv:forward} and \ref{eqv:reverse}.\end{proof}
\end{proof}

\subsection{Phase-sensitive approximation}

\begin{lemma}[From retained isometry error to output error]
\label{eqv:error-lemma}
Let $V,W$ be isometries with the same finite input and retained output
spaces. Apply to their outputs the same channel made of isometric
initializations, unitary gates, classical-output copying and partial
traces. The resulting channels $\Phi_V,\Phi_W$ satisfy
$d_{\rm out}(\Phi_V,\Phi_W)\leq2\|V-W\|$.
\end{lemma}
\begin{proof}
 \step{eqv-error-product}{For every input state $\rho$ including a
 reference, the retained density operators differ in trace norm by
 at most $2\|V-W\|$.}
 \begin{proof}\pf\ Suppress identity factors on the reference and
 write their difference as
 $(V-W)\rho V^\dagger+W\rho(V^\dagger-W^\dagger)$.
 For matrices $A,X,B$, $\|AXB\|_1\leq\|A\|\|X\|_1\|B\|$:
 expand $X$ in singular-value rank-one terms and use the triangle
 inequality and $\|uv^\dagger\|_1=\|u\|\|v\|$.
 Now $\|V\|=\|W\|=\|\rho\|_1=1$, and tensoring an operator
 with an identity leaves its operator norm unchanged.\end{proof}
 \step{eqv-error-trace}{The prescribed common output operations
 cannot increase that trace norm.}
 \begin{proof}\pf\ Unitary conjugation and isometric embedding
 preserve singular values apart from added zeros. For Hermitian
 $X$, the identity
 $\|X\|_1=\max_{-I\leq B\leq I}\tr(BX)$ follows by diagonalizing
 $X$ and choosing its sign matrix. Consequently
 $\|\tr_E X\|_1\leq\|X\|_1$, since each admissible output $B$
 lifts to admissible $B\otimes I_E$. Copying a classical output
 before tracing is an isometric embedding followed by a unitary.
 Apply these facts successively, also with a retained reference.
\end{proof}
 \qedstep
 \begin{proof}\pf\ Take the supremum in Definition~\ref{def:registers}
 after \stepref{eqv-error-product} and \stepref{eqv-error-trace}.
\end{proof}
\end{proof}

\begin{theorem}[Quantified approximation with the stated certificates]
\label{eqv:approx}
Under exactly the certificate hypotheses of
Definition~\ref{def:certificates}, take $N_g\geq1$ unitary factors on the
common retained logical register. For every factor $A_i$ and rational
$0<\eta<1$, assume a uniform procedure supplies a clean implementation
of a logical unitary $\widetilde A_i$ with actual-matrix error at most
$\eta$ and length at most $c_*\log^\alpha(2+1/\eta)$, where
$c_*>0,\alpha\geq0$ are fixed independently of $i,N_g,\eta$. Every controlled
factor is included as its entire controlled matrix, and every repeated
or inverse occurrence is counted. Let $\epsilon$ be rational with
$0<\epsilon<1$ and set $\eta=\epsilon/(2N_g)$. Then the output error and gate count satisfy:
\[
 d_{\rm out}(\Phi,\widetilde\Phi)\leq\epsilon,
 \qquad G_{\rm app}=O\bigl(N_g\log^\alpha(2+N_g/\epsilon)\bigr).
\]
The retained isometry error is at most $\epsilon/2$ on the specified
zero-workspace input subspace. If the logical retained width is $n_L$
and the largest certificate workspace is $a_*$, sufficient width is
$n_L+a_*$. Unused workspace may be removed, giving
$a_*=O(\log^\alpha(2+N_g/\epsilon))$ for gate-list certificates.

Let $\tau_i(\eta)$ be the time taken to produce the explicit certificate
circuit for occurrence $i$, and $d_1$ the bit length of the expanded
logical description and request. A classical construction-time bound is
\[
 \sum_{i=1}^{N_g}\tau_i(\eta)
   +O\bigl((d_1+G_{\rm app}+n_L+a_*+1)^4\bigr).
\]
No polynomial bound on the first term follows merely from uniformity.
Whenever a factor is used in reverse, reverse the same chosen
approximating circuit, rather than choose an unrelated approximation.
Repeated use is counted with multiplicity in $N_g$.
\end{theorem}
\begin{proof}
 \step{eqv-app-product}{The products of logical unitaries differ by
 at most $N_g\eta$ in operator norm.}
 \begin{proof}\pf\ Index the factors by their occurrences in the expanded list; where an occurrence is the inverse of an earlier chosen factor, take for it the adjoint of that same approximating circuit, so the hypothesis is applied once per distinct factor and adjunction supplies the inverse occurrences. For $U=A_{N_g}\cdots A_1$ and
 $\widetilde U=\widetilde A_{N_g}\cdots\widetilde A_1$, insert and
 subtract the intermediate products replacing one factor at a time.
 Each resulting summand is a product of unitary factors and one
 difference $\widetilde A_i-A_i$. Submultiplicativity and the
 triangle inequality bound the sum by $N_g\eta$.
 Also $\|\widetilde A_i^\dagger-A_i^\dagger\|
 =\|\widetilde A_i-A_i\|$, so the same estimate applies to every
 inverse occurrence. Controlled factors are covered by the expressly
 required bound on their entire actual matrix.\end{proof}
 \step{eqv-app-space}{The physical circuit has this action on its
 zero-workspace subspace and the claimed space and gate costs.}
 \begin{proof}\pf\ Lemma~\ref{lem:scratch} permits all clean
 implementations, including adjoints, to reuse workspace of the
 maximum width. Their restriction is precisely $\widetilde U$;
 no bound on arbitrary nonzero workspace states is asserted.
 Add the $N_g$ length bounds. Every elementary gate touches at most
 two wires, so a certificate of length $l$ needs at most $2l$
 distinct touched ancillary wires. Delete untouched ancillary wires
 before taking the maximum.\end{proof}
 \step{eqv-app-output}{The final error is at most $\epsilon$.}
 \begin{proof}\pf\ The initialized ideal and approximating
 isometries differ by at most $N_g\eta=\epsilon/2$ by
 \stepref{eqv-app-product} and \stepref{eqv-app-space}.
 Lemma~\ref{eqv:error-lemma} applies to the same specified output
 extraction and trace.\end{proof}
 \step{eqv-app-time}{The classical-time assertion holds.}
 \begin{proof}\pf\ Put ${\mathfrak m}=d_1+G_{\rm app}+n_L+a_*+1$. After certificate production, the emitting machine writes one gate at a time, scanning the supplied description and output so far once and doing binary counter arithmetic on at most ${\mathfrak m}$ bits. Each scan and arithmetic cost at most $O({\mathfrak m}^2)$ steps, so $G_{\rm app}\le {\mathfrak m}$ gates cost $O({\mathfrak m}^3)$ steps after an initial scan; the displayed fourth power also covers relabelling and reversal.
Run each supplied uniform procedure, then
 concatenate, relabel and reverse its finite gate lists where
 required. Full scans and binary counter arithmetic have the
 displayed conservative polynomial cost, as in
 Theorem~\ref{eqv:forward}. The hypotheses place no growth bound
 on the procedures' individual runtimes $\tau_i$.\end{proof}
 \qedstep
 \begin{proof}\pf\ By \stepref{eqv-app-space},
 \stepref{eqv-app-output}, and \stepref{eqv-app-time}.
 An irrational error tolerance admits a smaller rational request;
 uniform input requests themselves are rational by
 Definition~\ref{def:rtm}.\end{proof}
\end{proof}

\begin{counterexample}[Removing phases before a control is invalid]
\label{eqv:phase-obstruction}
The one-qubit matrices $I$ and $-I$ have zero distance modulo scalar
phase. Their singly controlled matrices, on a control initialized to
$\ket+=(\ket0+\ket1)/\sqrt2$ and any target, give respectively
$\ket+$ and $\ket-=(\ket0-\ket1)/\sqrt2$ on the control. Their
output channels have $d_{\rm out}=2$.
\end{counterexample}
\begin{proof}
 \step{eqv-phase-block}{The controlled matrices are $I$ and
 $Z\otimes I$.}
 \begin{proof}\pf\ Insert $I$ and $-I$ into
 $\proj0\otimes I+\proj1\otimes A$.\end{proof}
 \step{eqv-phase-distance}{Their outputs have trace distance $2$.}
 \begin{proof}\pf\ The two displayed control vectors are
 orthogonal. The difference of their projectors has eigenvalues
 $1,-1$ and hence trace norm $2$. The distance between any two
 states is at most $2$ by the triangle inequality.\end{proof}
 \qedstep
 \begin{proof}\pf\ By \stepref{eqv-phase-block} and
 \stepref{eqv-phase-distance}.\end{proof}
\end{proof}

\subsection{A local-source synthesis guarantee with actual phases}
Everything from here through Theorem~\ref{eqv:rotations} concerns only
the optional ideal rotations; the exact equivalence above is already
complete without it. The care below is needed because phases discarded
before adding a control change the controlled operation
(Counterexample~\ref{eqv:phase-obstruction}).


\begin{lemma}[An inverse-closed determinant-one instruction set]
\label{eqv:su2-set}
Define the following actual one-qubit matrices, with $S=T^2$:
\[
 V=SHTHS^\dagger,\quad J=VSV^\dagger,\quad
 A=T^\dagger HTH,\quad B=J^\dagger A J,
 \quad\mathcal F=\{A,A^\dagger,B,B^\dagger\}.
\]
This is a fixed finite inverse-closed subset of $SU(2)$ whose generated
subgroup is dense in $SU(2)$, and each element has a constant-length exact-phase
implementation over $\mathcal G^{\pm}$ without workspace.
\end{lemma}
\begin{proof}
 \step{eqv-su-matrices}{These are the two determinant-one rotations
 in Boykin et al., Section 3, equations (5)--(6), PDF p.~4.}
 \begin{proof}\pf\ Direct multiplication, also used in
 Lemma~\ref{lem:synthesis}, gives $VZV^\dagger=H$ and hence
 $J^2=H$. Algebraically, $\sigma_x^{1/4}=H\sigma_z^{1/4}H=HTH$ and $\sigma_z^{1/2}=S$ under the cited power definitions, so $V=SHTHS^\dagger=\sigma_z^{1/2}\sigma_x^{1/4}\sigma_z^{-1/2}=\sigma_y^{1/4}$ and therefore $J=VSV^\dagger=\sigma_y^{1/4}\sigma_z^{1/2}\sigma_y^{-1/4}$, which is the cited definition of $H^{1/2}$. More precisely $J$ acts by $1$ and $i$ on the $+1$ and
 $-1$ eigenspaces of $H$, respectively. Thus $J$ is exactly the
 paper's $H^{1/2}$ under its spectral convention in Section 2
 \cite[PDF pp.~3--4, power identities preceding Section~3]{BoykinEtAl1999}.
 Its $\sigma_z^{1/4}$ is $T$ and its $\sigma_x^{1/4}$ is $HTH$.
 Substitution in the cited equations gives $A$ and $B$, with no
 omitted scalar multiplier.\end{proof}
 \step{eqv-su-density}{The group generated by $A,B$ is dense in
 $SU(2)$.}
 \begin{proof}\pf\ The local proof in
 \cite[Section 3, PDF p.~4, equations (7)--(12)]{BoykinEtAl1999}
 gives these two matrices as rotations through an irrational multiple
 of $\pi$ about orthogonal axes, and expresses each $SU(2)$ matrix
 as three rotations about those axes. Powers approximate each of
 the three rotations by equations (10)--(11); if the three approximating powers are within $\epsilon/3$ of the three factors in operator norm, the telescoping estimate for a product of three unitaries puts their product within $\epsilon$ of the target. This is the
 precise density assertion used here.\end{proof}
 \step{eqv-su-hypotheses}{All other instruction-set hypotheses hold.}
 \begin{proof}\pf\ The matrices are unitary products, and
 $\det A=e^{-i\pi/4}(-1)e^{i\pi/4}(-1)=1$;
 conjugation gives $\det B=1$. Adjoints have determinant one too.
 Inverses are explicitly included. Each displayed word and its
 adjoint has constant length; $S^{\pm1}=T^{\pm2}$ and $H^\dagger=H$.
\end{proof}
 \qedstep
 \begin{proof}\pf\ By \stepref{eqv-su-density} and
 \stepref{eqv-su-hypotheses}. This application uses actual
 determinant-one matrices, not a claim that the original
 $H,T,\operatorname{CNOT}$ already satisfy $SU(d)$ hypotheses.
\end{proof}
\end{proof}

\begin{theorem}[Certified approximation for the optional rotations]
\label{eqv:rotations}
For $U\in SU(2)$ given by an effective entry-approximation procedure,
and rational $0<\eta<1$, an exact-phase word over $\mathcal G^{\pm}$
approximates $U$ within $\eta$ in operator norm, with length
$O(\log^4(2+1/\eta))$. A uniform procedure finds such a word.
The same statement holds for its singly controlled matrix with
constant clean workspace. In particular it supplies the certificates
of Theorem~\ref{eqv:approx} for the optional $R_z(\theta)$ gates when each such gate is supplied with the dyadic entry-approximation procedure hypothesised above for its matrix, their adjoints and their singly controlled
versions. A fixed bound on the number of controls can also be allowed,
with constants depending on that bound.

For completeness a conservative explicit classical-time guarantee is
available without assuming a unit-cost real-number model. Encode
$\eta=u/v$ by positive binary integers $0<u<v$ separated by delimiters;
reduction to lowest terms is not required. Let $r_\eta$ be the full
request length. For this bit-time bound require the entry procedure to
return dyadic approximations $\tilde U$ of the four complex entries with $|\tilde U_{ij}-U_{ij}|\le2^{-b}$ in complex modulus for each entry,
using $O(b)$ bits per entry. Let $\mu_U(b)$ bound its full multitape-TM
cost, including reading $U$ and writing those approximations; conversion
from any other entry representation must be charged in $\mu_U$.
There are absolute constants such that, putting
\[
 m=\min\{j\geq1:2^{-j}\leq\eta\},\quad
 \delta=2^{-m},\quad b=m+10,\quad
 L_{\max}=\lceil C_{\mathcal F}(m+4)^4\rceil,
\]
where $C_{\mathcal F}$ is a fixed constant for $\mathcal F$, exhaustive certified
search takes
\[
 O\bigl((r_\eta+1)^3+\mu_U(b)+4^{L_{\max}+1}(L_{\max}+b+1)^6\bigr)
\]
classical steps and emits a word of length $O(L_{\max})$.
Here $\delta\leq\eta<2\delta$, so
$L_{\max}=O(\log^4(2+1/\eta))$. Arbitrarily long binary rational requests
are allowed and their processing cost is included.
This time bound is not claimed polynomial in $\log(1/\eta)$.
No runtime bound on $\mu_U$ follows merely from computability of
the angle. General effectively specified $U(d)$ matrices with arbitrary
scalar phases are not included in this theorem.
\end{theorem}
\begin{proof}
 \step{eqv-rot-sk}{The word-length guarantee applies with actual
 operator norm and exponent at most $4$.}
\begin{proof}
 \step{eqv-rot-sk-apply}{Dawson--Nielsen's length theorem applies
 to $\mathcal F$ at dimension $2$ and accuracy $\delta/8$.}
 \begin{proof}\pf\ Lemma~\ref{eqv:su2-set} supplies a finite,
 inverse-closed instruction set dense in $SU(2)$. These are
 Definition 1's hypotheses for Theorem 1
 \cite[PDF pp.~2, 4]{DawsonNielsen2006}. Both the dimension and
 instruction set are fixed, so the length bound is uniform in $U$.
 \end{proof}
 \step{eqv-rot-sk-norm}{The guarantee is in actual operator norm,
 with exponent less than $4$.}
 \begin{proof}\pf\ The norm convention on PDF p.~3 is operator
 norm without identifying scalar phases. Section 3, equation (8),
 PDF p.~7 gives exponent $\log5/\log(3/2)\simeq3.97<4$.
 \end{proof}
 \step{eqv-rot-sk-bound}{\pick\ a fixed constant $C_{\mathcal F}$ giving
 $L_{\max}=\lceil C_{\mathcal F}(m+4)^4\rceil$ as an upper bound for every $m\geq1$.}
 \begin{proof}\pf\ The cited asymptotic bound supplies a uniform
 constant at sufficiently small accuracies. At any fixed larger
 accuracy, density gives open balls centered at words covering
 $SU(2)$; compactness gives a finite subcover and hence a finite
 maximum word length. Enlarge the constant to cover the remaining
 accuracies. Since $\delta/8=2^{-(m+3)}$, the displayed fourth
 power bounds all these lengths.\end{proof}
 \step{eqv-rot-sk-expand}{Elementary expansion preserves the norm
 guarantee and gives length $O(L_{\max})$.}
 \begin{proof}\pf\ Each of the finitely many $\mathcal F$ letters
 expands into its fixed exact-phase word. Substitution leaves the
 product matrix unchanged and multiplies length by at most a fixed
 constant.\end{proof}
 \qedstep
 \begin{proof}\pf\ By \stepref{eqv-rot-sk-apply}--\stepref{eqv-rot-sk-expand}.
 \end{proof}
\end{proof}
 \step{eqv-rot-search}{There is a certified finite search with the
 advertised bit cost.}
 \begin{proof}
  \step{eqv-rot-request}{Computing $m,\delta,b,L_{\max}$ costs
  $O((r_\eta+1)^3)$ steps.}
  \begin{proof}\pf\ Read the entire delimited request. Test
  $2^j u\geq v$ for successive $j$, starting at $1$, by binary
  shifts and comparisons. The first successful $j$ is $m$ and
  $m\leq r_\eta$. Even full scans and grade-school integer
  arithmetic for the fixed fourth power fit the displayed bound.
  Minimality gives $\delta\leq\eta<2\delta$.
  \end{proof}
  \step{eqv-rot-candidates}{Each candidate has a certified distance
  upper bound within $\delta/8$ of its true Frobenius distance,
  at cost $O((L_{\max}+b+1)^6)$.}
\begin{proof}
 \step{eqv-rot-candidates-exact}{A candidate of at most $L_{\max}$ letters
 has an exact representation with $O(L_{\max}+1)$-bit coefficients.}
 \begin{proof}\pf\ Enumerate words in $\mathcal F$ by increasing
 length. Entries of their elementary expansions lie in
 $\mathbb Z[1/\sqrt2,i]$, represented by four integer coefficients
 and a power-of-two denominator. Multiplication by a fixed letter
 uses a fixed number of integer operations and increases bit length
 by at most a constant. Choose a fixed $c_0$ so all coefficient
 lengths are at most $c_0L_{\max}+c_0$.\end{proof}
 \step{eqv-rot-candidates-matrices}{There are rational approximations
 $\widetilde W,\widetilde U$ with certified Frobenius errors
 $e_W\leq2^{-b-2}$ and $e_U\leq2^{1-b}$.}
 \begin{proof}\pf\ Bisect rationally to approximate $\sqrt2$ with
 error at most $2^{-(c_0L_{\max}+b+c_0+4)}$. Multiplication by any coefficient
 contributes at most $2^{-b-4}$. In the four-coefficient representation
 each real or imaginary component has one such irrational term;
 thus an entry error is at most $\sqrt2\,2^{-b-4}$ and the
 Frobenius error is at most $2\sqrt2\,2^{-b-4}\leq2^{-b-2}$.
 The four target approximations, computed once at $b=m+10$, each
 have complex-modulus error at most $2^{-b}$, giving $e_U\leq2^{1-b}$.
 \end{proof}
 \step{eqv-rot-candidates-suffices}{\suffices{Find a rational upper
 bound $u$ for $\|\widetilde W-\widetilde U\|_F$ exceeding it
 by at most $2^{-b}$, and return $u+e_W+e_U$ within the advertised arithmetic cost.}}
 \begin{proof}\pf\ The triangle inequality makes this an upper
 bound for $\|W-U\|_F$. Its excess is at most
 $2(e_W+e_U)+2^{-b}\leq(11/2)2^{-b}<\delta/8$, since
 $b=m+10$ and $\delta=2^{-m}$. The stipulated cost gives the
 remaining resource claim.\end{proof}
 \step{eqv-rot-candidates-root}{Such a rational $u$ is obtained by
 bisection without a positive lower bound on the norm.}
 \begin{proof}\pf\ Sum exactly the squares of the eight real
 components of $\widetilde W-\widetilde U$. The sum is rational.
 Both exact matrices are unitary, so their Frobenius distance is at
 most $2\sqrt2$; the stated approximation errors put the approximated
 distance in $[0,4]$. Bisect this interval using exact comparisons
 of squared midpoints with the sum, until its width is at most
 $2^{-b}$. Its upper endpoint is the required $u$, also when the
 sum is zero.\end{proof}
 \step{eqv-rot-candidates-cost}{The per-candidate arithmetic costs
 $O((L_{\max}+b+1)^6)$ multitape steps.}
 \begin{proof}\pf\ Exact multiplication of the $O(L_{\max})$ fixed
 letters, the two bisections and the fixed-size matrix arithmetic
 use polynomially many integer operations. There are $O(L_{\max}+b+1)$
 bisections and $O(L_{\max}+b+1)$-bit intermediate integers: rational
 approximations may use dyadic denominators, and only a fixed
 number of their squares and products are combined. Full scans,
 grade-school multiplication and division are bounded by the stated
 conservative sixth power. The target evaluation is charged once
 as $\mu_U(b)$ in the enclosing proof.\end{proof}
 \qedstep
 \begin{proof}\pf\ By \stepref{eqv-rot-candidates-matrices},
 \stepref{eqv-rot-candidates-suffices},
 \stepref{eqv-rot-candidates-root} and
 \stepref{eqv-rot-candidates-cost}.\end{proof}
\end{proof}
  \step{eqv-rot-accept}{Some word of at most $L_{\max}$ letters is accepted,
  and every accepted word meets the original request.}
  \begin{proof}\pf\ Accept precisely when the rational upper bound
  is less than $\delta$. Since operator norm is at most Frobenius
  norm, acceptance certifies error below $\delta\leq\eta$.
  Step~\stepref{eqv-rot-sk} supplies a word with operator error
  below $\delta/8$, hence Frobenius error below
  $\sqrt2\delta/8$. Its computed upper bound is at most
  $(\sqrt2+1)\delta/8<\delta$, so it is accepted. Comparisons
  are rational with a strict margin; no exact real equality is tested.
  \end{proof}
  \qedstep
  \begin{proof}\pf\ There are fewer than $4^{L_{\max}+1}$ candidates.
  Add \stepref{eqv-rot-request}, the single cost $\mu_U(b)$,
  and the per-candidate cost in \stepref{eqv-rot-candidates};
  \stepref{eqv-rot-accept} guarantees termination. For a uniform
  implementation one may simply enumerate increasing lengths until
  acceptance; $L_{\max}$ is an upper bound, not a required stopping test.
  \end{proof}
 \end{proof}
 \step{eqv-rot-control}{Controlled approximations preserve the
 phase-sensitive estimate.}
 \begin{proof}\pf\ For a word $W$ approximating $U$, replace each
 elementary gate by its exact controlled implementation from
 Lemma~\ref{lem:synthesis}. Their product is
 $\proj0\otimes I+\proj1\otimes W$, with no extra phase.
 Its difference from the controlled $U$ is
 $\proj1\otimes(W-U)$, of norm $\|W-U\|$.
 Workspace is constant and clean. For a fixed number of controls,
 compute their conjunction by a finite Toffoli circuit, control
 the whole word, then undo the conjunction; the control inputs
 are preserved, so the workspace returns to zero.\end{proof}
 \step{eqv-rot-angles}{The optional gates satisfy the entry hypotheses.}
 \begin{proof}\pf\ $R_z(\theta)$ has determinant one exactly, and its inverse is $R_z(-\theta)$, whose entries are the complex conjugates of the supplied approximations. Reversing the same approximating circuit gives the required correlated adjoint choice in Theorem~\ref{eqv:approx}. No entry procedure is derived here from a description of $\theta$ alone; Definition~\ref{def:circuit} requires that procedure explicitly, with its cost charged to $\mu_U$.\end{proof}
 \qedstep
 \begin{proof}\pf\ By \stepref{eqv-rot-sk},
 \stepref{eqv-rot-search}, \stepref{eqv-rot-control}, and
 \stepref{eqv-rot-angles}. Approximate universality does not replace
 the exact obstruction in Counterexample~\ref{eqv:exact-obstruction}.
\end{proof}
\end{proof}

\begin{remark}[What the equivalence does and does not say]\label{eqv:status}
Theorems~\ref{eqv:forward} and \ref{eqv:reverse} establish the bounded equivalence of Definition~\ref{eqv:def} with its stipulated finite capacities,
separately scheduled timeout, actual gate alphabet, and output-instrument meaning of
equivalence. They make explicit that state preparation and extraction
are included using $T\geq1$, that literal-table compilation need not
enumerate all configurations, and that drop does not reclaim storage.
The converse supplies the explicit layout bound
$s\leq(b+2)w+q+8$ and comparison bound $O(N(s+w+1))$.
The global canonical-decoder and live-flag conditions of
Definition~\ref{def:device},
all-code exclusion of read mode, valid initial configuration, initially
exact input live set, final-time liveness of $O$ on every reachable
outcome sequence, blank distinct forward records, and trace of the
entire complement of $OY$ are hypotheses throughout. They are not
claims about arbitrary controllers with the same instruction names.
Theorem~\ref{eqv:approx} proves the approximation conditional on Definition~\ref{def:certificates}, specifying its additional classical and
workspace costs. Theorem~\ref{eqv:rotations} discharges the certificates
for the optional one-qubit determinant-one rotations and bounded
controls, charging the full binary rational request length and dyadic
entry-computation cost in its explicit time bound. It does not remove
that hypothesis for every additional
matrix of arbitrary arity or scalar phase. The inverse-closed set
actually used for Solovay--Kitaev is $\mathcal F\subset SU(2)$;
replacing $T^\dagger$ by $T^7$ afterward is exact. None of these
statements imports a general dynamic mark or stopping semantics. Observational equivalence preserves final instruments only, not the tick at which a byte is written or a run stops. The certificate words for ideal rotations are found by exhaustive certified search with the exponential-in-$L_{\max}$ time bound of Theorem~\ref{eqv:rotations}; polynomial classical synthesis time in $\log(1/\eta)$ is not proved here, and certificates for matrices beyond determinant-one one-qubit rotations with a fixed number of controls remain a hypothesis of Theorem~\ref{eqv:approx}.
\end{remark}

\notedsection{Soundness for structured programs (Tier~A)}{sec:soundness}
  {This section studies \earlierforms.}

\subsection{Recursive soundness with retained registers}

\begin{theorem}[Structured soundness]\label{thm:sA-sound}
Fix the finite retained list $L=QM$ of Definition~\ref{def:tiera}.
Let $P$ be a finite term of its fixed-space grammar, with the specified
\emph{total} atom completions and total keep functions. Impose all the
restrictions in that definition, including static mark pairing and no
allocation, drop, observed measurement, external console observation, or
actual stop inside a marked body. Then $\mathcal U_P$ is a uniquely
specified unitary on all of $\cH_L$, at every finite nesting depth
(the algebraic part is Lemma~\ref{lem:term-types}; new here are the clean
implementations and the menu extension).
In particular, tick decoders obey Definition~\ref{def:device} and
exclude measurement and allocation modes on every configuration code;
continuation addresses fit their field. A declared timeout horizon
includes its explicitly scheduled final snapshot and its records.
For every marked subterm, with $U=\mathcal U_R$,
\[
 \mathcal U_{\mathsf M(R,K)}=U^\dagger KU,
 \qquad (U^\dagger KU)^\dagger=U^\dagger K^\dagger U.
\]
In particular, gates changing previously read qubits are permitted.

If each atom has a specified exact clean implementation over
$\mathcal G^{\pm}$, the recursively expanded circuit implements
$\psi\mapsto\mathcal U_P\psi\otimes\ket0_Z$ for every logical input
$\psi$, with implementation workspace $Z$ initialized to zero. All
logical records remain in $L$. No action on nonzero $Z$ is asserted.
Well-typed outer programs of Definition~\ref{def:tiera} denote isometries
on the entire retained output, and tracing designated output complements
gives channels, also on inputs entangled with an external reference.

The same conclusions hold for the extension by occurrence-typed derived
menu atoms of Definition~\ref{def:tierb}, including their use in
sequences, marked bodies and outer runs. Selectors are disjoint from
payloads and retained for enclosing inverses, entries are base Tier~A
terms, and all capacity, occurrence and phase hypotheses there apply.
The implementation conclusion requires the supplied exact clean
$\mathcal G^{\pm}$ lists for menu entries. No exact finite-alphabet
implementation of an arbitrary ideal-angle atom is asserted. Joint
ticks, Boolean atoms, the listed controlled gates and keeps have such
implementations by Lemmas~\ref{lem:joint}, \ref{lem:synthesis}
and \ref{lem:keep}; other admitted atoms retain the explicit
implementation hypothesis. These assertions
supply neither ordinary forward semantics for nonblank logical records
nor cleanliness of $M$, and cover no unrestricted dynamic source marks.
\end{theorem}
\begin{proof}
 \step{sA-sound-ind}{Each derived menu atom is a specified total
 unitary, with a clean implementation under the stated list hypothesis.}
 \begin{proof}\pf\ Its entries are base Tier~A terms and hence total
 unitaries by Lemma~\ref{lem:term-types}. Lemma~\ref{lem:menu}
 applies to their direct sum and, when supplied, their exact circuit
 lists. This establishes the additional atom case without assuming
 anything about a term containing menus.
 \end{proof}
 \step{sA-sound-struct}{Every fixed-space term, also in the menu
 extension, is a uniquely specified total unitary; if its atoms have
 the specified implementations, it and its adjoint have finite clean
 implementations.}
 \begin{proof}
  \step{sA-sound-base}{\case{$P=I$ or $P=a$.}}
  \begin{proof}
   \step{sA-sound-base-total}{The denotation is a specified total unitary.}
\begin{proof}
 \step{sA-sound-base-total-gates}{\case{Identity or an elementary gate.}}
 \begin{proof}\pf\ Their matrices are unitary, including the
 explicitly declared ideal matrices. Tensor identity on the
 complementary retained wires.\end{proof}
 \step{sA-sound-base-total-machine}{\case{A completed RTM step or
 ordinary joint tick.}}
 \begin{proof}\pf\ Lemmas~\ref{lem:completion} and \ref{lem:joint}
 establish totality on the specified supports, with identity on the
 complement. The tick restrictions on every code are supplied data,
 rather than promises about particular inputs.\end{proof}
 \step{sA-sound-base-total-cont}{\case{A continuation.}}
 \begin{proof}\pf\ The map $(c,j)\mapsto(j\oplus h_\ell(c),c)$
 has inverse $(c',j')\mapsto(j',c'\oplus h_\ell(j'))$.
 Substitution in either order cancels two copies of $h_\ell$.
 No position-restoration claim about a preceding mark is used.
 \end{proof}
 \step{sA-sound-base-total-timeout}{\case{A scheduled timeout snapshot.}}
 \begin{proof}\pf\ The same inverse calculation with $h_\tau$
 in place of $h_\ell$, leaving the sample register unchanged,
 makes this a total permutation and hence a unitary.\end{proof}
 \step{sA-sound-base-total-boolean}{\case{A Boolean atom.}}
 \begin{proof}\pf\ Lemma~\ref{lem:synthesis} supplies its
 implementation, with the specified total logical action.\end{proof}
 \step{sA-sound-base-total-menu}{\case{A derived menu.}}
 \begin{proof}\pf\ Apply \stepref{sA-sound-ind}, which establishes
 its specified total unitary before the structural induction.\end{proof}
 \qedstep
 \begin{proof}\pf\ These cases give the allowed atoms. A finite
 composition of unitary atoms has inverse the product of their
 inverses in reverse order, by adjacent cancellation. Thus any
 supplied finite composition also has its specified total action.
 \end{proof}
\end{proof}
   \qedstep
   \begin{proof}\pf\ The empty list implements identity. For an atom
   use its assumed clean circuit. Lemma~\ref{lem:scratch} gives its
   clean adjoint circuit by reversal. The gate alphabet is closed under
   adjoints, since $H$ and CNOT are self-adjoint and both $T,T^\dagger$
   are available.
   \end{proof}
  \end{proof}
  \step{sA-sound-seq}{\case{$P=R;S$, with both assertions true for
  its proper subterms.}}
  \begin{proof}
   \step{sA-sound-seq-unit}{Writing $U=\mathcal U_R$ and
   $V=\mathcal U_S$, the denotation $VU$ is unitary.}
   \begin{proof}\pf\ $(VU)^\dagger(VU)=U^\dagger V^\dagger VU=\id$
   and $(VU)(VU)^\dagger=VUU^\dagger V^\dagger=\id$.
   \end{proof}
   \qedstep
   \begin{proof}\pf\ Concatenate the circuit for $R$ with that for
   $S$. Each returns workspace to zero on every logical input, so
   Lemma~\ref{lem:scratch} permits padded shared workspace. Reverse
   the resulting circuit for the adjoint. Finite concatenation preserves
   finiteness and the actual matrix phases.
   \end{proof}
  \end{proof}
  \step{sA-sound-mark}{\case{$P=\mathsf M(R,K)$, with both assertions
  true for $R$.}}
  \begin{proof}
   \step{sA-sound-mark-unit}{For $U=\mathcal U_R$, $W=U^\dagger KU$
   is unitary and $W^\dagger=U^\dagger K^\dagger U$.}
   \begin{proof}\pf\ Lemma~\ref{lem:keep} makes $K$ unitary.
   Reversal of factors under the adjoint gives the claimed identity;
   multiplication gives $W^\dagger W=U^\dagger K^\dagger KU=\id$
   and $WW^\dagger=U^\dagger KK^\dagger U=\id$.
   \end{proof}
   \step{sA-sound-mark-circuit}{The ordered list for $R$, then $K$,
   then the reversed adjoint list for $R$, implements $W$ cleanly.}
   \begin{proof}\pf\ The induction hypothesis supplies both lists
   for $R$. Lemma~\ref{lem:keep} supplies the keep circuit. Each is
   clean on the full logical space, so its successor receives zero
   workspace regardless of the logical state produced by $K$.
   Lemma~\ref{lem:scratch} permits their composition and reversal.
   Reversing the full list includes the keeps of completed inner terms.
   \end{proof}
   \qedstep
   \begin{proof}\pf\ By \stepref{sA-sound-mark-unit} and
   \stepref{sA-sound-mark-circuit}. No part of this calculation requires
   preservation of a previously read register or blank logical records.
   \end{proof}
  \end{proof}
  \qedstep
  \begin{proof}\pf\ These are all constructors of the finite grammar.
  Structural induction therefore proves the assertions and uniqueness:
  every constructor applies a specified operation to its already specified
  proper subterms.
  \end{proof}
 \end{proof}
 \step{sA-sound-outer}{Every well-typed outer program is an isometry.}
 \begin{proof}\pf\ For an allocation $A_B$, $A_B^\dagger A_B=\id$
 because the appended zero vector has norm one. Identity and run are
 isometries by \stepref{sA-sound-struct}. For composable isometries
 $V_1,V_2$, $(V_2V_1)^\dagger V_2V_1=V_1^\dagger V_1=\id$.
 Here the fixed-space induction holds for each finite list separately,
 so each run uses its own current list. A term on a smaller list may
 be lifted by the tensor-identity rule of Definition~\ref{def:tiera}.
 Induction on the typed outer grammar proves the assertion, also with
 derived menus admitted as atoms. A physical run using allocated
 workspace agrees with the lifted logical run on zero workspace by the padding rule of Definition~\ref{def:tiera}, with Definition~\ref{def:scratch} supplying the zero-workspace equality for each implementation; no equality off that subspace is used.
 \end{proof}
 \step{sA-sound-trace}{Tracing a designated complement of an isometry
 gives a channel, including on inputs with a reference.}
 \begin{proof}\pf\ For an isometry $V:\cH_{\rm in}\to\cH_O\otimes
 \cH_F$, choose an orthonormal basis $(f_j)$ of $F$ and set
 $B_j=(\id_O\otimes\bra{f_j})V$. Partial trace is
 $\Phi(\rho)=\sum_jB_j\rho B_j^\dagger$ and
 $\sum_jB_j^\dagger B_j=V^\dagger V=\id$. Each summand remains
 positive after tensoring with any identity, and the last equality
 preserves trace. This proves complete positivity and trace preservation.
 \end{proof}
 \step{sA-sound-menu}{The finite-menu extension is sound.}
 \begin{proof}\pf\ The derived atom case was established in
 \stepref{sA-sound-ind} and included in \stepref{sA-sound-struct}.
 Its adjoint is the same selector sum with each payload replaced by
 its adjoint. Reversing its controlled circuit preserves the selector
 and reuses the same occurrence-labelled records, as required by
 Definition~\ref{def:tierb}; Lemma~\ref{lem:scratch} proves clean
 reversal. Thus the sequence and marked cases already apply to it.
 Outer composition and tracing apply by
 \stepref{sA-sound-outer} and \stepref{sA-sound-trace}.
 \end{proof}
 \qedstep
 \begin{proof}\pf\ By \stepref{sA-sound-struct}--\stepref{sA-sound-menu}.
 All reasoning is on total operators; no interpretation of a nonblank
 record, or of an operation outside the grammar, is inferred.
 \end{proof}
\end{proof}


\notedsection{Cleanliness and disposal of auxiliary registers}{sec:cleanliness}
  {This section studies \earlierforms.}
Soundness (Section~\ref{sec:soundness}) says a marked term is a unitary;
it does not say the machine register can be discarded. This section and
Definition~\ref{def:scratch} separate five assertions. A \emph{clean
implementation} restores only its workspace $Z$, without a reset promise
for logical records. A \emph{clean unitary} returns $M$ to its exact
initial vector, $WE=EA$, so $M$ may be dropped and the induced operation
on $Q$ is the unitary $A$. \emph{Reset of a factor $F$ on a subspace
$\mathcal S$} means only that $F$ is $\ket0$ after $W$ for inputs in
$\mathcal S$---the right notion inside a larger computation, where
$\mathcal S$ is whatever the preceding steps actually reach.
\emph{Marginal restoration} means only that $M$'s reduced state is
unchanged, which permits neither deletion nor replacement by fresh zeros
when $M$ is correlated with data. Finally the reduced operation on $Q$
can be unitary while $M$ ends in a fixed vector $\eta\neq m_0$ ((SA2)):
discarding $M$ is then harmless for $Q$ alone, but $M$ has not been reset
and cannot be reused as blank. Theorems~\ref{thm:sA-clean} and
\ref{thm:sA-recursive} give exact tests for logical restoration, subspace
reset and marginal restoration; Counterexample~\ref{ex:sA-rec-limits}
shows the limits of their implications. The collective disposal test for
the dynamic interpreter in Theorem~\ref{thm:tB-clean} has the same
fixed-pure-factor scope as (SA2).

\subsection{Exact tests for cleanliness and for a unitary reduced channel}

\begin{theorem}[Exact cleanliness]\label{thm:sA-clean}
Let $Q,M$ be finite registers, $\ket{m_0}_M$ a fixed unit vector,
$E\psi=\psi\otimes\ket{m_0}$ and $P_0=EE^\dagger$ as in
Definition~\ref{def:clean}. Let $U,K$ be total unitaries on $QM$, and
put $W=U^\dagger KU$, $\Pi=UP_0U^\dagger$. The following are equivalent:
\begin{enumerate}
 \item $WE=EA$ for a unitary $A$ on $Q$;
 \item $(\id-P_0)WE=0$;
 \item $(\id-\Pi)K\Pi=0$;
 \item $[K,\Pi]=0$, where $[K,\Pi]=K\Pi-\Pi K$.
\end{enumerate}
When these hold, $A=E^\dagger WE$ is unique. Equivalently, for every
input density operator $\rho$ on $Q$, the final state
$WE\rho E^\dagger W^\dagger$ has $M$ marginal $\proj{m_0}$.
For a single input density operator $\rho$, this return condition is exactly
\[
 \tr\bigl[(\id-P_0)WE\rho E^\dagger W^\dagger\bigr]=0.       \tag{SA1}
\]
For every finite untouched reference $R$, the reference forms use
$E\otimes\id_R$, $W\otimes\id_R$, $A\otimes\id_R$ and
$P_0\otimes\id_R$, with factors reordered by their names. In (SA1)
the input may then be any density operator on $QR$, including an
entangled one. The universal equivalences are unchanged.

The channel $\Phi(\rho)=\tr_M(WE\rho E^\dagger W^\dagger)$ on $Q$
is unitary if and only if there are a unitary $A$ on $Q$ and a fixed
unit vector $\eta\in\cH_M$ such that
\[
 WE\psi=A\psi\otimes\eta\quad\hbox{for every }\psi\in\cH_Q.
                                                               \tag{SA2}
\]
Here $\eta$ need not be $m_0$: a unitary reduced channel need not be clean.
If (SA2) holds, then tensoring it with $\id_R$ gives
$(W\otimes\id_R)(E\otimes\id_R)\Psi=(A\otimes\id_R)\Psi\otimes\eta$ for
every $\Psi\in\cH_{QR}$, and the reduced channel on $QR$ is conjugation
by $A\otimes\id_R$. Conversely, items 1–4 and the (SA2) equivalence for
the tensored data follow by applying this theorem with $QR$ in place of
$Q$ and $U\otimes\id_R$, $K\otimes\id_R$ in place of $U$, $K$; the
uniqueness in item 1 then forces the induced unitary on $QR$ to be
$A\otimes\id_R$.
\end{theorem}
\begin{proof}
 \step{sA-clean-restrict}{Conditions 1 and 2 are equivalent, with the
 stated unique $A$.}
 \begin{proof}
  \step{sA-clean-12}{\assume{Condition 1.}\prove{Condition 2.}}
  \begin{proof}\pf\ $(\id-P_0)E=0$, so multiply $WE=EA$ by
  $\id-P_0$ on the left.
  \end{proof}
  \step{sA-clean-21}{\assume{Condition 2.}\prove{Condition 1.}}
  \begin{proof}
   \step{sA-clean-A}{Set $A=E^\dagger WE$. Then $EA=WE$ and
   $A^\dagger A=\id_Q$.}
   \begin{proof}\pf\ Condition 2 says $P_0WE=WE$, whence
   $EA=P_0WE=WE$ and
   $A^\dagger A=E^\dagger W^\dagger P_0WE=E^\dagger E=\id_Q$.
   \end{proof}
   \qedstep
   \begin{proof}\pf\ A square isometry on a finite-dimensional space
   is surjective: its kernel is zero, so its rank equals the dimension.
   Thus also $AA^\dagger=\id_Q$. Multiplication by $E^\dagger$ shows
   uniqueness.
   \end{proof}
  \end{proof}
  \qedstep
  \begin{proof}\pf\ By \stepref{sA-clean-12} and
  \stepref{sA-clean-21}.
  \end{proof}
 \end{proof}
 \step{sA-clean-conjugate}{Conditions 2 and 3 are equivalent.}
 \begin{proof}\pf\ The range of $E$ is the range of $P_0$, so
 condition 2 is equivalent to $(\id-P_0)WP_0=0$. Conjugate this
 equality by $U$ and substitute $W=U^\dagger KU$ to obtain
 $(\id-\Pi)K\Pi=0$. Conjugation is invertible.
 \end{proof}
 \step{sA-clean-commute}{Conditions 3 and 4 are equivalent.}
 \begin{proof}
  \step{sA-clean-invariant}{\assume{Condition 3.}
  \prove{$K\operatorname{ran}\Pi=\operatorname{ran}\Pi$.}}
  \begin{proof}\pf\ Condition 3 gives inclusion. The restriction of
  the unitary $K$ is injective, so its image has the same finite
  dimension as $\operatorname{ran}\Pi$, proving equality.
  \end{proof}
  \step{sA-clean-reduce}{Under condition 3 the orthogonal complement
  of $\operatorname{ran}\Pi$ is invariant too.}
  \begin{proof}\pf\ For $x\perp\operatorname{ran}\Pi$ and
  $y\in\operatorname{ran}\Pi$, write $y=Kz$ with
  $z\in\operatorname{ran}\Pi$ using \stepref{sA-clean-invariant}.
  Then $\langle y,Kx\rangle=\langle z,x\rangle=0$.
  Thus $K$ preserves both summands of the orthogonal decomposition,
  and commutes with their projection $\Pi$.
  \end{proof}
  \qedstep
  \begin{proof}\pf\ This proves 3 implies 4. Conversely
  $K\Pi=\Pi K$ implies $(\id-\Pi)K\Pi=0$ by multiplication.
  \end{proof}
 \end{proof}
 \step{sA-clean-one}{Equation (SA1) is equivalent to return of $M$
 for that input.}
 \begin{proof}
  \step{sA-clean-positive}{For a positive operator $\sigma$, and an
  orthogonal projection $P$, $\tr[(\id-P)\sigma]=0$ if and only if
  $\operatorname{supp}\sigma\subseteq\operatorname{ran}P$.}
  \begin{proof}\pf\ A spectral decomposition
  $\sigma=\sum_j\lambda_j\proj{v_j}$ with $\lambda_j>0$ on its
  support gives the trace as
  $\sum_j\lambda_j\|(\id-P)v_j\|^2$. Its nonnegative summands
  vanish precisely when every $v_j$ lies in that range.
  \end{proof}
  \qedstep
  \begin{proof}\pf\ Apply \stepref{sA-clean-positive} to
  $\sigma=WE\rho E^\dagger W^\dagger$ and $P=P_0$.
  Support in $\cH_Q\otimes\ket{m_0}$ is equivalent to marginal
  $\proj{m_0}$: one direction follows by partial trace; the other
  follows from
  $\id-P_0=\id_Q\otimes(\id_M-\proj{m_0})$ and
  $\tr[(\id-P_0)\sigma]
  =\tr[(\id_M-\proj{m_0})\tr_Q\sigma]$.
  Testing all pure inputs yields condition 2. Tensoring these
  identities with $\id_R$ gives the stated reference form; the
  implication back to the unreferenced test uses a one-dimensional $R$. Conversely, items 1--4 and the (SA2) equivalence for the tensored data follow by applying this theorem with $QR$ in place of $Q$ and $U\otimes\id_R$, $K\otimes\id_R$ in place of $U$, $K$; the uniqueness in item 1 then forces the induced unitary on $QR$ to be $A\otimes\id_R$.
  \end{proof}
 \end{proof}
 \step{sA-clean-channel}{The reduced channel is unitary exactly under (SA2).}
 \begin{proof}
  \step{sA-clean-channel-if}{\assume{Equation (SA2).}
  \prove{$\Phi(\rho)=A\rho A^\dagger$.}}
  \begin{proof}\pf\ Expand $\rho$ in matrix units and trace the
  fixed factor $\proj\eta$, whose trace is one.
  \end{proof}
  \step{sA-clean-channel-only}{\assume{$\Phi(\rho)=A\rho A^\dagger$
  for a unitary $A$ and every $\rho$.}\prove{Equation (SA2).}}
  \begin{proof}
   \step{sA-clean-basis}{\pick\ an orthonormal basis
   $e_1,\ldots,e_d$ of $Q$. With $V=(A^\dagger\otimes\id_M)WE$,
   there are unit vectors $\eta_j$ such that $Ve_j=e_j\otimes\eta_j$.}
   \begin{proof}\pf\ A finite Hilbert space has an orthonormal basis.
   The $Q$ marginal of $\proj{Ve_j}$ is $\proj{e_j}$; the support
   argument of \stepref{sA-clean-one}, with the factors exchanged,
   forces $Ve_j$ into $\ket{e_j}\otimes\cH_M$. Norm preservation
   makes $\eta_j$ unit.
   \end{proof}
   \step{sA-clean-same}{Every $\eta_j$ is the same vector.}
   \begin{proof}\pf\ If $i\ne j$, apply the unchanged $Q$ marginal
   property of $V$ to $(e_i+e_j)/\sqrt2$. Its $(i,j)$ matrix element
   is $\langle\eta_j,\eta_i\rangle/2$, and must be $1/2$.
   Hence $\|\eta_i-\eta_j\|^2=2-2\operatorname{Re}
   \langle\eta_j,\eta_i\rangle=0$. For $d=1$ there is just one
   vector and the assertion already holds.
   \end{proof}
   \qedstep
   \begin{proof}\pf\ Set $\eta=\eta_1$ and extend the basis
   equality by linearity; multiply $V$ by $A\otimes\id_M$.
   \end{proof}
  \end{proof}
  \qedstep
  \begin{proof}\pf\ By \stepref{sA-clean-channel-if} and
  \stepref{sA-clean-channel-only}.
  \end{proof}
 \end{proof}
 \qedstep
 \begin{proof}\pf\ The four equivalences follow from
 \stepref{sA-clean-restrict}--\stepref{sA-clean-commute}; the
 input and channel assertions follow from \stepref{sA-clean-one}
 and \stepref{sA-clean-channel}.
 \end{proof}
\end{proof}

\begin{definition}[Logical occurrences and retained segments]
\label{def:sA-reset-points}
For a finite fixed-space term define two annotated ordered expansions
$\mathcal E_+(P)$ and $\mathcal E_-(P)$. Lists here are in execution
order; the minus sign denotes the adjoint. Identity gives an empty list.
For an atom $a$ the lists are $(A_a)$ and $(A_a^\dagger)$.
For a sequence, concatenate in the orders
\[
 \mathcal E_+(R;S)=(\mathcal E_+(R),\mathcal E_+(S)),\qquad
 \mathcal E_-(R;S)=(\mathcal E_-(S),\mathcal E_-(R)).
\]
For each expanded occurrence of a marked term set
\[
 \mathcal E_\epsilon(\mathsf M(R,K))
   =[\mathcal E_+(R),K_\epsilon,\mathcal E_-(R)]_{o,\epsilon},
 \qquad K_+=K,\quad K_-=K^\dagger,
 \quad\epsilon\in\{+,-\}.
\]
The brackets are annotations, not operations: $o$ distinguishes the
individual occurrence in the finite expansion, even when several
occurrences come from the same syntactic subterm. Its entry is immediately
before its bracketed list, its exit immediately after. A \emph{marked
reset point} is that exit together with its annotated occurrence and
a specified retained factor $F_o$ to test for zero. The atom and record
identifiers stay those of the original syntax in both expansions.
Outer allocations are left as zero embeddings, outer sequences are
concatenated, and each run is replaced by its plus expansion.
Atoms, including derived menu atoms, are indivisible at this logical
level: no marked occurrences inside their supplied implementations or
inside the menu entries are enumerated here.

More generally, in a specified retained circuit of total unitary gates
and fresh zero embeddings, a \emph{segment} is a consecutive list of
operations with entry and exit boundaries, possibly the empty list.
Its product is an isometry $V:\cH_A\to\cH_B$ on its typed input
and output lists. This definition also applies to a supplied physical
implementation, using its full register lists including workspace.
Thus a boundary inside an implemented atom or keep has a meaning only
after that implementation has been specified. A reset test certifies
the state already produced there; it inserts no operation and deletes
no register.
\end{definition}

\begin{theorem}[Recursive reset and marginal-restoration tests]
\label{thm:sA-recursive}
Let $L$ be a finite retained register list, $U,K$ total unitaries on it,
$W=U^\dagger KU$, and $\mathcal S\subseteq\cH_L$ a specified subspace
with orthogonal projection $P_{\mathcal S}$. Let $F\subseteq L$ be the
auxiliary factor proposed for reset, and set
$P_F=\id_{L\setminus F}\otimes\proj0_F$ (the projector $P_A$ of
Definition~\ref{def:clean}, with $F$ in place of $A$). For arbitrary sublists,
these tensor expressions use the canonical permutation of factors from
their order in $L$ into the displayed order, and are transported back
by its inverse. The same convention applies to $M'$ below.
Then $F$ is zero after $W$ on every input
supported on $\mathcal S$ if and only if
\[
 (\id-UP_FU^\dagger)KU P_{\mathcal S}=0.                  \tag{SA3}
\]
For reuse as an initially blank record the additional incoming condition
$\mathcal S\subseteq\operatorname{ran}P_F$ is required. If
$\mathcal S=\operatorname{ran}P_F$, (SA3) is equivalent to
$[K,UP_FU^\dagger]=0$.

Let $M'\subseteq L$ instead be a factor whose marginal is to be restored,
with no product-state assumption. Fix a basis of $M'$ and its matrix
units $B_{ij}=\ket i\bra j$. Preservation of the joint reduced state
on $M'$ and every finite untouched external reference $R$, for every
joint density operator supported on $\mathcal S\otimes\cH_R$ under
$W\otimes\id_R$, is equivalent to
\[
 P_{\mathcal S}W^\dagger(\id_{L\setminus M'}\otimes B_{ij})
 WP_{\mathcal S}
 =P_{\mathcal S}(\id_{L\setminus M'}\otimes B_{ij})
 P_{\mathcal S}\quad\hbox{for all }i,j.                 \tag{SA4}
\]

For any specified segment isometry $V:\cH_A\to\cH_B$, input
subspace $\mathcal S\subseteq\cH_A$, and factor $F\subseteq B$,
set $P_F^B=\id_{B\setminus F}\otimes\proj0_F$. Reset at its exit
on every input supported on $\mathcal S$ is equivalent to
\[
 (\id_B-P_F^B)VP_{\mathcal S}=0.                       \tag{SA3a}
\]
This includes physical segment boundaries inside atoms, with all
workspace included in $A,B$ and no zero-workspace assumption there.

For a finite Tier~A term or typed outer program, also with the derived
menus of Definition~\ref{def:tierb}, fix an initial subspace
$\mathcal S_0$ and a finite collection of marked reset points as
defined in Definition~\ref{def:sA-reset-points}. Marked occurrences inside the entries of a derived menu atom are not
marked reset points at this logical level; a boundary there is tested by
(SA3a) on the corresponding segment of the specified controlled circuit,
with the selector included in its register lists. For occurrence $(o,\epsilon)$ of $\mathsf M(R,K)$, let $L_o$ be the
register list at its entry in the expansion of the whole program, let
$U$ and $K_\epsilon$ denote the denotations of $R$ and of the central
operation on $L_o$ (tensor-identity lifts of
Definition~\ref{def:tiera} included), let $F_o\subseteq L_o$, and let
$J_o$ be the ordered product of all expansion operations before its
entry, so that the exact entry subspace is
$\mathcal S_o=J_o\mathcal S_0$. Equation (SA3), with $K_\epsilon$ in
place of $K$, is jointly necessary and sufficient for zero at all
those exits for every allowed initial input. For the keeps of
Definition~\ref{def:keep}, $K^\dagger=K$, but the plus and minus
occurrences can have different entry subspaces. For any specified retained circuit with a specified initial subspace
$\mathcal S_0^{\mathrm{circ}}$ of its own input list (for a clean
implementation of a term, $\mathcal S_0\otimes\ket0_Z$ with workspace
named), the same conjunction assertion holds for a finite collection of
segment exits tested by (SA3a), with each segment's input subspace the
exact image of $\mathcal S_0^{\mathrm{circ}}$ under the prefix ending at
that segment's entry. These assertions include inputs supported on
$\mathcal S_0\otimes\cH_R$ and untouched references.
Replacing exact entry subspaces by containing subspaces preserves
sufficiency; necessity for actual reachable inputs is then not asserted.
No equivalence between final cleanliness and cleanliness of every
syntactic subterm is claimed. Blank-record use requires
zero at its actual entry, supplied by fresh zero allocation, an earlier
proved reset and preservation until use, or another proof of that entry
condition. Neither (SA4) nor a reset alone permits deletion of a factor
still required by an enclosing inverse.
\end{theorem}
\begin{proof}
 \step{sA-rec-reset}{Reset on $\mathcal S$ is equivalent to (SA3).}
 \begin{proof}
  \step{sA-rec-support}{Reset on every density operator supported on
  $\mathcal S$ is equivalent to $(\id-P_F)WP_{\mathcal S}=0$.}
  \begin{proof}\pf\ Test rank-one inputs $\proj\psi$ with
  $\psi\in\mathcal S$. Their output has zero $F$ marginal exactly
  when $W\psi\in\operatorname{ran}P_F$, by the positive-support
  calculation in Theorem~\ref{thm:sA-clean}. This gives the operator
  equality; its converse for mixed inputs follows by a spectral
  decomposition supported on $\mathcal S$.
  \end{proof}
  \qedstep
  \begin{proof}\pf\ Multiply the equality of
  \stepref{sA-rec-support} by $U$ on the left, use
  $U(\id-P_F)=(\id-UP_FU^\dagger)U$ and $UW=KU$.
  Multiplication by $U$ is reversible, proving the equivalence.
  \end{proof}
 \end{proof}
 \step{sA-rec-full}{For $\mathcal S=\operatorname{ran}P_F$ the
 commutator criterion is equivalent.}
 \begin{proof}\pf\ Set $\Pi_F=UP_FU^\dagger$. Equation (SA3)
 becomes $(\id-\Pi_F)K\Pi_F=0$ after right multiplication by
 $U^\dagger$. The finite-dimensional invariant-subspace argument in
 Theorem~\ref{thm:sA-clean} gives $[K,\Pi_F]=0$ and its converse.
 \end{proof}
 \step{sA-rec-marginal}{Marginal restoration is equivalent to (SA4).}
 \begin{proof}
  \step{sA-rec-marginal-need}{\assume{Restoration for all states
  supported on $\mathcal S$.}\prove{Equation (SA4).}}
  \begin{proof}\pf\ Write
  $D=W^\dagger(\id\otimes B_{ij})W-\id\otimes B_{ij}$.
  Equality of the marginals implies $\langle\psi,D\psi\rangle=0$
  for every $\psi\in\mathcal S$, by homogeneity from normalized
  inputs. For $u,v\in\mathcal S$, apply this to $u+v$ and $u+iv$,
  and subtract the diagonal terms: one obtains respectively
  $\langle u,Dv\rangle+\langle v,Du\rangle=0$ and
  $i\langle u,Dv\rangle-i\langle v,Du\rangle=0$.
  Hence both cross terms vanish, without requiring $D$ Hermitian.
  Thus $P_{\mathcal S}DP_{\mathcal S}=0$.
  \end{proof}
  \step{sA-rec-marginal-suff}{\assume{Equation (SA4).}
  \prove{Restoration including every external reference $R$.}}
  \begin{proof}\pf\ For any matrix $C$ on $R$, tensor the compressed
  equality with $C$. Every joint density operator $\sigma$ with support
  in $\mathcal S\otimes\cH_R$ satisfies
  \[
   \tr\!\left[\sigma\bigl(W^\dagger(\id\otimes B_{ij})W
          -\id\otimes B_{ij}\bigr)\otimes C\right]=0.
  \]
  Matrix units $B_{ij}$ and matrix units $C$ span all operators on
  $M'R$. Equality of their expectations is equality of the reduced
  density operators. This proves the claimed restoration.
  \end{proof}
  \qedstep
  \begin{proof}\pf\ Necessity holds by taking the reference
  one-dimensional in \stepref{sA-rec-marginal-need}; sufficiency is
  \stepref{sA-rec-marginal-suff}.
  \end{proof}
 \end{proof}
 \step{sA-rec-segment}{The general segment criterion is (SA3a), also
 with untouched references.}
 \begin{proof}\pf\ For a unit vector $\psi\in\mathcal S$, the
 isometric output $V\psi$ has zero $F$ marginal exactly when
 $(\id_B-P_F^B)V\psi=0$, by the positive-support argument of
 Theorem~\ref{thm:sA-clean}. Testing all such vectors gives (SA3a);
 spectral decomposition proves its converse for density operators.
 Tensoring the zero operator with $\id_R$ gives the same assertion
 on $\mathcal S\otimes\cH_R$. Conversely a one-dimensional reference
 recovers (SA3a). No intermediate workspace condition was used.
 \end{proof}
 \step{sA-rec-prefix}{Exact entry subspaces can be determined without
 assuming any intermediate cleanliness.}
 \begin{proof}
  \step{sA-rec-prefix-rules}{For a fixed-space term define
  $\mathcal T_P(\mathcal S)=\mathcal U_P\mathcal S$. Then
  \[
  \begin{aligned}
   \mathcal T_I(\mathcal S)&=\mathcal S,\\
   \mathcal T_a(\mathcal S)&=A_a\mathcal S,\\
   \mathcal T_{P;R}(\mathcal S)&=\mathcal T_R(\mathcal T_P(\mathcal S)),\\
   \mathcal T_{\mathsf M(P,K)}(\mathcal S)
      &=\mathcal U_P^\dagger K\mathcal U_P\mathcal S.
  \end{aligned}
  \]}
  \begin{proof}\pf\ Substitute the denotations in
  Definition~\ref{def:tiera}, also for the derived menu atom case.
  For an outer allocation use $A_B\mathcal S$; for an adjoint atom
  use $A_a^\dagger\mathcal S$. At an annotated occurrence
  $(o,\epsilon)$ of $\mathsf M(P,K)$, the body starts in
  $\mathcal S_o$, its central operation starts in
  $\mathcal U_P\mathcal S_o$, and the inverse body starts in
  $K_\epsilon\mathcal U_P\mathcal S_o$.
  The expansion rules multiply to $\mathcal U_P^\dagger
  K_\epsilon\mathcal U_P$: for the minus occurrence this is
  $(\mathcal U_P^\dagger K\mathcal U_P)^\dagger$, as proved in
  Theorem~\ref{thm:sA-sound}. The reversed sequence rule likewise
  multiplies to the adjoint. Induction on the term therefore validates
  both signed expansions, including every nested occurrence. Prefixes
  inside an inverse body are evaluated from its actual entry subspace.
  \end{proof}
  \step{sA-rec-prefix-projector}{If a prefix isometry is $V$ and the
  initial subspace is $\mathcal S_0$, its exact entry projector is
  $VP_{\mathcal S_0}V^\dagger$.}
  \begin{proof}\pf\ This operator is self-adjoint and its square is
  itself because $V^\dagger V=\id$. Its range is precisely
  $V\mathcal S_0$. Each logical prefix is a product of total unitary
  atoms, keeps, their adjoints, and zero embeddings; a physical prefix
  is a product of total gates and zero embeddings. Each is an isometry
  by the same multiplication as in Theorem~\ref{thm:sA-sound}.
  Moreover every density operator $\rho$ supported on $V\mathcal S_0$
  is reached from $V^\dagger\rho V$, which has trace one, is supported
  on $\mathcal S_0$, and is taken back to $\rho$ by $V$.
  The same statements hold after tensoring with $\id_R$.
  \end{proof}
  \qedstep
  \begin{proof}\pf\ The syntax and its expansion are finite, so
  \stepref{sA-rec-prefix-rules} and \stepref{sA-rec-prefix-projector}
  determine every requested entry subspace by finitely many operations
  on the specified finite matrices. For a supplied physical circuit,
  take its actual gate prefixes on all registers, including workspace.
  These are exact algebraic prescriptions; no efficient procedure for
  deciding their zero equalities is claimed.
  \end{proof}
 \end{proof}
 \step{sA-rec-all}{The tests have exactly the asserted logical scope.}
\begin{proof}
 \step{sA-rec-all-mark}{\assume{The specified exit is a marked
 occurrence.}\prove{Its (SA3) test with $K_\epsilon$ is necessary
 and sufficient for zero at that exit for every allowed initial state.}}
 \begin{proof}\pf\ Step~\stepref{sA-rec-prefix} identifies the
 actual operator as $U^\dagger K_\epsilon U$ and makes every state
 on the exact entry subspace reachable. Apply
 \stepref{sA-rec-reset} with $K_\epsilon$; Lemma~\ref{lem:keep}
 gives $K_-=K_+$ for a keep. The resulting test is equivalent to
 zero at this exit for every allowed initial state.\end{proof}
 \step{sA-rec-all-segment}{\assume{The specified exit is another
 segment exit.}\prove{Its (SA3a) test on the exact entry subspace
 is necessary and sufficient for zero there for every allowed initial state.}}
 \begin{proof}\pf\ Step~\stepref{sA-rec-segment} gives exactly
 the zero-exit test on its entry subspace. Reachability from
 \stepref{sA-rec-prefix} makes it necessary and sufficient for
 every allowed initial state.\end{proof}
 \step{sA-rec-all-conjunction}{The finite conjunction requires exactly
 that all specified exits be zero for every allowed initial state,
 also in the reference forms.}
 \begin{proof}\pf\ Apply \stepref{sA-rec-all-mark} and
 \stepref{sA-rec-all-segment} separately to each exit. Universal
 quantification over initial states commutes with finite conjunction.
 Reference forms follow from \stepref{sA-rec-segment} and
 \stepref{sA-rec-prefix}.\end{proof}
 \step{sA-rec-all-containing}{\assume{$\mathcal S\subseteq
 \widehat{\mathcal S}$.}\prove{Vanishing on $\widehat{\mathcal S}$
 implies vanishing on $\mathcal S$.}}
 \begin{proof}\pf\ Every vector in $\mathcal S$ is a vector in
 $\widehat{\mathcal S}$, so restriction of the zero operator to
 $\mathcal S$ is zero.\end{proof}
 \step{sA-rec-all-entry}{Forward-entry blankness is an additional
 support condition, and the tests do not specify missing factors
 in later inverses.}
 \begin{proof}\pf\ Definition~\ref{def:clean} requires incoming
 support on the zero-auxiliary subspace. A prior reset establishes
 that support only if intervening operations preserve it. Every
 tested equality is an equality of retained operators, whose
 domains still contain all retained factors. It therefore supplies
 no action on a register list from which a factor has been removed.
 \end{proof}
 \qedstep
 \begin{proof}\pf\ The preceding steps establish the scope and
 entry requirements without assuming intermediate cleanliness.
 \end{proof}
\end{proof}
 \qedstep
 \begin{proof}\pf\ By \stepref{sA-rec-reset}--\stepref{sA-rec-all}.
 \end{proof}
\end{proof}

\begin{counterexample}[Limits of local reset tests]\label{ex:sA-rec-limits}
The following finite Tier~A examples show that (i) replacing a reached
subspace by the whole zero-auxiliary subspace can destroy necessity,
and the commutator criterion need not be necessary on a proper subspace;
(ii) final cleanliness does not imply cleanliness of every inner mark;
(iii) marginal restoration does not permit discarding correlations; and
(iv) a unitary reduced channel need not reset to its initial machine state.
Without an incoming-zero hypothesis, the full-space commutator alone
need not even imply reset on a specified input subspace.
In (i)--(ii) let $Q=r$, $M=a$, and use
\[
 U=\operatorname{CNOT}_{r\to a},\qquad
 K=\operatorname{CNOT}_{a\to r},\qquad
 P_a=\id_r\otimes\proj0_a.
\]
Here and below $K_{0,t}$ denotes the single keep entry with target $t$
and constant-zero Boolean function on the entire machine list.
In (iii) take $Q=q$, $M'=a$, and $\mathcal S$ spanned by
$\Omega=(\ket{00}+\ket{11})_{qa}/\sqrt2$, with $U=\id$ and
$K=K_{0,q}=\id$, so $W=U^\dagger KU=\id$.
In (iv) take $Q=r$, $M=a$ initially zero, body $U$ as above and constant
keep $K=X_r$.
\end{counterexample}
\begin{proof}
 \step{sA-limits-swap}{$U^\dagger KU=\operatorname{SWAP}_{ra}$.}
 \begin{proof}\pf\ On a basis pair $(y,z)$ in order $r,a$, the
 three CNOTs give successively $(y,z\oplus y)$,
 $(z,z\oplus y)$, and $(z,y)$. These are the SWAP basis values;
 linearity gives the operator equality.
 \end{proof}
 \step{sA-limits-subspace}{Item (i): the reached-subspace test need
 not agree with the full zero-subspace test.}
 \begin{proof}\pf\ For $\mathcal S=\operatorname{span}\{\ket{00}\}$,
 SWAP fixes its vector, so reset holds. But
 $\ket{10}\in\operatorname{ran}P_a$ is taken to $\ket{01}$,
 outside that range. Thus reset fails on the containing space
 $\operatorname{ran}P_a$, and Theorem~\ref{thm:sA-recursive}
 gives $[K,UP_aU^\dagger]\ne0$. This also proves that a containing
 subspace test is not necessary for actual reachable inputs.
 For the same SWAP, the subspace spanned by $\ket{01}$ resets to
 $a=0$, but is not initially blank: reset alone does not establish
 the incoming-zero requirement. Separately, take $U=\id$ and
 $K=K_{0,r}=\id$ with
 that same input subspace. Then $[K,UP_aU^\dagger]=0$ while $a$
 remains one, disproving sufficiency of the commutator without the
 incoming-zero hypothesis.
 \end{proof}
 \step{sA-limits-nesting}{Item (ii): final cleanliness does not imply
 inner cleanliness.}
 \begin{proof}\pf\ Let $R=\mathsf M(U,K)$, identifying the displayed
 Boolean gate $U$ with its atom, and let $K_0=K_{0,r}$.
 The nested term $\mathsf M(R,K_0)$ has denotation
 $\mathcal U_R^\dagger\mathcal U_R=\id$. It is therefore clean on
 every $r$ input with $a=0$. Nevertheless, its inner forward term
 $R$ takes input $\ket{10}$ to $\ket{01}$ by
 \stepref{sA-limits-swap}, so that inner mark is not clean.
 No logical cell is discarded or reused as blank in this example.
 \end{proof}
 \step{sA-limits-reference}{Item (iii): marginal restoration does not
 permit discarding correlations.}
 \begin{proof}\pf\ Identity satisfies (SA4) on every subspace and
 leaves the $a$ marginal of $\proj\Omega$ equal to $\id_a/2$.
 Dropping $a$ and allocating a fresh wire $a'$ in zero gives
 $(\id_q/2)\otimes\proj0_{a'}$. After identifying its two-dimensional
 space with that of $a$ solely to compare states, this differs from
 $\proj\Omega$:
 its expectation of $\proj\Omega$ is $1/4$, whereas the original
 expectation is $1$. Even replacement by a state with the same
 marginal $\id_a/2$ gives $\id_{qa}/4$ and the same discrepancy.
 Thus preservation of a marginal does not preserve the joint state
 under deletion and replacement.
 \end{proof}
 \step{sA-limits-unitary}{Item (iv): a unitary reduced channel need
 not reset the machine.}
 \begin{proof}\pf\ With the constant keep $X_r$, conjugation by
 $\operatorname{CNOT}_{r\to a}$ gives $W=X_rX_a$: on a basis pair
 $(y,z)$ the result is $(y\oplus1,z\oplus1)$. Hence
 $W(\psi_r\ket0_a)=X\psi_r\ket1_a$. The reduced channel is
 conjugation by $X$, but $a$ is never returned to zero.
 \end{proof}
 \qedstep
 \begin{proof}\pf\ Steps \stepref{sA-limits-subspace}--\stepref{sA-limits-unitary}
 establish the four assertions. All keeps are of
 Definition~\ref{def:keep}, including constant-zero and constant-one
 Boolean functions, and each body is identity or an explicitly
 displayed Boolean-gate term admitted by Definition~\ref{def:tiera}.
 \end{proof}
\end{proof}

\subsection{A classical sufficient condition and its limits}

The compute--copy--uncompute organization is stated in Bennett's abstract
\cite[abstract, journal p.~525, PDF p.~1]{Bennett1973}.
The following quantum extension and its limitations are proved directly.

\begin{theorem}[Compute--copy--uncompute with an untouched target]
\label{thm:sA-classical}
Let $X,r,M$ be finite pairwise disjoint registers, and let $m_0$ be a
computational-basis configuration of $M$. Suppose $U=U_{XM}\otimes\id_r$
is a total unitary, and for every computational-basis value $x$ of $X$
there are a unit vector $\phi_x\in\cH_X$ and a computational-basis
configuration $m_x$ of $M$ such that, in order $X,M$,
\[
 U_{XM}(\ket x\ket{m_0})=\ket{\phi_x}\ket{m_x}.          \tag{SA5}
\]
Let $v:\{0,1\}^{|M|}\to\{0,1\}^{|r|}$ be a total Boolean function
with the specified finite circuit of Definition~\ref{def:keep}, and
let $K=K_{v,r}$, with $M$ the whole keep-control list. Any additional
machine controls must be included in $M$ and in (SA5).
Then $U^\dagger KU$ is clean and induces the unitary
\[
 \ket{x,y}\longmapsto\ket{x,y\oplus v(m_x)}             \tag{SA6}
\]
on $Xr$, for every quantum input on $Xr$ and any external reference.
The action on inputs with $M$ in $\ket{m_0}$, and hence (SA6),
depends only on (SA5) and independence of $r$, not on the choice of
total extension of $U_{XM}$ away from these inputs. No independence of
the total operator $U^\dagger KU$ on other input states is asserted.

In particular, (SA5) holds with $\phi_x=\ket x$ for a classical
computation preserving the input register $X$. It also
holds whenever $U_{XM}$ is a total computational-basis permutation,
including a fixed finite computation composed only of reversible classical
steps, independent of the unread keep target $r$. Literal preservation
of $X$ is therefore unnecessary for the stronger sufficient condition.

Two enlargements of this universal claim are false. First, allowing
non-basis machine output in (SA5) fails for $X$ empty, $M=a$, $m_0=0$,
$U=H_a\otimes\id_r$, $K=\operatorname{CNOT}_{a\to r}$:
\[
 W\ket{0,0}_{ar}=
   (\ket0_a\ket+_r+\ket1_a\ket-_r)/\sqrt2,
 \qquad\ket\pm=(\ket0\pm\ket1)/\sqrt2.                \tag{SA7}
\]
Second, allowing $U\ket{x,y,m_0}=\ket{x,y,m_{x,y}}$ to depend on the
keep target fails for $M=a$, $U=\operatorname{CNOT}_{r\to a}$,
$K=\operatorname{CNOT}_{a\to r}$: input $\ket{1,0}_{ra}$ ends at
$\ket{0,1}_{ra}$. These are counterexamples to enlarged universal
claims, not necessity claims for each individual clean operation.
For an entangled-output failure upon allowing unrestricted outputs in
place of (SA5), see the separately proved Theorem~\ref{thm:sA-four}.
\end{theorem}
\begin{proof}
 \step{sA-class-basis}{For every basis pair $(x,y)$ the three stages
 have the asserted action.}
 \begin{proof}
  \step{sA-class-compute}{The forward stage gives
  $\ket{\phi_x}_X\ket y_r\ket{m_x}_M$.}
  \begin{proof}\pf\ This is (SA5) and $U$ acting as identity on $r$.
  \end{proof}
  \step{sA-class-copy}{The keep gives
  $\ket{\phi_x}_X\ket{y\oplus v(m_x)}_r\ket{m_x}_M$.}
  \begin{proof}\pf\ The machine factor is one basis configuration,
  so in the sector sum defining $K$ only its $m_x$ summand acts.
  \end{proof}
  \step{sA-class-uncompute}{The inverse gives
  $\ket x_X\ket{y\oplus v(m_x)}_r\ket{m_0}_M$.}
  \begin{proof}\pf\ Total unitarity and (SA5) imply
  $U_{XM}^\dagger(\ket{\phi_x}\ket{m_x})=\ket x\ket{m_0}$.
  Independence of $r$ permits applying this equality to the modified
  target value. No unspecified forward value of $U_{XM}$ is used.
  \end{proof}
  \qedstep
  \begin{proof}\pf\ Compose \stepref{sA-class-compute}--\stepref{sA-class-uncompute}.
  \end{proof}
 \end{proof}
 \step{sA-class-unitary}{The induced operation on $Xr$, with $M$
 initially in $\ket{m_0}$, is unitary on every $Xr$ input and is
 independent of the total extension.}
 \begin{proof}\pf\ For a fixed $x$, XOR by $v(m_x)$ permutes the
 basis of $r$ and squares to identity. Consequently (SA6) is a
 permutation of the full basis of $Xr$. The equality in
 \stepref{sA-class-basis} extends by linearity, also after tensoring
 an identity on a reference. Its expression uses only $m_x$; the
 inverse equality there holds for every unitary extension agreeing
 with (SA5). Thus $M$ is returned to $m_0$ with the claimed induced
 unitary.
 \end{proof}
 \step{sA-class-permutation}{Reversible classical steps with an
 untouched target satisfy the sufficient condition.}
 \begin{proof}\pf\ A finite product of computational-basis
 permutations is a computational-basis permutation. For each input
 $(x,m_0)$ its output is one basis pair $(z_x,m_x)$, so take
 $\phi_x=\ket{z_x}$. This verifies (SA5), whether or not $z_x=x$.
 \end{proof}
 \step{sA-class-H}{Equation (SA7) disproves deletion of the basis-output
 requirement.}
 \begin{proof}\pf\ On $\ket{00}_{ar}$ the forward $H_a$ gives
 $(\ket{00}+\ket{10})/\sqrt2$. The keep gives
 $(\ket{00}+\ket{11})/\sqrt2$. Applying $H_a$ again gives
 $(\ket{00}+\ket{10}+\ket{01}-\ket{11})/2$, which is (SA7).
 Its $a=1$ probability is $1/2$. The forward machine state was the
 non-basis vector $\ket+$, while the keep target was untouched by $U$.
 \end{proof}
 \step{sA-class-target}{The target-dependent example disproves deletion
 of the unread-target requirement.}
 \begin{proof}\pf\ The three CNOTs give
 $\ket{10}_{ra}\mapsto\ket{11}_{ra}\mapsto\ket{01}_{ra}
 \mapsto\ket{01}_{ra}$. Thus $a$ is one at the end although it
 was initially zero. The forward rule is nevertheless the
 basis-valued rule $\ket{y,0}\mapsto\ket{y,y}$.
 \end{proof}
 \qedstep
 \begin{proof}\pf\ By \stepref{sA-class-basis}--\stepref{sA-class-target}.
 The failures arise after deleting the indicated hypothesis; they do
 not contradict the sufficient statement with all hypotheses retained.
 \end{proof}
\end{proof}

\subsection{A four-qubit two-copy circuit, its leakage, and its channels}

\begin{theorem}[Four-qubit calculation]\label{thm:sA-four}
Let $q,r,a,b$ be distinct qubits, with $Q=qr$ and $M=ab$ initially
$\ket{00}$. The copying operations below are the explicit Boolean-gate
atoms of Definition~\ref{def:tiera}, not the full-snapshot source-read
translation of Definition~\ref{def:joint-tick}. Set
\[
 U=\operatorname{CNOT}_{q\to b}H_q\operatorname{CNOT}_{q\to a},
 \quad K=\operatorname{CNOT}_{b\to r},\quad W=U^\dagger KU,
 \quad P_\pm=\proj\pm_r.
\]
Define $E_{00}\xi=\xi_{qr}\otimes\ket{00}_{ab}$. On every input
$\xi\in\cH_{qr}$,
\[
 WE_{00}\xi=B_0\xi\otimes\ket{00}_{ab}
             +B_1\xi\otimes\ket{10}_{ab},\qquad
 B_0=\id_q\otimes P_+,\quad B_1=X_q\otimes P_-.         \tag{SA8}
\]
For any finite untouched reference $R$, the same equality holds for
$\xi\in\cH_{qrR}$ with $W,E_{00},B_0,B_1$ each tensored with
$\id_R$, and the factors canonically reordered by their names.
Consequently the following assertions hold.
\begin{enumerate}
 \item For $r$ initially $0$ and every unit vector $\psi_q$, the final
 state, in order $q,r,a,b$, is
 \[
  \frac{\ket\psi_q\ket+_r\ket0_a+
              X\ket\psi_q\ket-_r\ket1_a}{\sqrt2}
          \otimes\ket0_b.                              \tag{SA9}
 \]
 The probability that $M\ne00$ is $1/2$, independently of the input
 density operator of $q$ and of any reference: the reference input
 is $\rho_{qR}\otimes\proj0_r$, with $\rho_{qR}$ arbitrary.
 The factor $b$ is always zero.
 \item Dropping $M$ gives the channel on $qr$
 \[
  \Phi(\sigma)=B_0\sigma B_0^\dagger+B_1\sigma B_1^\dagger.
                                                               \tag{SA10}
 \]
 It is not a unitary channel on the full space $\cH_{qr}$. For $r=0$,
 \[
  \Phi(\rho\otimes\proj0_r)
   =\tfrac12\rho\otimes P_++\tfrac12X\rho X\otimes P_-.
                                                               \tag{SA11}
 \]
 Every pure $q$ input yields a rank-two $qr$ output of purity $1/2$.
 \item For any fixed density operator $\tau$ initializing $r$, with
 product input $\rho_q\otimes\tau_r$, put
 $p_\pm=\tr(P_\pm\tau)$. Then
 \[
 \begin{split}
  \Phi(\rho\otimes\tau)&=p_+\rho\otimes P_+
                         +p_-X\rho X\otimes P_-,\\
  \tr_r\Phi(\rho\otimes\tau)&=p_+\rho+p_-X\rho X.
 \end{split}                                                  \tag{SA12}
 \]
 The channel in the second line is unitary on $q$ if and only if
 $\tau=P_+$ or $\tau=P_-$. The map from $q$ to retained $qr$ in
 the first line is isometric if and only if the same condition holds;
 the word unitary here applies only after identifying the fixed $r$
 factor with a one-dimensional space. For $\tau=P_+$ the action is
 identity on $q$ and $M$ resets. For $\tau=P_-$ the action is $X$
 on $q$ and $M$ is the fixed state $\ket{10}$, so it does not reset.
 \item When $r=0$, the $q$ marginal is $(\rho+X\rho X)/2$.
 It equals $\rho$ precisely when $X\rho X=\rho$; among pure states
 these are exactly the $X$ eigenstates. This does not make (SA11)
 pure. If the intervening $H_q$ is omitted, the marked operation on
 the same initialized machine is instead CNOT from $q$ to $r$,
 with $M$ reset on every $qr$ input.
\end{enumerate}
Moreover
$U(\ket0_q\ket{00}_{ab})
=(\ket0_q\ket{00}_{ab}+\ket1_q\ket{01}_{ab})/\sqrt2$, which is not a
product vector across $q$ and $ab$; hence, with $X=q$, $M=ab$ and target
$r$ untouched by $U$, unrestricted forward outputs cannot replace (SA5)
in Theorem~\ref{thm:sA-classical}.
The example therefore refutes universal cleanliness for Tier~A, even
with an unread keep target, while preserving its total unitary soundness.
\end{theorem}
\begin{proof}
 \step{sA-four-conjugate}{On the full four-qubit space,
 $W=P_+\otimes\id_{qab}+P_-\otimes X_qX_aZ_b$, with the
 tensor factors placed on their indicated wires.}
 \begin{proof}
  \step{sA-four-K}{$K=P_+\otimes\id_{qab}+P_-\otimes Z_b$,
  with identities on $q,a$ in the second summand.}
  \begin{proof}\pf\ The target Pauli $X_r$ acts by $+1$ on
  $\ket+$ and by $-1$ on $\ket-$. In the latter sector,
  controlled-$X_r$ contributes $(-1)^b$, the action of $Z_b$.
  \end{proof}
  \step{sA-four-Z}{$U^\dagger Z_bU=X_qX_aZ_b$.}
  \begin{proof}\pf\ Conjugating $Z_b$ by
  $\operatorname{CNOT}_{q\to b}$ gives $Z_qZ_b$: on a basis
  pair the multiplier is $(-1)^{b\oplus q}=(-1)^b(-1)^q$.
  Next $H_qZ_qH_q=X_q$ by the matrices in
  Definition~\ref{def:circuit}. Finally conjugating $X_q$ by
  $\operatorname{CNOT}_{q\to a}$ gives $X_qX_a$, because
  the basis pair $(q,a)$ becomes $(q\oplus1,a\oplus1)$.
  The other factors commute with gates on disjoint wires.
  \end{proof}
  \qedstep
  \begin{proof}\pf\ $U$ does not act on $r$, so conjugate
  \stepref{sA-four-K} by $U$ and use \stepref{sA-four-Z}.
  \end{proof}
 \end{proof}
 \step{sA-four-isometry}{Equation (SA8) holds.}
 \begin{proof}\pf\ In \stepref{sA-four-conjugate}, the first
 summand leaves $\ket{00}_{ab}$ fixed. The second takes it to
 $\ket{10}_{ab}$, since $Z_b\ket0_b=\ket0_b$ and
 $X_a\ket0_a=\ket1_a$. The remaining operators are $B_0,B_1$.
 Each identity remains valid after tensoring with a reference identity.
 \end{proof}
 \step{sA-four-state}{Equation (SA9) and the leakage probability hold.}
 \begin{proof}\pf\ Substitute $\xi=\psi\otimes\ket0_r$ in
 (SA8) and use $P_\pm\ket0=\ket\pm/\sqrt2$.
 The two machine vectors in (SA8) are orthogonal, and
 $B_1^\dagger B_1=\id_q\otimes P_-$. Thus the $a=1$
 probability for any $qr$ input $\sigma$ is
 $\tr[(\id_q\otimes P_-)\sigma]$. For $r=0$ this is $1/2$,
 even if $q$ has a reference, since its input trace is one.
 Both summands have $b=0$.
 \end{proof}
 \step{sA-four-channel}{Equations (SA10)--(SA12) hold and define
 trace-preserving completely positive maps.}
 \begin{proof}\pf\ Tracing the orthogonal machine vectors in
 (SA8) removes their cross terms and yields (SA10).
 Moreover $B_0^\dagger B_0+B_1^\dagger B_1
 =\id_q\otimes(P_++P_-)=\id_{qr}$; the Kraus calculation in
 Theorem~\ref{thm:sA-sound} proves the channel properties.
 For every density operator $\tau$ on $r$,
 $P_\pm\tau P_\pm=\tr(P_\pm\tau)P_\pm$ because these
 projectors have rank one. Substitution gives (SA12); taking
 $\tau=\proj0$ gives $p_+=p_-=1/2$ and (SA11).
 \end{proof}
 \step{sA-four-rank}{The full $qr$ channel is not unitary, and every
 pure input on $q$ with $r=0$ has the asserted rank and purity.}
 \begin{proof}\pf\ In (SA11) for $\rho=\proj\psi$, the unit
 vectors $\psi\otimes\ket+$ and $X\psi\otimes\ket-$ are
 orthogonal because their $r$ factors are orthogonal. The two nonzero
 eigenvalues are $1/2,1/2$, so the rank is two and purity is
 $1/4+1/4=1/2$. Unitary conjugation preserves rank-one projections,
 whereas this channel takes the pure input $\proj\psi\otimes\proj0$
 to rank two. It cannot be unitary on $qr$.
 \end{proof}
 \step{sA-four-unitary-when}{The exact unitary and isometric conditions
 for a fixed target input are those in assertion 3.}
 \begin{proof}
  \step{sA-four-extreme}{\case{$p_+=1$ or $p_-=1$.}}
  \begin{proof}\pf\ Positivity of $\tau$ and its zero expectation
 on the complementary projection force its support into the respective
 rank-one eigenspace, by the positive-support argument of
 Theorem~\ref{thm:sA-clean}. Thus $\tau=P_+$ or $P_-$ respectively.
 Equation (SA8) becomes $\psi\ket+\ket{00}$ or
 $X\psi\ket-\ket{10}$ respectively. These are the stated isometries,
 induced unitaries on $q$, and machine outputs.
  \end{proof}
  \step{sA-four-interior}{\case{$0<p_+<1$, so $0<p_-<1$.}}
  \begin{proof}\pf\ For input $\rho=\proj0_q$, the $q$ output
  in (SA12) is $p_+\proj0+p_-\proj1$, of rank two; it cannot
  be a unitary channel. The retained $qr$ output has two nonzero
  eigenvalues $p_+,p_-$ on orthogonal vectors, so the map from $q$
  to $qr$ cannot be isometric either. The cases exhaust all possible
  $p_\pm$ because they are nonnegative and sum to one.
  \end{proof}
  \qedstep
  \begin{proof}\pf\ By \stepref{sA-four-extreme} and
  \stepref{sA-four-interior}.
  \end{proof}
 \end{proof}
 \step{sA-four-fixed}{The fixed states of the $q$ marginal for $r=0$
 are exactly those asserted.}
 \begin{proof}\pf\ The equality $(\rho+X\rho X)/2=\rho$ is
 equivalent to $X\rho X=\rho$. For a rank-one projection this
 says $X\psi=\lambda\psi$ for a scalar of modulus one; applying
 $X$ twice gives $\lambda^2=1$. Conversely either $X$ eigenstate
 satisfies the equality. The rank calculation in
 \stepref{sA-four-rank} still applies to its retained $qr$ output.
 \end{proof}
 \step{sA-four-noH}{Deleting the intermediate Hadamard gives a clean
 CNOT on $qr$.}
 \begin{proof}\pf\ Put $U_0=\operatorname{CNOT}_{q\to b}
 \operatorname{CNOT}_{q\to a}$. On every basis input
 $\ket{x,y,0,0}_{qrab}$, the forward circuit gives
 $\ket{x,y,x,x}$. The keep gives $\ket{x,y\oplus x,x,x}$,
 and the inverse gives $\ket{x,y\oplus x,0,0}$.
 Linear extension proves the claimed action on arbitrary $qr$ inputs.
 \end{proof}
 \step{sA-four-product-limit}{The forward body can entangle data
 and machine while leaving the keep target untouched.}
 \begin{proof}\pf\ On $\ket0_q\ket{00}_{ab}$ its action is
 $(\ket0_q\ket{00}_{ab}+\ket1_q\ket{01}_{ab})/\sqrt2$.
 Tracing $ab$ gives $\id_q/2$, of rank two, whereas a product pure
 vector has rank-one $q$ marginal. Thus this forward vector is not
 of the form (SA5). The operator $U$ never acts on $r$, yet for
 $r=0$ the leakage is $1/2$ by \stepref{sA-four-state}. Hence
 unrestricted forward outputs cannot replace (SA5) in the universal
 cleanliness assertion. This does not claim (SA5) necessary for each
 individually clean operation.
 \end{proof}
 \qedstep
 \begin{proof}\pf\ By \stepref{sA-four-isometry}--\stepref{sA-four-product-limit}.
 The nonzero leakage in \stepref{sA-four-state} disproves universal
 cleanliness, while all factors are total unitaries covered by
 Theorem~\ref{thm:sA-sound}. Their explicit copying atoms do not
 supply an equality with a different source-read translation.
 \end{proof}
\end{proof}

\notedsection{Resources for structured programs (Tier~A)}{sec:resources}
  {This section studies \earlierforms.}

\subsection{Exact resources of the prescribed expansion}

\begin{remark}[Reading the resource bounds]
The content of Theorem~\ref{thm:resA-resources} for most readers is (RA1),
(RA2) and (RA4): in the expansion that literally writes body, keep,
reversed body at every mark, work doubles per nesting level, exactly;
retained wires do not multiply, because inverses reuse the same named
records. Two warnings govern everything in this section. First, ``tight''
always means tight for that prescribed, unsimplified expansion: with an
empty keep the expanded circuit can already denote the identity, so
nothing here lower-bounds the best equivalent circuit, parallel depth,
or time at bounded space. Second, the remaining clause, (RA3), is a
routine but lengthy verification that a classical Turing machine can
print the expansion without random access; it is logically separate from
(RA1)--(RA2) and can be skipped on a first reading. No pebble-game
trade-off continuum is proved in this paper: the only proved alternatives
to eager expansion are the restricted delayed construction of
Theorem~\ref{thm:resA-delayed} and deletion of an inner inverse exactly
when the commutator test (RA9) (or its subspace form (RA10)) holds.
\end{remark}

\begin{theorem}[Expansion resources, not optimal synthesis]
\label{thm:resA-resources}
Assume all hypotheses of Theorem~\ref{thm:sA-sound}, including finite
static syntax, total completions, retained logical registers, the
exclusions of Definition~\ref{def:tiera}, and specified exact clean
circuits over $\mathcal G^{\pm}$. Use Definition~\ref{def:costs}, with
no cancellation or circuit simplification. In particular all tick
decoders satisfy the canonical-field, live-flag and fault requirements
on every code, and exclude allocation and measurement modes on every
code. Marks pair statically; marked bodies contain no allocation, drop,
observed measurement, external console observation or actual stop.
The menu extension requires base Tier~A entries with supplied exact
finite-alphabet lists and retained disjoint selectors. A compiled menu
is one leaf for $C,L,d$ and for the following sum: its payload gate
costs are not counted again. For an atom occurrence $v$
let $h(v)$ count its enclosing marks. For a keep occurrence $v$ let
$h(v)$ count its strictly enclosing marks, excluding the mark whose
keep it is. Write $c_v$ for its supplied circuit length. Then
\[
 \begin{gathered}
 C(I)=0,\qquad C(a)=c_a,\qquad C(P;R)=C(P)+C(R),\\
 C(\mathsf M(P,K))=2C(P)+c_K,\qquad
 C(P)=\sum_v2^{h(v)}c_v\leq2^{d(P)}L(P).             \tag{RA1}
 \end{gathered}
\]
The sum is an exact count, hence both an upper and a matching lower
bound for this expansion. Under serial execution, time in gate ticks is
$C(P)$ by the definition of that unit. Define $q_{\rm in}$ as the
number of data input wires in the enclosing outer program, so that its
fixed data pool has $q=q_{\rm in}+q_{\rm new}(P)$ wires.
Classical preparation of the supplied gate lists is a separate cost;
every gate in those lists is included in the execution count.

For the snapshot representation assume one shared $s$-bit configuration,
common sample width $k\geq1$, and common upper bounds $\beta,\Delta$
on the device circuits of Definition~\ref{def:device} across all tick
occurrences. Use globally distinct syntactic records. A sufficient
total wire count is
\[
 \begin{split}
 N(P)={}&q_{\rm in}+q_{\rm new}(P)+s+(s+k)n_{\rm tick}(P)\\
       &{}+s n_{\rm cont}(P)+m_{\rm extra}
         +\sum_\nu\delta_\nu+a_* .
 \end{split}                                         \tag{RA2}
\]
Here the sum counts disjoint retained menu selectors of widths
$\delta_\nu$, excluded from $m_{\rm extra}$ by convention in this
display. Every payload tick and continuation occurrence is counted once
for record allocation, across all entries and menu occurrences. If a retained
selector is computed by a subterm of $P$, that subterm's tick and continuation
records are already included in $n_{\rm tick}(P)$ and $n_{\rm cont}(P)$,
any other logical wires it retains are included in $m_{\rm extra}$,
and its clean scratch is included in $a_{\rm other}$ (or in the tick
and keep terms of (RA2a) when it consists of those constructions).
A selector supplied from outside $P$ contributes only its retained width
$\delta_\nu$ here; the cost of computing it is not part of this display. Inverses reuse these names.
The maximum local clean scratch requirement is $a_*$. For the concrete
implementations constructed below it can be chosen to satisfy
\[
 a_*\leq\max\left\{
 a_{\rm other},\ c_0(\beta+\Delta+s+k+1),\
 \max_K c_0\kappa_K,\
 \max_\nu\bigl(a_\nu^{\rm pay}+c_0(\delta_\nu+1)\bigr)
 \right\}.                                         \tag{RA2a}
\]
The absolute constant $c_0$ is sufficiently large. For a keep $K$,
$\kappa_K$ is the maximum, over its Boolean constituent circuits, of
one plus gate count plus input and output widths; an empty maximum
is zero. The number $a_\nu^{\rm pay}$ is the maximum clean workspace
of the supplied entry lists. Its base Tier~A constituents obey the
same bounds without the menu term. The declared number $a_{\rm other}$
bounds scratch for all other supplied atoms, including completed RTM
steps; no complexity bound for arbitrary supplied completions is assumed.
If arbitrary supplied circuits are used instead of these constructions,
their declared scratch widths must replace the corresponding bounds.
If every non-input wire in (RA2) is freshly zero-initialized and each
initialization costs one tick, add $N(P)-q_{\rm in}$ ticks; preparation
of a nonzero initial machine configuration is additional.
For
\[
 \mathcal R=\beta+\Delta+s+k+(q^2+kq+1)(s+k+1),
\]
a supplied fragment of $T_0$ joint ticks, including any declared timeout
tick, has $T_0(s+k)$
record bits and an implementation of $O(T_0\mathcal R)$ gates, before
multiplication by enclosing-mark factors in (RA1). The claim of ordinary
behavior at each forward tick additionally requires its logical records
blank at that use. A static menu with $0\leq j\leq2^\delta$ entries
of supplied exact clean $\mathcal G^{\pm}$ lengths $g_i$ and
a preserved $\delta$-bit selector disjoint from all payload and scratch
wires costs
$O(\sum_{i<j}g_i+j(\delta+1))$ gates, with clean scratch for one
selector flag and its computation of width $O(\delta+1)$, in addition
to the maximum payload scratch used simultaneously with that flag.

Use an explicit uncompressed ordered syntax tree encoded in prefix order
by constructor tokens with their arities and delimited atom/keep lists. Gate descriptions
are nonempty finite strings with a constant-width gate tag and binary
wire labels. If $\sigma$ is the total bit length of this input and
$\ell$ the maximum gate-description length (zero if there are no gates),
a multitape classical TM can print the
expansion in
\[
 O((\sigma+C(P)+1)^3(\ell+1)^2)                       \tag{RA3}
\]
steps. This deliberately loose bound excludes construction of the supplied
lists, and does not assume unit-cost random access.
For $P_0=a$, $c_a>0$, and $P_{j+1}=\mathsf M(P_j,K)$ with a fixed
keep circuit of length $c_K\geq0$,
\[
 C(P_j)=2^jc_a+(2^j-1)c_K.                            \tag{RA4}
\]
Thus the exponential dependence on nesting depth is attained for the
specified expansion. These statements give no lower bound on the
smallest equivalent circuit, circuit depth under parallel execution,
or optimal time under a space restriction.
\end{theorem}
\begin{proof}
 \step{resA-count-ind}{The recurrences and weighted sum (RA1) hold.}
 \begin{proof}
  \step{resA-count-base}{\case{Identity or an atom.}}
  \begin{proof}\pf\ Identity prints no gates. An atom prints its
  supplied list once, and has $h=0$.\end{proof}
  \step{resA-count-seq}{\case{A sequence, with the assertion for its
  two proper subterms.}}
  \begin{proof}\pf\ Concatenation adds lengths and retains each
  occurrence's number of enclosing marks.\end{proof}
  \step{resA-count-mark}{\case{$\mathsf M(P,K)$, with the assertion
  for $P$.}}
  \begin{proof}\pf\ The list is the body, keep, and reversed adjoint
  body. The alphabet is closed under adjoints, so reversal preserves
  length. Every body occurrence has one additional enclosing mark and
  appears twice as often; the new keep appears once with $h=0$.
  This proves both the recurrence and the weighted sum.\end{proof}
  \qedstep
  \begin{proof}\pf\ Structural induction gives the equality. For
  every occurrence $h(v)\leq d(P)$ and $c_v\geq0$, so summing
  $2^{h(v)}c_v\leq2^{d(P)}c_v$ gives the inequality.\end{proof}
 \end{proof}
 \step{resA-count-wires}{Equation (RA2) suffices without assuming
 logical machine cleanliness.}
\begin{proof}
 \step{resA-count-wires-logical}{The logical-wire contribution in
 (RA2) counts each required retained factor once.}
 \begin{proof}\pf\ Count the data wires, shared $s$-bit
 configuration, one $s+k$-bit cell per joint-tick occurrence, one
 $s$-bit cell per continuation occurrence and declared extra logical
 wires. An occurrence's adjoint reuses its names, so repetitions
 in (RA1) add no wires. All logical wires remain present.\end{proof}
 \step{resA-count-wires-selectors}{Menu selectors and payload records
 have exactly the declared additional allocation.}
 \begin{proof}\pf\ Add the disjoint selector widths explicitly.
 Selector-computation records and other logical wires are already
 included in the declared counts. Although a menu gate list is
 one leaf, its payload record declarations remain logical wires
 and are counted as stipulated in (RA2).\end{proof}
 \step{resA-count-wires-scratch}{One implementation pool of width
 $a_*$ suffices independently of logical-record reset.}
 \begin{proof}\pf\ Every atom/keep implementation and its inverse
 restores implementation scratch on every logical input.
 Lemma~\ref{lem:scratch} therefore permits sequential reuse of
 a pool of the maximum width. This assumption concerns
 implementation scratch, not cleanliness of logical machine records.
 \end{proof}
 \step{resA-count-wires-init}{The optional zero-initialization charge
 is $N(P)-q_{\rm in}$.}
 \begin{proof}\pf\ These are exactly the non-input wires in (RA2),
 each charged one tick under the stated convention. Preparation of
 a nonzero machine configuration is additional.\end{proof}
 \qedstep
 \begin{proof}\pf\ Add \stepref{resA-count-wires-logical}--\stepref{resA-count-wires-scratch}
 for the declared allocation. The following construction bounds the
 scratch width; \stepref{resA-count-wires-init} gives the optional
 initialization cost.\end{proof}
\end{proof}
 \step{resA-count-tick}{The asserted tick and menu bounds hold.}
 \begin{proof}
  \step{resA-count-ticklocal}{One joint tick costs $O(\mathcal R)$.}
  \begin{proof}
   \pf\ We reuse the circuit of Lemma~\ref{eqv:tick-cost}, here with the
   sharper simultaneous-scratch count needed for (RA2a).
   \step{resA-count-decoder}{For a device tick, its quantum action
 $D_t$ has gate cost
 $O(\Delta+(q^2+kq+1)(s+k+1))$ and scratch
 $O(\Delta+s+k+1)$.}
\begin{proof}
 \step{resA-count-decoder-eval}{Computing the decoder costs
 $O(\Delta)$ gates and workspace.}
 \begin{proof}\pf\ Apply Lemma~\ref{lem:synthesis} to the canonical
 decoder on unchanged $C$; $\Delta$ includes its input and output
 ports. No global validity test is added: canonical output, legal
 live operands and fault handling on every code are hypotheses.
 Measurement and allocation modes are excluded here.\end{proof}
 \step{resA-count-decoder-predicates}{There are $O(q^2+kq+1)$
 predicates, each of size and workspace $O(s+k+1)$.}
 \begin{proof}
  \step{resA-count-decoder-gatepred}{There are $O(q^2+1)$ gate
  predicates, each of Boolean size $O(a_q+1)$.}
  \begin{proof}\pf\ Enumerate $O(q^2+1)$ gate actions. A predicate
  compares only mode, gate tag and at most two addresses. Equality
  is a conjunction of bit comparisons, of size $O(a_q+1)$.\end{proof}
  \step{resA-count-decoder-readpred}{There are $kq$ read
  predicates, each of Boolean size $O(a_q+a_k+1)$.}
  \begin{proof}\pf\ Enumerate $kq$ position/address pairs. Test mode,
  the one read address and whether the fixed position is below the
  read length. Scan the length from most to least significant bit,
  retaining equality and less-than bits. This has Boolean size
  $O(a_q+a_k+1)$.\end{proof}
  \qedstep
  \begin{proof}\pf\ Definition~\ref{def:device} gives $a_q\leq s$,
  and $a_k=\lceil\log_2(k+1)\rceil\leq k$ for $k\geq1$.
  Lemma~\ref{lem:synthesis} gives the same order of workspace,
  independently of the entire decoder's width.\end{proof}
 \end{proof}
 \step{resA-count-decoder-apply}{Each selected action has constant
 additional gate cost and scratch.}
 \begin{proof}\pf\ Hold its predicate flag and apply controlled-$H$,
 controlled-$T$ or controlled-CNOT by Lemma~\ref{lem:synthesis}.
 Targets lie in $QB_t$, disjoint from $C$, decoder, flag and evaluator
 workspace. Gate mode selects one gate; read mode selects its listed
 copies, whose distinct sample targets make their order immaterial.
 \end{proof}
 \step{resA-count-decoder-erase}{One predicate computation and
 constant gate scratch suffice alongside the decoder word.}
 \begin{proof}\pf\ Erase each predicate before reusing its workspace;
 its decoder inputs have not changed. After all actions erase the
 decoder, whose input $C$ has not changed. Each erasure is the
 inverse evaluation and has the same cost.\end{proof}
 \qedstep
 \begin{proof}\pf\ Sum the gate costs over
 $O(q^2+kq+1)$ predicates and add decoder evaluation and erasure.
 Sequential reuse in \stepref{resA-count-decoder-erase} gives
 scratch $O(\Delta+s+k+1)$.\end{proof}
\end{proof}
   \step{resA-count-update}{The update and the timeout obey the
   remaining asserted bounds.}
   \begin{proof}\pf\ After $D_t$, evaluate $f(C,B_t)$ into $J_t$,
   undo its Boolean workspace, then swap $C,J_t$. By
   Lemma~\ref{lem:synthesis} this costs $O(\beta+s+k)$ gates and
   scratch; a swap is three CNOTs per bit. It is the required $F_t$
   even on nonblank logical records. The separate timeout has
   $D_t=\id$ and update $h_\tau$: equality of the status with its
   active code, followed by conditional constant replacement, has
   $O(s)$ Boolean size. The same snapshot construction costs
   $O(s+k)$ and keeps its declared $s+k$ recording bits. A continuation
   replaces a fixed position field and costs $O(s)$ gates and scratch.
   Each is bounded by the tick term in (RA2a).\end{proof}
   \qedstep
   \begin{proof}\pf\ By \stepref{resA-count-decoder} and
   \stepref{resA-count-update}, sequential reuse gives $O(\mathcal R)$
   gates and $O(\beta+\Delta+s+k+1)$ scratch per tick.
   The $T_0$ listed ticks, including any timeout, have distinct cells.
   Their lengths add; blankness for forward interpretation is exactly
   the promise in Definition~\ref{def:joint-tick}.\end{proof}
  \end{proof}
  \step{resA-count-menu}{A menu entry costs $O(g_i+\delta+1)$.}
\begin{proof}
 \step{resA-count-menu-flag}{Compute equality of the selector with
 $i$ in one flag using $O(\delta+1)$ gates and workspace.}
 \begin{proof}\pf\ Conjoin the bit equalities and apply the clean
 Boolean evaluator of Lemma~\ref{lem:synthesis}. The capacity
 $j\leq2^\delta$ makes every $i<j$ a selector value, including
 the zero-width selector when it has an entry.\end{proof}
 \step{resA-count-menu-control}{Hold the flag while controlling the
 $g_i$ elementary gates, at constant overhead per gate.}
 \begin{proof}\pf\ Each supplied gate is $H,T,T^\dagger$ or CNOT.
 Lemma~\ref{lem:synthesis} supplies its exact controlled version.
 Selector and predicate workspace are disjoint from all targets,
 including payload implementation wires. Thus simultaneous
 workspace is $a_\nu^{\rm pay}+O(\delta+1)$, including the
 constant scratch of the controlled elementary gate.\end{proof}
 \step{resA-count-menu-false}{\assume{The flag is zero.}
 \prove{The complete controlled list is identity in this sector.}}
 \begin{proof}\pf\ Every controlled gate is identity in this
 sector; consequently the complete list is identity.\end{proof}
 \step{resA-count-menu-true}{\assume{The flag is one.}
 \prove{The complete controlled list equals the supplied payload
 circuit, including its phase, and restores its implementation scratch.}}
 \begin{proof}\pf\ The controlled list is exactly the supplied
 clean payload circuit, including its phase. Its implementation
 scratch is zero after the complete list.\end{proof}
 \step{resA-count-menu-reuse}{Erase the flag and reuse the scratch
 for the next entry.}
 \begin{proof}\pf\ The selector has not changed, so inverse
 evaluation erases the flag. The two sectors above restore payload
 scratch; linearity gives the same restoration on arbitrary selector
 superpositions. Lemma~\ref{lem:scratch} permits sequential reuse.
 \end{proof}
 \qedstep
 \begin{proof}\pf\ Add flag computation and erasure to the $g_i$
 controlled gates, then sum over entries for the asserted menu bound.
 Only supplied exact finite lists are used, so this covers no
 ideal-angle payload lacking such a list.\end{proof}
\end{proof}
  \step{resA-count-keep}{The remaining scratch bounds in (RA2a) hold.}
  \begin{proof}\pf\ For each Boolean constituent of a keep use the
  evaluator of Lemma~\ref{lem:synthesis} directly into the disjoint
  XOR target and erase its stored Boolean gate values. Its cost and
  scratch are $O(\kappa_K)$. Constituent keeps execute sequentially,
  so their maximum scratch suffices. Arbitrary other atom scratch is
  the supplied bound $a_{\rm other}$. Menu entries are base terms;
  Lemma~\ref{lem:scratch} bounds their scratch by the maximum of
  their base constituents. These observations and the preceding two
  steps establish (RA2a) for one fixed sufficiently large $c_0$.
  \end{proof}
  \qedstep
  \begin{proof}\pf\ By \stepref{resA-count-ticklocal} and
  \stepref{resA-count-menu} and \stepref{resA-count-keep}.\end{proof}
 \end{proof}
 \step{resA-count-compile}{A sequential-tape compiler has bound (RA3).}
 \begin{proof}
  \step{resA-compile-encoding}{\pick\ $N$ syntax-node identifiers
 of common width $N+1$, a child-before-parent order, and fixed-width
 gate fields of width $w=O(\ell+1)$. They can be prepared in
 $O((\sigma+1)^3)$ multitape steps, with $N=O(\sigma+1)$.}
\begin{proof}
 \step{resA-compile-encoding-ids}{There are $N=O(\sigma+1)$ nodes;
 assign the $i$th constructor the identifier $1^i0^{N+1-i}$.}
 \begin{proof}\pf\ Count constructor tokens in one input scan.
 Each consumes input symbols. The displayed distinct unary strings
 all have width $N+1$, and writing all of them costs $O((N+1)^2)$.
 \end{proof}
 \step{resA-compile-encoding-order}{A stack scan emits each node
 after its children, in $O((\sigma+1)^2)$ steps.}
 \begin{proof}
  \step{resA-compile-encoding-frames}{Use fixed-width $O(N+1)$ frames
  holding the identifier, constructor tag, up to two child identifiers
  and the number of children still awaited.}
  \begin{proof}\pf\ The prefix encoding supplies each constructor's
  arity. All identifier fields have width $N+1$, and the remaining
  fields have bounded width. Keep the stack head at its top.\end{proof}
  \step{resA-compile-encoding-leaf}{\assume{A leaf is completed.}
  \prove{Its record can be emitted, its frame popped, and its
  identifier recorded in the current parent frame, if any.}}
  \begin{proof}\pf\ It awaits no child. Emit its record, pop its
  frame and record its identifier in the current parent frame, if
  there is one.\end{proof}
  \step{resA-compile-encoding-child}{\assume{A nonleaf receives a
  completed child's identifier.}\prove{The child can be recorded;
  on the last child the node can be emitted and popped after all
  its children, and its identifier passed to its parent.}}
  \begin{proof}\pf\ Store the identifier and decrease the number
  awaited. On the last child, emit and pop the parent frame and
  pass its identifier to its parent. Thus all children have already
  been emitted when a node is emitted.\end{proof}
  \qedstep
  \begin{proof}\pf\ Each frame is pushed and popped once; there
  are $O(N)$ events. Reading or updating its constant number of
  fields, copying an identifier and returning the auxiliary tape
  head cost $O(N+1)$ per event. Add the input scan.\end{proof}
 \end{proof}
 \step{resA-compile-encoding-lists}{The owner-indexed tape of padded
 supplied lists has length $O((\sigma+1)^2)$.}
 \begin{proof}\pf\ Copy each delimited list under its owner
 identifier when encountered. Their original total length and the
 number of gates are at most $\sigma$, since descriptions are
 nonempty. A scan finds $\ell$; right-padding and copying cost
 $O((\sigma+1)(\ell+1))$. Owner headers cost $O((N+1)^2)$.
 Since $N,\ell\leq O(\sigma+1)$, the stated length follows.
 Fixed delimiters distinguish padded field ends.\end{proof}
 \step{resA-compile-encoding-lookup}{Retrieving a node's supplied
 list costs $O((\sigma+1)^2)$; all such retrievals cost
 $O((\sigma+1)^3)$.}
 \begin{proof}\pf\ Rewind and scan the owner-indexed tape.
 Compare unary fields synchronously and return the query head in
 $O(N+1)$ per header. This is proportional to header length, so
 it is included in the complete scan. Sum the tape-length bound
 over the $N$ nodes; charge each retrieval to its node below.\end{proof}
 \step{resA-compile-encoding-adjoint}{Adjoint conversion changes
 only a constant number of tag symbols in a gate field.}
 \begin{proof}\pf\ The constant-width tag ranges over
 $\{H,T,T^\dagger,\mathrm{CNOT}\}$. Its adjoint is a fixed table
 in finite control. Wire labels are unchanged, including the two
 labels of self-adjoint CNOT.\end{proof}
 \qedstep
 \begin{proof}\pf\ The preceding scans use a fixed finite number
 of tapes and no indexed access. They prepare the stated choices
 within $O((\sigma+1)^3)$ steps, with the subsequent list retrievals
 bounded separately by \stepref{resA-compile-encoding-lookup}.\end{proof}
\end{proof}
  \step{resA-compile-size}{For every processed node $R$,
  $C(R)\leq C(P)$; the expansion table has length at most
  $T=O((\sigma+1)(C(P)(\ell+1)+\sigma+1))$.}
  \begin{proof}\pf\ Every child length is at most its parent's
  length by (RA1) and nonnegativity. Following the finite ancestor
  chain proves the first assertion. Encode an expansion-table entry
  as its width-$(N+1)$ identifier, a delimiter, its expanded gate
  fields, and an end delimiter; table delimiters are distinct from
  gate symbols. Each entry has length $O(N+1+C(P)w)$. There are
  at most $N$ entries, giving $T$. Empty lists still have a header
  and an end delimiter.\end{proof}
  \step{resA-compile-node}{One node's expanded list is constructed
  and appended in $O(T+(\sigma+1)^2+(C(P)+1)(\ell+1)^2+1)$ steps.}
  \begin{proof}
   \step{resA-compile-lookup}{Every child lookup costs $O(T)$ steps.}
   \begin{proof}\pf\ Rewind the table tape and scan it sequentially.
   At each header compare its identifier against the query tape, then
   rewind the query head. This takes $O(N+1)$ per header. A failed
   match skips the entry, scanning each field once; a successful match
   copies the list to a separate tape as it is scanned. The header
   comparison and head return cost is bounded by a constant times
   header length. Thus a complete lookup, including the initial rewind,
   costs $O(T)$, not $O(NT)$. There are at most two child lookups
   per node.\end{proof}
   \step{resA-compile-reverse}{Reversed adjoint copying of an $m$-gate
   list costs $O((m+1)w^2)$.}
   \begin{proof}\pf\ First scan to the end, costing $O((m+1)w)$.
   From the end of each fixed-width gate field move left to its start,
   copy its symbols left-to-right to the output while changing its tag
   to the adjoint tag, then move left by at most $2w+O(1)$ to the
   preceding field's start. A fixed-width position counter may cost
   $O(w)$ per copied symbol even under a unary representation; resetting
   it costs $O(w)$ per field. This gives the stated conservative
   $O((m+1)w^2)$ bound including head movements. Each gate is visited
   once, in reverse gate order, while its wire-label bits retain their
   order. The zero-gate case only reads the delimiters.\end{proof}
   \qedstep
   \begin{proof}\pf\ Identity emits an empty list; an atom copies
   its supplied list; a sequence concatenates its two retrieved lists;
   a mark concatenates its body, supplied keep, and the reversed
   adjoint body. Locating the node's supplied list on the owner-indexed
   tape costs $O((\sigma+1)^2)$ by \stepref{resA-compile-encoding}.
   Only a constant number of tapes are needed because
   this is one node at a time, after its children have been stored.
   Every relevant supplied list and child list has at most $C(P)$
   gates. Sequential copying and the preceding reversal bound control
   all work except lookup. Rewinding the table and moving to its end
   to append the new entry takes $O(T)$ further steps. Apply
   \stepref{resA-compile-lookup} and \stepref{resA-compile-reverse}.
   \end{proof}
  \end{proof}
  \step{resA-compile-total}{The total bound is (RA3), including
  $C(P)=0$.}
  \begin{proof}\pf\ Process the prepared order once. If $t_i$ is
  the cumulative time after $i$ nodes, the preceding step gives
  $t_{i+1}-t_i=O(T+(\sigma+1)^2+(C(P)+1)(\ell+1)^2+1)$ uniformly in $i$.
  Sum over $N=O(\sigma+1)$ nodes and add preprocessing and the final
  root-list copy. The result is
  \[
   O\bigl((\sigma+1)T+(\sigma+1)^3+
       (\sigma+1)(C(P)+1)(\ell+1)^2\bigr).
  \]
  Substitute the bound on $T$ from \stepref{resA-compile-size};
  each summand is bounded by a constant times
  $(\sigma+C(P)+1)^3(\ell+1)^2$. When $C(P)=0$, there are still
  only $N$ node-processing events; empty body lists are stored once,
  not recursively regenerated for their two uses.\end{proof}
  \qedstep
  \begin{proof}\pf\ By \stepref{resA-compile-encoding}--\stepref{resA-compile-total}.
  \end{proof}
 \end{proof}
 \step{resA-count-tight}{Equation (RA4) and the claimed limitations hold.}
 \begin{proof}\pf\ The recurrence is $C_{j+1}=2C_j+c_K$ with
 $C_0=c_a$. Substitution and induction give (RA4). In particular,
 $C_j\geq2^jc_a$. Each counted gate must occur in an unchanged
 prescribed list, proving its matching lower bound and serial time.
 Printing requires at least one output action per gate because each
 gate description is nonempty. Definition~\ref{def:keep} permits a
 finite product with no entries: choose its empty circuit, so
 $K=\id$ and $c_K=0$. Already $P_1$ denotes
 $A_a^\dagger A_a=\id$, and every $P_j$, $j\geq1$, denotes
 identity. Its optimal gate count is zero although the expansion has
 $2^jc_a$ gates. Therefore an optimal-circuit lower bound does not
 follow. With $c_K=0$ the upper bound $2^dL$ itself is attained.
 \end{proof}
 \qedstep
 \begin{proof}\pf\ By \stepref{resA-count-ind}--\stepref{resA-count-tight}
 and the clean correctness of the expansion in
 Theorem~\ref{thm:sA-sound}.\end{proof}
\end{proof}

\subsection{Retaining intermediate computations}

\begin{theorem}[Time versus retained workspace]
\label{thm:resA-delayed}
Let $d\geq1$, let $b_0,\ldots,b_d,u_1,\ldots,u_d$ be nonnegative
integer widths, and let
$f_j:\{0,1\}^{b_{j-1}}\to\{0,1\}^{b_j}$ be total deterministic
functions. Let $X,R$ be disjoint external registers with
$|X|=b_0$, $|R|=b_d$. Supply total evaluator circuits $A_j$ of lengths
$a_j$ over $\mathcal G^{\pm}$, counted as in Definition~\ref{def:costs}.
On input $\ket x\ket{0^{b_j}}\ket{0^{u_j}}$ they produce
$\ket x\ket{f_j(x)}\ket{g_j(x)}$ for specified Boolean strings
$g_j(x)$ of width $u_j$. All evaluator implementation wires are included
in $u_j$; there is no additional implementation pool in this theorem.
For each $j$ supply a scratch-free $\mathcal G^{\pm}$ circuit of length
$k_j$ acting only on a $b_j$-bit output and a disjoint $b_j$-bit target,
with exact action $\ket{z,y}\mapsto\ket{z,y\oplus z}$ on every input.
The usual $b_j$ CNOTs are an admissible such list. Inverse lists have
the same length because this gate alphabet is adjoint-closed.
All evaluator output and private wires initially are zero.
Each evaluator acts only on its own input, output and private wires.
Only intermediate outputs, never the final external target, become
inputs of later evaluators. Earlier evaluator private wires remain
untouched until their specified inverse. No additional operation changes
retained values. Use the eager and delayed compositions of
Definition~\ref{def:delayed}, under these explicit hypotheses.
Let $F_j=f_j\circ\cdots\circ f_1$.

Both the eager and delayed constructions implement, on all states of
$XR$ with their specified auxiliary zero initialization, the same retained
isometry after padding their auxiliary lists to a common width and
permuting auxiliary factors to identify their names:
\[
 \ket{x,y}\ket0\longmapsto
 \ket{x,y\oplus F_d(x)}\ket0,                         \tag{RA5}
\]
including arbitrary external references. They need not be the same
unitary on nonzero auxiliary inputs. The exact unsimplified gate counts
and sufficient allocated auxiliary widths, excluding $X,R$, are
\[
 \begin{gathered}
 o_j=2a_j+k_j,\qquad E_1=o_1,\qquad E_j=2E_{j-1}+o_j,\\
 E_d=\sum_{j=1}^d2^{d-j}o_j,\qquad
 D_d=2\sum_{j=1}^da_j+k_d.                            \tag{RA6}
 \end{gathered}
\]
\[
 W_E=\sum_{j<d}b_j+\max_{1\leq j\leq d}(b_j+u_j),
 \qquad W_D=\sum_{j\leq d}(b_j+u_j).                 \tag{RA7}
\]
These are serial gate times as well as gate counts. If, uniformly in
$d,j$, $0<\alpha\leq a_j\leq A$ and $0\leq k_j\leq K$ for fixed
constants, then $E_d=\Theta(2^d)$ and $D_d=\Theta(d)$.
There is no universal strict separation between the two space counts;
they are sufficient allocations, not space lower bounds.

The same conclusions, including exact counts, hold for quantum
private outputs if the evaluator hypothesis is replaced by
\[
 A_j\bigl(\ket x\ket{0^{b_j}}\ket{0^{u_j}}\bigr)
   =\ket x\ket{f_j(x)}\ket{\gamma_{j,x}},             \tag{RA8}
\]
where each $\gamma_{j,x}$ is any unit vector, with the same widths,
alphabet, copying, locality, zero-initialization and no-intervention
assumptions. Thus quantum phases and non-basis private states are
permitted. The restrictions are exact preservation of the basis input,
an exact computational-basis output $f_j(x)$, and a product private
factor for each basis input. Disturbed inputs or an output entangled
with its private factor on a basis input are excluded.
\end{theorem}
\begin{proof}
 \step{resA-delay-oracle}{A compute--copy--inverse application of
 $A_j$ realizes a clean XOR by $f_j$ on its input and disjoint target.}
 \begin{proof}\pf\ On the basis input $\ket{x,y}$ and zero local
 wires, (RA8) gives
 $\ket x\ket{f_j(x)}\ket{\gamma_{j,x}}\ket y$.
 Copying changes only the last factor to $\ket{y\oplus f_j(x)}$.
 Applying $A_j^\dagger$ to the first three factors gives exactly
 $\ket x\ket0\ket0$, by total unitarity and (RA8).
 Linearity proves the assertion on arbitrary input/target states.
 The classical hypothesis is the special case
 $\gamma_{j,x}=\ket{g_j(x)}$.\end{proof}
 \step{resA-delay-eager}{The eager construction realizes (RA5).}
 \begin{proof}
  \step{resA-delay-eagerbase}{\case{$j=1$.}}
  \begin{proof}\pf\ Use \stepref{resA-delay-oracle}.\end{proof}
  \step{resA-delay-eagernext}{\assume{The previous composite is a
  clean XOR by $F_{j-1}$.}\prove{The level-$j$ composite is a clean
  XOR by $F_j$.}}
  \begin{proof}\pf\ Compute $F_{j-1}(x)$ into a zero intermediate.
  By the induction hypothesis every private wire used for that operation
  is reset. Apply the clean composite for $f_j$ with that intermediate as
  input and the current external target as target. It preserves the
  intermediate and changes the target by $f_j(F_{j-1}(x))$.
  Invert the previous composite to erase the intermediate; it acts
  independently of the changed target. This proves the basis action
  and zero restoration, then linearity proves the full isometry.\end{proof}
  \qedstep
  \begin{proof}\pf\ Induction using
  \stepref{resA-delay-eagerbase} and \stepref{resA-delay-eagernext}.
  \end{proof}
 \end{proof}
 \step{resA-delay-late}{The delayed construction realizes (RA5).}
 \begin{proof}\pf\ Fix $x$ and write $z_0=x$, $z_j=f_j(z_{j-1})$.
 After all forward evaluators, the retained vector is
 $\ket x\bigotimes_{j=1}^d
 (\ket{z_j}\ket{\gamma_{j,z_{j-1}}})$, with the target unchanged.
 This follows by induction from (RA8) and locality. Copy $z_d$ into
 the target. The factors of the last evaluator still have exactly
 its forward output, so its inverse erases them. Repeat in reverse
 order; every earlier input was preserved. The target is outside all
 these inverses. Hence all local wires are zero and (RA5) holds on
 each basis pair. Extend linearly and tensor with identity on a reference.
 No private vector is copied or discarded.\end{proof}
 \step{resA-delay-counts}{Equations (RA6)--(RA7) hold.}
\begin{proof}
 \step{resA-delay-counts-oracle}{The $j$th compute--copy--inverse composite has length
 $o_j=2a_j+k_j$.}
 \begin{proof}\pf\ Count its length-$a_j$ evaluator, length-$k_j$
 copy and length-$a_j$ inverse. Adjoint closure preserves list length.
 \end{proof}
 \step{resA-delay-counts-eager}{The eager counts are
 $E_1=o_1$, $E_j=2E_{j-1}+o_j$ and
 $E_d=\sum_{j=1}^d2^{d-j}o_j$.}
 \begin{proof}\pf\ The eager procedure uses the preceding
 composite twice and the current compute--copy--inverse composite once. The base is one
 oracle. Substituting the sum for $E_{j-1}$ into the recurrence
 doubles its coefficients and appends $o_j$, proving the formula
 by induction.\end{proof}
 \step{resA-delay-counts-delayed}{The delayed count is
 $D_d=2\sum_{j=1}^da_j+k_d$.}
 \begin{proof}\pf\ Each evaluator and its inverse occurs exactly
 once, and only the last copy occurs. Sum their list lengths.
 \end{proof}
 \step{resA-delay-counts-pool}{The eager procedure can share a local
 pool of width $\max_j(b_j+u_j)$.}
 \begin{proof}\pf\ Keep all intermediate outputs outside this
 pool. Step~\stepref{resA-delay-eager} restores the pool before
 the level-$j$ compute--copy--inverse composite; \stepref{resA-delay-oracle} restores it
 before the preceding composite's inverse. Lemma~\ref{lem:scratch}
 gives zero restoration on that inverse pass as well. Induction
 therefore justifies each reuse at every level. Pad unused pool
 wires by identity.\end{proof}
 \step{resA-delay-counts-widths}{The sufficient widths are
 $W_E=\sum_{j<d}b_j+\max_j(b_j+u_j)$ and
 $W_D=\sum_{j\leq d}(b_j+u_j)$.}
 \begin{proof}\pf\ For eager execution add the reserved intermediate
 wires to the shared pool of \stepref{resA-delay-counts-pool}.
 Delayed execution reserves a separate $b_j+u_j$ block for each
 evaluator. Neither allocation counts $X,R$. Renaming wires of
 a zero pool incurs no gates.\end{proof}
 \qedstep
 \begin{proof}\pf\ The first three steps give (RA6);
 \stepref{resA-delay-counts-widths} gives (RA7).\end{proof}
\end{proof}
 \step{resA-delay-orders}{The time orders and space qualifications hold.}
 \begin{proof}\pf\ The assumptions bound $o_j$ between $2\alpha$
 and $2A+K$. Sum $\sum_{j=1}^d2^{d-j}=2^d-1$ to obtain the eager
 order. Also $2\alpha d\leq D_d\leq2Ad+K$, proving the delayed
 order. For $b_j=b$, $u_j=u$, the displayed widths are
 $db+u$ and $d(b+u)$. They coincide when $u=0$ and differ by
 $(d-1)u$ otherwise. For general varying widths the sufficient
 formula $W_E$ need not even be less than $W_D$; no inequality
 between them has been used.\end{proof}
 \qedstep
 \begin{proof}\pf\ By \stepref{resA-delay-oracle}--\stepref{resA-delay-orders}.
 Classical basis private outputs satisfy (RA8), so this also proves
 the original restricted delayed-erasure target with its hypotheses.
 \end{proof}
\end{proof}

\subsection{The exact condition for deleting an inner inverse}

\begin{theorem}[A quantum deletion criterion and exact savings]
\label{thm:resA-deletion}
Let $A,P,B$ be specified total unitaries on a common finite retained
list $L=QM$, realized by Tier~A terms; let $K,K'$ be keeps of
Definition~\ref{def:keep}. There are no allocation, drop or observed
measurement operations in these terms. Consider the nested expression
\[
 \mathsf M\bigl(A;\mathsf M(P,K);B,K'\bigr).
\]
Here and below the same letters denote the terms' logical unitaries.
Delete the inner inverse $P^\dagger$ after $K$, retaining its records
until the outer inverse, to obtain $\mathsf M(A;P;K;B,K')$, where
$K$ is now a standalone atom. Write
\[
 U=BP^\dagger KPA,\quad V=BKPA,\quad H_0=B^\dagger K'B,
 \quad W=U^\dagger K'U,\quad \widetilde W=V^\dagger K'V.
\]
Then, with exact matrix phases,
\[
 W=\widetilde W\quad\Longleftrightarrow\quad
 [P,H_0]=0.                                          \tag{RA9}
\]
In particular it suffices that there is a tensor decomposition
$\cH_L=\cH_S\otimes\cH_{L\setminus S}$ over disjoint wire lists
with $P=P_S\otimes\id$ and $H_0=\id\otimes H_{L\setminus S}$,
even if $A,P,B$ contain quantum gates. More generally,
for any specified input isometry $J:\cH_{\rm in}\to\cH_L$,
equality of the retained isometries $WJ=\widetilde WJ$ is equivalent to
\[
 (PH_0P^\dagger-H_0)KPAJ=0.                          \tag{RA10}
\]
This weaker condition may suffice when (RA9) fails. These statements
are about retained operators first; any identical final partial trace
then preserves equality. They imply no reset of logical records on
their own and do not identify unequal dilations merely from equal
output channels.

For the resource assertion assume supplied exact clean circuits over
the adjoint-closed alphabet $\mathcal G^{\pm}$ as in
Definition~\ref{def:costs}. If $a,p,b,k,k'$ are the lengths of the
supplied expansions for $A,P,B,K,K'$, the two unsimplified counts are respectively
\[
 2(a+2p+k+b)+k',\qquad 2(a+p+k+b)+k',                 \tag{RA11}
\]
so precisely $2p$ gates and serial ticks are saved whenever the
transformation is justified. On a fixed retained $n$-wire logical list,
both can be implemented with $n+a_*^{\rm old}$ wires. Here
$a_*^{\rm old}$ is the original list's maximum local clean scratch;
the modified list uses only occurrences from that list, so
$a_*^{\rm new}\leq a_*^{\rm old}$.
Thus this local rule alone does not force a space increase. If it
replaces previously certified reuse of logical zero cells, those cells
must instead be retained separately or that reuse proved valid anew.
\end{theorem}
\begin{proof}
 \step{resA-delete-factor}{Put $F=KPA$. Then
 $W=F^\dagger PH_0P^\dagger F$ and
 $\widetilde W=F^\dagger H_0F$.}
 \begin{proof}\pf\ Substitute $U=BP^\dagger F$ and $V=BF$
 in the definitions and reverse factors under the adjoint.\end{proof}
 \step{resA-delete-full}{Equation (RA9) holds.}
 \begin{proof}\pf\ $F$ is unitary. Multiplying the two expressions
 in \stepref{resA-delete-factor} on the left by $F$ and right by
 $F^\dagger$ shows that equality is equivalent to
 $PH_0P^\dagger=H_0$. Multiplication on the right by $P$ makes
 this equivalent to $PH_0=H_0P$. Operators on disjoint tensor factors
 commute, proving the stated sufficient special case.\end{proof}
 \step{resA-delete-subspace}{Equation (RA10) is equivalent to
 $WJ=\widetilde WJ$.}
 \begin{proof}\pf\ Subtract the expressions in
 \stepref{resA-delete-factor}, multiply by $J$ on the right and
 by the invertible $F$ on the left. The result is exactly (RA10).
 The equality remains true when tensored with a reference identity;
 linearity and a final partial trace give the corresponding channel
 equality. Neither calculation asserts a zero record output.\end{proof}
 \step{resA-delete-resources}{The exact counts are (RA11).}
 \begin{proof}\pf\ The original outer body has length
 $a+2p+k+b$; the modified body has length $a+p+k+b$.
 Each is followed by its keep and its complete inverse, so apply
 (RA1). The logical names are unchanged; Lemma~\ref{lem:scratch}
 supplies the common clean scratch pool. Deletion removes occurrences
 both in the forward body and in its outer inverse; it introduces no
 new local implementation, proving the scratch inequality.
 Its hypothesis concerns
 implementation scratch, and does not certify logical cell reuse.
 \end{proof}
 \qedstep
 \begin{proof}\pf\ By \stepref{resA-delete-full}--\stepref{resA-delete-resources}.
 \end{proof}
\end{proof}

\begin{corollary}[Certified removal at multiple levels]
\label{cor:resA-flat}
Let a finite base Tier~A term, without derived menu atoms, have one
outermost mark with keep $K_*$. All resource lists are exact clean
$\mathcal G^{\pm}$ lists counted without simplification.
Let $S$ be the sum of supplied gate lengths of every atom and every
inner keep, each counted once. Suppose its expanded outer body can
be reduced to the list visiting each such occurrence once by successive
deletions of inner inverse lists. At each deletion, in the current
body write the unitary as $BP^\dagger F$, where the removed list
implements $P^\dagger$, and require
$[P,B^\dagger K_*B]=0$, or require (RA10) with $F$ in place of $KPA$
for the same fixed input embedding $J$. Retain every logical wire and
keep the supplied exact clean local circuits. Then the final outer
operation agrees with the original (on all retained inputs in the
commutator case, or on $J$ in the embedding case). Its count is
$2S+c_{K_*}$. Each deletion of a length-$p$ inverse list saves exactly
$2p$ gates. The time is therefore linear in this once-counted body size,
conditional on all the stated tests; no universal deletion rule is asserted.
\end{corollary}
\begin{proof}
 \step{resA-flat-one}{Each deletion preserves the specified equality.}
 \begin{proof}\pf\ The proof of Theorem~\ref{thm:resA-deletion}
 uses $F$ only as a unitary factor, so applies to the current factorization.
 Its two outer lists differ in length by twice the deleted length.
 \end{proof}
 \step{resA-flat-all}{The final count is $2S+c_{K_*}$.}
 \begin{proof}\pf\ Equality is transitive through the finite
 sequence of deletions. By hypothesis the final body contains each
 designated occurrence once, so has length $S$. Its outer expansion
 consists of this list, $K_*$, and its reversed adjoint list.\end{proof}
 \qedstep
 \begin{proof}\pf\ By \stepref{resA-flat-one} and \stepref{resA-flat-all}.
 \end{proof}
\end{proof}

\begin{counterexample}[A clean inner evaluation does not justify deletion]
\label{ex:resA-quantum}
Use five distinct qubits $q,r,t,a,b$, with $Q=qrt$, $M=ab$,
$A=\id$, and
\[
 P=\operatorname{CNOT}_{q\to a},\quad
 K=\operatorname{CNOT}_{a\to r},\quad
 B=\operatorname{CNOT}_{q\to b}H_q,\quad
 K'=\operatorname{CNOT}_{b\to t}.                     \tag{RA12}
\]
The Boolean copying gates here are explicit atoms, not full snapshot
translations of source reads. All atoms and keeps are total and allowed.
On the subspace supported on $a=0$, namely pure inputs
$\ket\psi_{qrtb}\otimes\ket0_a$ with factors placed by name and
their density-operator extensions, the inner marked operation
$P^\dagger KP$ is a clean CNOT from $q$ to $r$.
Nevertheless deleting its
inverse changes the outer retained isometry and even the reduced channel
on $qrt$. Here each channel uses the fixed isometry
$E:\cH_{qrt}\to\cH_{qrtab}$,
$E\psi=\psi_{qrt}\otimes\ket{00}_{ab}$ with factors placed by name,
then the respective retained operator, then $\tr_{ab}$; as channels
they are $\rho\mapsto\tr_{ab}(WE\rho E^\dagger W^\dagger)$ and
the same with $\widetilde W$. Let $\Pi_\pm=\proj\pm_t$ with
$\ket\pm=(\ket0\pm\ket1)/\sqrt2$. With identities and tensor
positions implicit, the full retained operators are
\[
 W=\Pi_++\Pi_-X_qX_rZ_b,\qquad
 \widetilde W=\Pi_++\Pi_-X_qX_aZ_b.                  \tag{RA13}
\]
For the input $\ket{00000}_{qrtab}$, their outputs, written in order
$q,r,t,a,b$, are respectively
\[
 \begin{split}
 W\ket{00000}&=
 (\ket{00}_{qr}\ket+_t+\ket{11}_{qr}\ket-_t)
               \ket{00}_{ab}/\sqrt2,\\
 \widetilde W\ket{00000}&=
 (\ket{00}_{qr}\ket+_t\ket0_a+
  \ket{10}_{qr}\ket-_t\ket1_a)\ket0_b/\sqrt2.
 \end{split}                                        \tag{RA14}
\]
In the first case the probability $r=1$ is $1/2$; in the second it is
zero, including after $M$ is traced. The prescribed counts are $11$
and $9$ gates, both on five logical wires with zero implementation
scratch. Omitting the Hadamard in $B$ makes the deletion valid on all
retained inputs. Thus its failure here depends on the intervening
quantum gate, rather than merely on nonzero initial auxiliary values.
\end{counterexample}
\begin{proof}
 \step{resA-ex-inner}{The inner marked operation is
 $R=P^\dagger KP=\operatorname{CNOT}_{a\to r}
                       \operatorname{CNOT}_{q\to r}$.}
 \begin{proof}\pf\ On a basis triple $(q,a,r)=(x,z,y)$, the three
 gates give $(x,z\oplus x,y)$, then
 $(x,z\oplus x,y\oplus z\oplus x)$, then
 $(x,z,y\oplus z\oplus x)$. This is the stated product.
 For $z=0$ it changes only $r$ by XOR with $x$ and restores $a$;
 linear extension gives cleanliness on the entire stated $a=0$
 subspace, including entanglement with untouched factors or a reference.
 \end{proof}
 \step{resA-ex-outer}{The conjugated outer keep is
 $H_0=B^\dagger K'B=\Pi_++\Pi_-X_qZ_b$.}
 \begin{proof}\pf\ $K'=\Pi_++\Pi_-Z_b$, since its target $X_t$
 has eigenvalues $+1,-1$. Conjugating $Z_b$ by
 $\operatorname{CNOT}_{q\to b}$ gives $Z_qZ_b$, as follows by
 multiplying a basis vector by $(-1)^{b\oplus q}$.
 Conjugation by $H_q$ changes $Z_q$ to $X_q$.\end{proof}
 \step{resA-ex-two}{Equation (RA13) holds.}
 \begin{proof}\pf\ The original operator is $R^\dagger H_0R$.
 In the expression for $R$ from \stepref{resA-ex-inner}, its
 $a\to r$ factor commutes with $X_q$, and conjugation by its
 $q\to r$ factor sends $X_q$ to $X_qX_r$ by direct basis action.
 For the modified body $V=BKP$, its outer operation is
 $P^\dagger K^\dagger H_0KP$. The keep $K$ commutes with $X_q$,
 whereas conjugation by $P$ sends $X_q$ to $X_qX_a$. All these
 gates commute with $Z_b$ and $\Pi_\pm$. Substitute into
 \stepref{resA-ex-outer} to obtain (RA13).\end{proof}
 \step{resA-ex-input}{Equation (RA14), the probabilities, and the
 failure of the commutator test follow.}
 \begin{proof}\pf\ Apply (RA13) to the all-zero basis vector, using
 $\Pi_\pm\ket0_t=\ket\pm_t/\sqrt2$ and $Z_b\ket0_b=\ket0_b$.
 The two displayed summands in either output are orthogonal and have
 squared norm $1/2$. In the original output only the second has
 $r=1$; in the modified output both have $r=0$.
 Moreover $PH_0P^\dagger=\Pi_++\Pi_-X_qX_aZ_b\ne H_0$,
 so (RA9) fails.\end{proof}
 \step{resA-ex-noH}{Without $H_q$ the deletion is valid, and the
 counts with $H_q$ are $11$ and $9$.}
 \begin{proof}\pf\ Without $H_q$, the conjugated keep is
 $\Pi_++\Pi_-Z_qZ_b$, which commutes with $P$ because CNOT
 preserves its control's $Z_q$ and does not touch $b,t$.
 Apply (RA9). With $H_q$, the lengths in (RA11) are
 $(a,p,b,k,k')=(0,1,2,1,1)$, giving $11$ and $9$.
 All five wires were explicitly listed and only elementary H and
 CNOT gates occur.\end{proof}
 \qedstep
 \begin{proof}\pf\ By \stepref{resA-ex-inner}--\stepref{resA-ex-noH}.
 \end{proof}
\end{proof}

\subsection{Small exact resource instances}

\begin{example}[Three depths with explicit circuit lists]
\label{ex:resA-counts}
In family N let $P_0=\operatorname{CNOT}_{q\to a}$,
$P_{j+1}=\mathsf M(P_j,\operatorname{CNOT}_{a\to r})$ on
$Q=qr$, $M=a$. Each supplied atom and keep has cost one; there
are three logical wires and no implementation scratch.
In family D let all $b_j=u_j=1$, $f_j(x)=x$, and let evaluator $A_j$
apply CNOT from its input to its zero output, then CNOT from that input
to its zero private bit. Thus $a_j=2$, $k_j=1$, $o_j=5$.
The external input and target are one qubit each. Workspace below means
the allocated sufficient widths (RA7), not an optimal width. All counts
include every gate of the prescribed lists, even cancellable pairs.
\[
\begin{array}{c|rrr|rrrr|rr}
 d& L_N&C_N&n_N&E_D&D_D&W_E&W_D&n_E&n_D\\\hline
 1&2&3 &3& 5 &5 &2&2&4&4\\
 2&3&7 &3&15 &9 &3&4&5&6\\
 3&4&15&3&35&13 &4&6&6&8
\end{array}                                         \tag{RA15}
\]
Here $n_N$ is total logical width of family N, and $n_E,n_D$ add the
two external qubits to family D's auxiliary widths. Serial gate time
is the corresponding gate-count column. Both family D procedures
realize the clean CNOT from their external input to their final target.
\end{example}
\begin{proof}
 \step{resA-table-N}{The three family N rows are correct.}
 \begin{proof}\pf\ There is one original atom and $d$ keeps, so
 $L_N=d+1$. Equation (RA4) with $c_a=c_K=1$ gives
 $C_N=2^{d+1}-1=3,7,15$. The gates use only $q,a,r$ and require
 no scratch; inverse repetitions do not add wires.\end{proof}
 \step{resA-table-D}{The three family D rows are correct.}
 \begin{proof}\pf\ On input $(x,0,0)$ each evaluator yields
 $(x,x,x)$, verifying its classical hypothesis and its cost two.
 A one-bit copy has cost one. Substitute into (RA6)--(RA7):
 $E_D=5(2^d-1)$, $D_D=4d+1$, $W_E=d+1$, $W_D=2d$.
 For $d=1,2,3$ these are exactly the displayed integers.
 Adding the external two wires gives $n_E=d+3$ and $n_D=2d+2$.
 The composite Boolean function is identity, so (RA5) is clean CNOT.
 \end{proof}
 \qedstep
 \begin{proof}\pf\ By \stepref{resA-table-N} and \stepref{resA-table-D}.
 \end{proof}
\end{proof}

\begin{remark}[Scope of the resource comparison]
The exponential lower count here is for a prescribed expansion, not for
quantum or reversible computation as a model. The delayed construction
has a proved retained semantics under its explicit hypotheses, including
its quantum extension (RA8). For arbitrary intervening gates, (RA9) or
(RA10) is a separate obligation at each proposed deletion; (RA12)
shows why cleanliness of an inner mark alone is insufficient. No
pebble-game theorem or unverified external reference carries a proof
step here. General optimal time--space trade-offs and unrestricted
dynamic mark control are not settled by these results.
The evaluator comparison requires exact basis outputs and a preserved
input even in (RA8), includes all evaluator wires in $u_j$, and uses
scratch-free XOR-copy lists. The multiple-deletion corollary is restricted
to base Tier~A terms without menus. The general resource theorem permits
menus only with supplied exact finite-alphabet lists and the once-only
gate/record accounting stated there. All restrictions collected in
Remark~\ref{rem:r3-restrictions} continue to apply.
\end{remark}

\notedsection{Dynamic control (Tier~B)}{sec:tierB}
  {This section studies \earlierforms.}
% Tier~B: a certified bounded semantics, not a source compiler.
\subsection{A bounded reversible interpreter: scope and obligations}

\begin{remark}[Status of the extension]\label{rem:tB-status}\label{target:tierB}
Definition~\ref{def:tierb} admits a static menu only. The following is
an additional, conditional semantics for a bounded dynamic interpreter.
It specifies what an implementation certificate must establish; it does
not assert that every source program, or the existing executor, has such
a certificate. In particular a supplied sequence of unitary matrices
without the journal and routing obligations below is insufficient.
Every dynamic source interpretation and implementation bound below is
conditional on the supplied certificate and, for circuit costs, the
specified implementation. The static clock schedule and the algebraic
disposal test do not construct a dynamic certificate. General dynamic
source rewind, including the unconditional source-level extension in
Remark~\ref{problem:remaining}, remains open.
\end{remark}

\begin{definition}[Finite interpreter state]\label{def:tB-state}
Fix a finite quantum data register $Q$, an ordinary-work register $A$
encoded in qubits, a program $P$ of length $\ell\geq1$, and integers
$T\geq0$, $N\geq1$, $d\geq0$, and $w\geq1$.
All are supplied independently of the unknown input.
The interpreter register is
\[
 M=A\,C\,J\,S\,R\,Z\,\Gamma.
\]
Here $C$ contains the program position, direction/mode and finite
control; $J$ has $N$ journal cells of $w$ bits and a cursor;
$S$ has $d$ mark cells and a stack pointer; $R$ retains return addresses,
visit identifiers, saved cursors and any other routing information;
$Z$ is implementation workspace; and $\Gamma$ is a clock. All pointer,
identifier and counter widths are explicit. One mark cell records a
mark identifier, an activation identifier, a visit identifier, the
journal cursor, and a continuation address. Each field fits its declared
width. Repeated visits use distinct live identifiers or a proved reversible
reuse rule. Integers are finite binary strings; overflow is never silently
ignored.

A journal entry identifies a specified total unitary, its sign (forward
or adjoint), its operands and recording cells, and the old bits needed
to invert an ordinary update. A record of a gate includes its actual
matrix phase. No quantum state is copied into a journal: ordinary
basis information is copied coherently, and gates are specified by
classical descriptors. Descriptors needed later are retained or recovered
by proved reversible operations. Reading $Q$ means controlled reversible
copying into machine bits. Positions, modes, mark identifiers, cursors and
budgets may consequently be in superposition. The controlled-reversible
reading of primitive~10 has the preserved-sector scope of
Lemma~\ref{lem:read-product}; subsequent gates changing a read register
do not give a controlled description of the entire run.

All wires are reserved initially. There is no observed measurement,
drop, allocation, external output observation or physical stopping during
the $T$ ticks. Output tapes, faults and timeout indicators are retained
registers. A claim of successful source behavior requires that capacities,
valid mark identifiers, blank-cell conditions and the supplied horizon
suffice on the entire allowed input subspace. Completion on other states
is part of the semantics, not a claim of successful source execution.
\end{definition}

\begin{definition}[Local rules and their certificate]\label{def:tB-cert}
Write $B=QACJSR$ for the logical state without clock or clean workspace.
The set of local rules is finite. A certificate names a linear initial
subspace $\mathcal S_{\rm cert}\subseteq
\mathcal H_B\otimes\operatorname{span}\{\ket0_\Gamma\}$;
in an implementation its initial
embedding also appends $\ket0_Z$. The clock is omitted in the autonomous
case. The reachable-subspace condition in item 5 applies to this entire
linear space, not merely to a set of basis inputs. If capacity overflow,
an exhausted journal or stack, or an invalid mark identifier is reachable
within $T$ ticks, the instance has no certificate for successful source
execution on that subspace.
A local rule $e$ consists of an input projector $D_e$ on $B$ and a
partial isometry $V_e$ with
\[
 V_e^\dagger V_e=D_e,\qquad V_eV_e^\dagger=E_e.
\]
Input projectors for different rules are orthogonal, and so are their
output projectors $E_e$. These are \emph{separate obligations}.
Rules are selected by computational-basis mode, position, journal and
mark fields, including coherently computed predicates. Predicates that
are not already machine bits must first be computed reversibly.
Their computation and eventual erasure are themselves included in the
rules. Quantum gates can act within a selected sector.

Each rule supplies explicit coordinate identifications of its input and
output subspaces with the same finite payload space $\mathcal L_e$,
\[
 I_e^-:\mathcal L_e\longrightarrow\mathcal H_B,
 \quad I_e^+:\mathcal L_e\longrightarrow\mathcal H_B,
 \quad V_e=I_e^+F_e(I_e^-)^\dagger,                 \tag{B1}
\]
where $I_e^-$ and $I_e^+$ are isometries onto
$\operatorname{ran}D_e$ and $\operatorname{ran}E_e$, respectively, and $F_e$ is
a specified total unitary. The coordinate identifications are explicit
basis encodings, not unspecified identifications of abstract spaces.
They include every live journal cell and every retained old value.
The forward ordinary action of $F_e$ requires the stated blank cells;
on other codes its chosen total completion still defines the operation.
The snapshot tick of Definition~\ref{def:joint-tick} is one possible
$F_e$; its inverse is the stated $D_t^\dagger F_t^{-1}$, including its
sampling and snapshot registers, not just the inverse data gate.

A certificate verifies the following instruction meanings on a specified
input subspace and all its reachable states through tick $T$.
\begin{enumerate}
\item A forward rule performs its stated completed step and records the
information needed for its inverse, including control changes. A mark
pushes the tuple of Definition~\ref{def:tB-state} reversibly.
\item A keep is the total XOR operation of Definition~\ref{def:keep};
its controlling machine fields are preserved during that operation.
The subsequent change to undo mode retains its old control information.
\item An undo rule selects the last entry above the named mark. If that
entry denotes $F_e$ in (B1), the undo rule has isometric coordinate
encodings $\bar I_e^+,\bar I_e^-$ onto its input and output sectors,
respectively, and is exactly
\[
 \bar V_e=\bar I_e^- F_e^\dagger(\bar I_e^+)^\dagger.          \tag{B2}
\]
The bars permit different direction and continuation labels, subject to
the following matching requirement on the entire retained payload.
Write $\bar D_e=\bar V_e^\dagger\bar V_e$ and
$\bar E_e=\bar V_e\bar V_e^\dagger$. There are specified isometries
$\Lambda_e^+:\operatorname{ran}E_e\to\operatorname{ran}\bar D_e$
and $\Lambda_e^-:\operatorname{ran}\bar E_e\to\operatorname{ran}D_e$,
onto those sectors, such that
\[
 \bar I_e^+=\Lambda_e^+ I_e^+,
 \qquad I_e^-=\Lambda_e^-\bar I_e^-
 \quad\hbox{on all of }\mathcal L_e.                 \tag{B2a}
\]
Every routing or mode operation taking a forward output sector to its
undo input sector realizes $\Lambda_e^+$ on the whole sector; the
corresponding routing after undo realizes $\Lambda_e^-$ on the whole
undo output sector. These are payload-coordinate identities also for
vectors produced by a keep, not only for ordinary forward outputs.
Consecutive recorded entries must admit a common retained payload space
after identity padding. If the specified isometry $J_i$ routes the output of entry $i$ to
the input of entry $i+1$, the specified encodings obey
$J_i I_i^+=I_{i+1}^-$ on that whole space; the reverse routing obeys
$J_i^\dagger I_{i+1}^-=I_i^+$. All these routings retain every live
record and are among the certified operations. Their costs are charged.
The full operator $F_e^\dagger$,
including its journal and control action, is required. A signed adjoint
entry is undone with $F_e$. An inner rewind is represented by its entire
reversible operation, including its keep and routing, or by its sequence
of elementary signed entries; a compound entry is not a unit-cost gate.
\item Undo continues one elementary entry at a time until the recorded
cursor of the coherently named mark is reached. Equality is a reversible
Boolean test. Changing mode and installing a continuation preserve all
information needed to undo those changes. Popped records are retained
until they have been reversibly cleared. Repeated use of a mark and undo
of an earlier rewind obey these same rules; neither operation licenses
re-execution of different, newly selected instructions.
\item For every $0\leq t\leq T$, every vector in
$W^t\mathcal S_{\rm cert}$ lies in the sum of the asserted rule
domains, and each of its sector components satisfies the assertion for
that rule. Here $W$ is the logical tick specified below; in (B3) use
the input projectors $\proj c_\Gamma\otimes D_{c,e}$ for the
clock-$c$ rules, and in the autonomous case use $D_e$.
In an implementation whose single tick $W$ has the supplied clean
implementation $\widetilde W$ of Theorem~\ref{thm:tB-sound}, all
statements are lifted with the fixed zero workspace factor
$\ket0_Z$, which $\widetilde W$ restores at every tick. Thus this condition includes superpositions over
different rule sectors, every post-keep state and every partial undo.
Formula (B2) and the matching identities (B2a) must hold
on that whole domain, not just on states obtained with an empty keep.
All rules have orthogonal ranges as stipulated above. A stopped logical
output may idle only through a reversible rule, retaining any timing
information it requires.
\end{enumerate}
The identifications in (B1)--(B2a), the domain/range assertions, and
reachable-subspace assertions are mathematical proof obligations for an
instance. They are not established merely by naming an instruction.
For an efficient circuit claim, a total clean circuit realizing all
rules and their complement must additionally be supplied.
\end{definition}

\begin{remark}[Why the range obligation matters]\label{rem:tB-range}
For a map specified only on blank records, the allowed input and its
image need not be invariant under keep. The inverse on that image then
has no prescribed action after keep. A total $F_e$ fixes the algebraic
inverse, but a state outside a claimed ordinary sector need not have the
intended cursor, mark or continuation. Extending $F_e$ alone does not
prove the reachable-subspace obligation. Likewise, retaining an extra
copy of a control decision can prevent erasure: reversing the data gates
while leaving that copy is not reversal of the full forward operation.
\end{remark}

\begin{lemma}[Exact extension criterion]\label{lem:tB-extension}
Let $\mathcal H$ be finite dimensional and let $D_1,\ldots,D_k$ be
orthogonal projectors with pairwise orthogonal ranges. For each $i$ let
$V_i$ be an isometry from $D_i\mathcal H$ into $\mathcal H$, extended
by zero on the orthogonal complement. There exists a unitary $U$ on
$\mathcal H$ with $UD_i=V_i$ for every $i$ if and only if the ranges
of the $V_i$ are pairwise orthogonal. The assertion does not impose a
gate-alphabet or efficiency condition on $U$.
\end{lemma}
\begin{proof}
 \step{tB-ext-necessary}{\assume{Such a $U$ exists.}
 \prove{The ranges of the $V_i$ are pairwise orthogonal.}}
 \begin{proof}\pf\ For $x\in D_i\mathcal H$ and
 $y\in D_j\mathcal H$ with $i\ne j$,
 $\langle V_ix,V_jy\rangle=\langle Ux,Uy\rangle
 =\langle x,y\rangle=0$.\end{proof}
 \step{tB-ext-sufficient}{\assume{The ranges are pairwise orthogonal.}
 \prove{Such a $U$ exists.}}
 \begin{proof}\pf\ The direct sum of the $V_i$ is an isometry on
 $\sum_iD_i\mathcal H$. Its input and output dimensions are equal.
 Specify orthonormal bases of their equal-dimensional orthogonal
 complements and map one basis to the other. Combining the two maps
 gives a unitary with the required restrictions.\end{proof}
 \qedstep
 \begin{proof}\pf\ The two implications establish the equivalence.
 \end{proof}
\end{proof}

\begin{definition}[One fixed tick and bounded idle]\label{def:tB-tick}
Set $D=\sum_eD_e$, $E=\sum_eE_e$ and $V=\sum_eV_e$.
A completion is a \emph{specified} unitary isomorphism
$V_\perp:(\id-D)\mathcal H_B\to(\id-E)\mathcal H_B$,
extended by zero elsewhere. Put $U=V+V_\perp$.
In the circuit version this completion must have its supplied clean
implementation; finite-dimensional existence is no efficiency claim.
One may supply such rules and a completion $U_c$ for each clock value
$c$, for example to reserve clock-indexed recording cells and bounded
idle intervals.
In that case $D_{c,e}$ and $V_{c,e}$ denote the input projector and
partial isometry of the clock-$c$ copy of rule $e$, and $U_c$ is the
specified completion formed from them as above.
Choose $L_{\Gamma}=2^b>T$, let $\Gamma$ have $b$ bits initially
zero, and define the single fixed logical tick
\[
 W=\sum_{c=0}^{L_{\Gamma}-1}\ket{c+1\bmod L_{\Gamma}}\bra c_\Gamma\otimes U_c.
                                                               \tag{B3}
\]
Identical $U_c=U$ give the time-independent microcode special case.
A finite autonomous microcode that is already a total unitary can
instead omit the separate clock ($b=0$); a stated horizon still applies.
In (B3) idle may have $U_c=\id$, but the clock still advances bijectively.
The logical run is exactly $W^T$. When a clean implementation of $W$
is supplied, its implementation workspace $Z$ returns to zero each tick.
There is no amplitude-dependent choice of the number of physical ticks.
Logical completion times, mark visits and undo lengths may differ.
The clock has no wrap during the promised interval; outside it the
cyclic continuation is still part of the unitary specification.
\end{definition}

\begin{theorem}[Conditional dynamic soundness and exact undo]
\label{thm:tB-sound}
Assume Definitions~\ref{def:tB-state}--\ref{def:tB-tick}, including all
certificate obligations, specified completions and the full reachable
input subspace $\mathcal S_{\rm cert}$. In particular assume the
matching-encoding identities (B2a), their whole-sector routing
realizations, and the consecutive-entry identifications of
Definition~\ref{def:tB-cert}, including after every keep.
Then $W^T$ is unitary on all retained logical registers.
If the single tick $W$ itself has a supplied clean implementation
$\widetilde W$ on $B\Gamma Z$ in the sense of
Definition~\ref{def:scratch}, namely
$\widetilde W(\psi\otimes\ket0_Z)=W\psi\otimes\ket0_Z$ for every
$\psi$, then $T$ concatenations give a clean implementation with
\[
 (\widetilde W)^T(\psi\otimes\ket0_Z)
       =W^T\psi\otimes\ket0_Z\quad\hbox{for every }\psi.
\]
No clean implementation of $W$ is inferred here from clean
implementations of the separate $U_c$ alone.
On the certified subspace it executes the specified dynamic instructions.
Each undo together with its specified label routing is the exact adjoint
of its corresponding completed forward operation: on all of
$\mathcal L_e$,
\[
 \Lambda_e^-\bar V_e\Lambda_e^+ I_e^+
       = I_e^-F_e^\dagger.                              \tag{B2b}
\]
This includes arbitrary intervening gates, completed inner rewinds and
unequal ticks of arrival at an undo instruction, provided the certificate
obligations hold after those operations. The statement does not assert
cleanliness or a certificate for every source program. General dynamic
source rewind remains open (Open problem~\ref{prob:tB-compiler}).
\end{theorem}
\begin{proof}
 \step{tB-unit-partial}{$V^\dagger V=D$ and $VV^\dagger=E$.}
 \begin{proof}\pf\ For $e\ne f$, orthogonality of output sectors gives
 $V_e^\dagger V_f=0$, and orthogonality of input sectors gives
 $V_eV_f^\dagger=0$. Sum the defining diagonal identities.\end{proof}
 \step{tB-unit-complete}{The completion exists and $U$ is unitary.}
 \begin{proof}\pf\ The preceding step gives equal ranks of $D,E$.
 Their orthogonal complements have equal finite dimension. Choose and
 specify orthonormal bases on those complements and map one to the other.
 Orthogonality eliminates cross terms, so
 $U^\dagger U=D+(\id-D)=\id$ and $UU^\dagger=\id$.
 This argument alone gives no efficient finite-alphabet synthesis.
 Indeed, when the complementary sector is nonzero, exact synthesis of
every prescribed completion over a fixed finite alphabet is
impossible: each finite circuit over that alphabet, with any finite
number of clean workspace wires, is specified by a finite string of
gate names and wire numerals, so these circuits and their exact
logical restrictions form a countable set; whereas, for fixed unit
vectors in the two complementary sectors, mapping the first to
$e^{i\theta}$ times the second and extending by one fixed basis map
elsewhere gives uncountably many distinct unitary completions of the
same $V$ as $\theta$ ranges over $[0,2\pi)$.
 \end{proof}
 \step{tB-unit-clock}{$W$ and $W^T$ are unitary.}
 \begin{proof}\pf\ Distinct summands of (B3) have orthogonal clock
 inputs and outputs. Therefore $W^\dagger W=WW^\dagger=\id$.
 A finite product of this unitary is unitary. For $b=0$ use the preceding
 step directly. If $\widetilde W$ is supplied as hypothesized, induction
 on $t$ gives $(\widetilde W)^t(\psi\ket0_Z)=W^t\psi\ket0_Z$:
 the preceding tick restores zero workspace, so the defining equality
 applies to $W^t\psi$ at the next tick. This also covers $t=0$.
 This product assertion is the elementary part of soundness.\end{proof}
 \step{tB-unit-inverse}{Undo together with its specified label routing
 reverses the entire local completed step.}
 \begin{proof}
  \step{tB-unit-cancel}{Equation (B2b) holds on all of $\mathcal L_e$;
  the routed undo equals $V_e^\dagger$ on $\operatorname{ran}E_e$.}
  \begin{proof}\pf\ For any $v\in\mathcal L_e$, (B2a) and (B2) give
  \[
  \begin{split}
   \Lambda_e^-\bar V_e\Lambda_e^+ I_e^+v
    &=\Lambda_e^-\bar V_e\bar I_e^+v\\
    &=\Lambda_e^-\bar I_e^-F_e^\dagger v
      =I_e^-F_e^\dagger v=V_e^\dagger I_e^+v.
  \end{split}
  \]
  The second equality uses
  $(\bar I_e^+)^\dagger\bar I_e^+=\id$; the last uses (B1).
  Since $I_e^+$ is onto, this proves the sector equality. Taking
  $v=F_eu$ gives $I_e^-u$, the original forward input. The same
  equality holds for every other $v$, including the post-keep vector,
  by the whole-sector routing hypothesis.\end{proof}
  \step{tB-unit-sequence}{For a recorded sequence
  $F_1,\ldots,F_k$ on a common retained payload, with identity padding
  and matching routing identifications included, the inverse payload action is
  $F_1^\dagger\cdots F_k^\dagger$.}
  \begin{proof}\pf\ For $k=0$ the product is identity. For $k>0$,
  write the vector at the last forward output as $I_k^+v$.
  Step~\stepref{tB-unit-cancel} takes it to
  $I_k^-F_k^\dagger v$. If $k>1$, the certified reverse routing
  $J_{k-1}^\dagger$ takes this to $I_{k-1}^+F_k^\dagger v$.
  Induction on the remaining $k-1$ entries gives
  $I_1^-F_1^\dagger\cdots F_k^\dagger v$.
  Forward routing $J_i I_i^+=I_{i+1}^-$ gives the forward payload
  product $F_k\cdots F_1$ by the same induction. Thus the two
  payload products are adjoints. The reachable-subspace obligation
  justifies every sector component used, including after keeps;
  the identities do not assume $v$ has a blank-record origin.
  An inner completed operation $A^\dagger K A$ has adjoint
  $A^\dagger K^\dagger A$ by multiplication of adjoints. The same
  identity applies to its retained control operations.\end{proof}
  \qedstep
  \begin{proof}\pf\ These identities are operator identities on the
  certified sectors. Inserting gates or a keep changes the vector to
  which the inverse is applied, not the inverse matrix. Returning the
  original vector additionally requires undoing all inserted operations;
  $F^\dagger K F$ need not be identity.\end{proof}
 \end{proof}
 \step{tB-unit-dynamic}{Different arrival times require no observation.}
 \begin{proof}\pf\ At any tick decompose the vector by the orthogonal
 rule-domain projectors. Apply the specified linear map to each component
 and sum. Induction over the fixed $T$ ticks uses the reachable-domain
 obligation at each stage. The selected mark and cursor comparisons are
 machine-register operations. No common number of logical undo entries
 is assumed.\end{proof}
 \qedstep
 \begin{proof}\pf\ The preceding steps prove the stated algebraic and
 certified operational assertions. Verification of domains, disjoint
 ranges, (B2) and the whole-sector matching in (B2a), rather than
 multiplication of unitaries, is the
 substantive implementation obligation.\end{proof}
\end{proof}

\begin{remark}[The physical tick is not reversed by advancing its clock]
An undo-mode application of $W$ is not the global operator $W^\dagger$:
the clock in (B3) still advances. The exact-inverse assertion concerns
the completed logical operation and its explicit encodings, including
all of its destroyed-information records. Reversing a physical tick,
clock included, requires $W^\dagger$. Omitting this distinction would
make the exact-inverse claim false even for a classical counter.
\end{remark}

\subsection{The static specialization and the disposal criterion}

\begin{theorem}[Embedding Tier~A, with its machine registers retained]
\label{thm:tB-A}
Let $P_A$ be a finite fixed-space Tier~A term satisfying all applicable
hypotheses of Theorem~\ref{thm:sA-sound}, with
its supplied expansion into logical atoms, keeps and adjoint logical
atoms $A_0,\ldots,A_{n-1}$ in execution order, each a total unitary on
$L_A=QM_A$, with $A_{n-1}\cdots A_0=\mathcal U_{P_A}$.
This is the logical expansion on $L_A$, not its elementary clean circuit
on $L_AZ$. It includes the original occurrence-labelled records and
any specified continuation atoms. Choose $L_{\Gamma}=2^b>T\geq n$ and a fresh
counter $\Gamma$ initially zero, and let $M_B$ be any disjoint auxiliary
register in a fixed unit vector $\ket{m_B}$. Take
$U_c=A_c\otimes\id_{M_B}$ for $0\leq c<n$ and
$U_c=\id_{L_AM_B}$ for $n\leq c<L_{\Gamma}$.
Then this is a static specialization of the tick form (B3), with no
dynamic mark or undo instruction executed, and for every $\psi_{L_A}$,
\[
 W^T(\psi_{L_A}\otimes\ket{m_B}\ket0_\Gamma)
 =\mathcal U_{P_A}\psi_{L_A}\otimes\ket{m_B}\ket T_\Gamma.
                                                               \tag{B4}
\]
The fixed list already specifies every action; no separate mutable mark
stack is used. This theorem supplies no certificate under
Definition~\ref{def:tB-cert}. Any dynamic source interpretation would
additionally require that certificate, whose general construction remains
open (Open problem~\ref{prob:tB-compiler}).
Here a clean implementation concerns only $Z$ as in
Definition~\ref{def:scratch}; logical cleanliness would restore $M_A$
exactly, in the sense of Definition~\ref{def:clean}.
Classical control alone does \emph{not} imply $M_A$ is clean, nor does it
make an arbitrary implementation's extra records independent of $Q$.
Equation (B4) is an existence statement for this implementation.
\end{theorem}
\begin{proof}
 \step{tB-A-word}{The expanded list has product $\mathcal U_{P_A}$.}
 \begin{proof}\pf\ The logical recursion in Definition~\ref{def:tiera}
 and Theorem~\ref{thm:sA-sound} expands identity as an empty list, an
 atom as itself, and a sequence by concatenation. A marked body expands
 as its logical list, its keep, and the reversed adjoint logical list
 using the same record names. Structural induction gives the displayed
 product on all of $\mathcal H_{L_A}$. This never substitutes the
 action of an elementary implementation on nonzero workspace.
 Distinct counter inputs and outputs make the summands of (B3)
 orthogonal. No dynamic certificate is asserted for this schedule.
 \end{proof}
 \step{tB-A-counter}{After $t\leq T$ ticks the clock equals $t$ and
 the logical action is $U_{t-1}\cdots U_0$.}
 \begin{proof}\pf\ Induct on $t$ using (B3), starting with the empty
 product. The assumption $T<L_{\Gamma}$ excludes wrap.\end{proof}
 \qedstep
 \begin{proof}\pf\ At $T$ the last $T-n$ factors are identity.
 Combine the preceding steps. Nothing erases $M_A$; its cleanliness is
 the separate condition in Theorem~\ref{thm:sA-clean}.\end{proof}
\end{proof}

\begin{theorem}[Exact synchronization and disposal]\label{thm:tB-clean}
Let $W$ be a finite unitary interpreter tick, let $T\geq0$, initialize
all of $M$ in a fixed pure state $m_0$, and allow every input on $Q$.
For a dynamic source application, $W$ and its interpretation must be
supplied under the certificate hypotheses of Theorem~\ref{thm:tB-sound}.
To claim that interpretation for every input considered here, under
the identification $QM=B\Gamma Z$ of
Definitions~\ref{def:tB-state}--\ref{def:tB-cert} with factors placed
by name, the fixed vector $m_0$ must have clock factor
$\ket0_\Gamma$ and workspace factor $\ket0_Z$, the tick must be the
logical tick $W$ of Theorem~\ref{thm:tB-sound} (respectively its
supplied clean implementation $\widetilde W$ when $Z$ is retained
in $M$), and $E_0\mathcal H_Q$, with the $Z$ factor suppressed as
in Definition~\ref{def:tB-cert}, must be contained in
$\mathcal S_{\rm cert}$.
The algebraic test below applies to any such finite unitary, without
constructing that certificate; general dynamic source rewind remains
open (Open problem~\ref{prob:tB-compiler}).
Set $E_0\psi=\psi\otimes m_0$ and $V=W^TE_0$.
For a specified unit vector $\eta\in\mathcal H_M$ the following are
equivalent:
\begin{enumerate}
\item There is a unitary $A$ on $Q$ such that $V\psi=A\psi\otimes\eta$
for every $\psi$.
\item $(\id_Q\otimes(\id_M-\proj\eta))V=0$.
\item The final $M$ marginal is $\proj\eta$ for every input density
operator on $Q$.
\end{enumerate}
Moreover the reduced channel $\rho\mapsto\tr_M(V\rho V^\dagger)$ is
unitary if and only if these conditions hold for some fixed pure $\eta$.
They continue to hold with an arbitrary untouched finite reference.
This disposal test is weaker than logical reset: $\eta$ need not equal
$m_0$. A clean implementation restores its workspace $Z$, but does not
by that fact restore the other logical interpreter factors.
Thus clock, position, journal, stack \emph{and all other retained
interpreter information} must collectively factor in a fixed pure state
if disposal is to implement a unitary on arbitrary $Q$.
Equal completion times alone are not a sufficient condition: the fixed
finite schedule of Theorem~\ref{thm:sA-four}, with keep target initially
zero, has one endpoint time while its final machine vectors
$\ket{00}$ and $\ket{10}$ remain correlated with the data, so collective
factorization fails there.
For a preserved computational-basis input with an orthonormal output
family, if
\[
 V\ket x=\ket{f_x}\ket{m_x},
\]
then an input matrix element $\ket x\bra y$ is sent after disposal to
$\langle m_y|m_x\rangle\ket{f_x}\bra{f_y}$. In particular orthogonal
clock/position records erase that off-diagonal element.
\end{theorem}
\begin{proof}
 \step{tB-clean-12}{Conditions 1 and 2 are equivalent.}
 \begin{proof}\pf\ Condition 1 implies 2 by projection. Under 2 set
 $A=(\id_Q\otimes\bra\eta)V$. Then $V\psi=A\psi\otimes\eta$
 for every $\psi$, and
 $A^\dagger A=V^\dagger V=\id$. Equal finite input and output
 dimensions make $A$ unitary.\end{proof}
 \step{tB-clean-23}{Conditions 2 and 3 are equivalent.}
 \begin{proof}\pf\ Condition 2 implies the stated pure marginal.
 Conversely apply 3 to every pure input. The squared norm of its
 component orthogonal to $Q\otimes\eta$ is the expectation of
 $\id\otimes(\id-\proj\eta)$, hence zero.\end{proof}
 \step{tB-clean-channel}{A unitary reduced channel implies a fixed pure
 environment vector.}
 \begin{proof}\pf\ If the reduced channel is conjugation by $A$,
 each pure output $V\ket i$ has $Q$ marginal
 $A\ket i\bra iA^\dagger$. The squared norm of its component
 orthogonal to $\operatorname{span}\{A\ket i\}\otimes\mathcal H_M$
 is the expectation of the complementary projection, hence zero.
 Therefore $V\ket i=A\ket i\otimes\eta_i$ for a unit vector
 $\eta_i$, on an orthonormal input basis. Density matrices span the
 full operator space: diagonal matrix units are pure projections and
 off-diagonal units are linear combinations of projections onto
 $(\ket i\pm\ket j)/\sqrt2$ and $(\ket i\pm i\ket j)/\sqrt2$.
 Channel equality thus holds on $\ket i\bra j$ and gives
 \[
 \langle\eta_j|\eta_i\rangle A\ket i\bra jA^\dagger
       =A\ket i\bra jA^\dagger.
 \]
 The rank-one operator on the right is nonzero, so the scalar is one.
 Hence $\|\eta_i-\eta_j\|^2
 =2-2\operatorname{Re}\langle\eta_j|\eta_i\rangle=0$.
 If $\dim Q=1$ there is only one vector. Conversely a common
 vector gives conjugation by $A$ on tracing.\end{proof}
 \step{tB-clean-elements}{The stated matrix-element formula holds.}
 \begin{proof}\pf\ Take the partial trace of
 $\ket{f_x}\bra{f_y}\otimes\ket{m_x}\bra{m_y}$.\end{proof}
 \qedstep
 \begin{proof}\pf\ Combine the steps. Tensoring the operator equality
 in condition 1 with the reference identity proves the reference claim.
 The equal-time failure follows from the nonzero, orthogonal machine
 terms and rank-two output proved in Theorem~\ref{thm:sA-four}.
 The condition is collective: constant individual mixed marginals do
 not establish factorization of the joint machine register.\end{proof}
\end{proof}

\begin{counterexample}[Exact undo at unequal ticks]\label{ex:tB-delay}
There are four qubits, in order $q,r,p_1,p_0$. The data are $Q=qr$;
$M=p_1p_0$ is a compact encoding of position, mode, the one-entry journal
and elapsed idle time. There is a single fixed mark, hence its identifier
and stack need zero additional bits. No separate clock or workspace is
used. Define $p=2p_1+p_0$, and the following complete permutations and
gate bits:
\[
\begin{array}{c|rrrr|rrrr}
 &\pi(0)&\pi(1)&\pi(2)&\pi(3)&a(0)&a(1)&a(2)&a(3)\\\hline
 q=0&1&2&3&0&1&1&0&0\\
 q=1&1&3&0&2&1&0&0&1
\end{array}
\]
The single fixed tick is
\[
 W\ket{q,r,p}=\ket{q,r\oplus a_q(p),\pi_q(p)}.          \tag{B5}
\]
At $p=0$ it performs a forward $X_r$ and records it in the position
encoding. For $q=0$ it undoes that entry at the second tick ($p=1$);
for $q=1$ it waits at $p=1$ and undoes at the third tick ($p=3$).
Completion is $p=2$ in either case. The transition $2\to3$ for $q=0$
is an idle tick recording elapsed completion time. Remaining transitions
complete the finite cycles; only $T=3$ ticks are claimed as execution.
For arbitrary amplitudes $\alpha,\beta$ and an arbitrary unit vector
$\chi\in\mathcal H_r$,
\[
 W^3\bigl((\alpha\ket0+\beta\ket1)_q\chi_r\ket0_p\bigr)
 =\alpha\ket0_q\chi_r\ket3_p+
   \beta\ket1_q\chi_r\ket2_p.                         \tag{B6}
\]
For $q=\ket+$ disposal of $p$ yields
$\tfrac12\id_q\otimes\proj\chi_r$, although each executed undo is
exact and $r$ is restored. The two completion ticks are 2 and 3.
\end{counterexample}
\begin{proof}
 \step{tB-delay-unit}{Equation (B5) is a permutation unitary.}
 \begin{proof}\pf\ For each unchanged $q$, $\pi_q$ is a permutation.
 From its output recover $p=\pi_q^{-1}(p')$ and then
 $r=r'\oplus a_q(p)$. This gives a unique inverse on all 16 basis states.
 Thus rule domains and ranges are disjoint.\end{proof}
 \step{tB-delay-evolve}{Equation (B6) holds.}
 \begin{proof}\pf\ For $q=0$, positions are $0,1,2,3$ and data gates
 are $X,X,\id$. For $q=1$, positions are $0,1,3,2$ and gates are
 $X,\id,X$. Both products are identity. Apply linearity. The inverse
 data action is $X^\dagger=X$. The compressed register $p$ is used
 only through the displayed total unitary (B5); no exhibition of this
 example as a certificate instance of Definition~\ref{def:tB-cert}
 is required for the disposal calculation.\end{proof}
 \qedstep
 \begin{proof}\pf\ The final machine states are orthogonal, so
 Theorem~\ref{thm:tB-clean} removes the cross terms. For $q=\ket+$
 the reduced $q$ state has purity $1/2$.\end{proof}
\end{proof}

\begin{lemma}[The global clock and absorbing completion]
Assume the clock initialization and tick (B3). Also let $U$ be any
unitary and $u,v$ any vectors in its Hilbert space.
For every initial state $\Phi\otimes\ket0_\Gamma$, the global clock
of (B3) is always the fixed basis
state $\ket T$ at the endpoint. It cannot itself cause the loss in (B6).
A data-dependent local elapsed-time register or position can do so.
Furthermore, if $Uu=v$ and $Uv=v$, then $u=v$.
\end{lemma}
\begin{proof}
 \step{tB-absorbing-clock}{At tick $t\leq T$, the clock factors as
 $\ket t$ from all other registers.}
 \begin{proof}\pf\ Induction in (B3) advances every input vector's
 clock by one, independent of its other registers. There is no wrap
 since $T<L_{\Gamma}$. A fixed pure factor has no correlations to discard.
 \end{proof}
 \step{tB-absorbing-injective}{$Uu=Uv$ implies $u=v$.}
 \begin{proof}\pf\ Multiply by $U^\dagger$.\end{proof}
 \qedstep
 \begin{proof}\pf\ Apply the preceding two steps; the local-position
 obstruction is Example~\ref{ex:tB-delay}.\end{proof}
\end{proof}

\begin{remark}[Checked halting and synchronization context]\label{rem:tB-halting}
Bernstein--Vazirani define exact halting by no final configurations before
$T$ and only final configurations at $T$
\cite[Definition 3.5.1, PDF p.~14]{BernsteinVazirani1997}.
Their synchronization theorem applies to polynomial-time classical
$f$ with output length depending only on input length. It supplies a
stationary normal-form reversible machine computing $(x,f(x))$ in time
depending only on input length
\cite[Theorem 4.1.3, PDF p.~17]{BernsteinVazirani1997}.
Its phrase ``whose running time depends only on the length of $x$''
does not supply a synchronization theorem for arbitrary coherent
mark selection after gates.

Linden--Popescu's specified auxiliary evolution retains elapsed stopping
time; they write ``ancilla contains a record of the time since the
computation halted''
\cite[PDF p.~2, equations (7)--(11) and text after (9)]{LindenPopescu1998}.
Example~\ref{ex:tB-delay} gives a finite bounded instance of the relevant
entanglement obstruction. We use neither a general nonexistence claim
nor a claim that unequal durations always destroy coherence.

Ozawa proves equality of monitored and unmonitored output probabilities
under his assumptions (I)--(III), including invariant halted tape output
\cite[PDF p.~2, equation (11); proof pp.~2--4,
equations (12)--(14) and following]{Ozawa1998}.
The phrase ``no more changes the halt qubit nor the tape string'' is
assumption (II), equation (6), on PDF p.~2. Equality of that output
distribution is distinct from unitary preservation of arbitrary quantum
output and does not contradict (B6).
Shi questions the composition argument for a particular universal-machine
proposal \cite[PDF pp.~7--8, Section 5, equations (12)--(14)]{Shi2002}.
This is contextual, not a theorem precluding bounded reversible circuits.
No halting paper carries a proof step of the results above.
\end{remark}

\subsection{Circuit costs and the unresolved compilation problem}

\begin{definition}[Explicit circuit cost parameters]\label{def:tB-cost}
For the following implementation bound, assume a supplied reversible
microcode circuit for (B3). It fetches a $w$-bit instruction from a
fixed classical table of length $\ell$ at an $a$-bit quantum address,
where $a=\lceil\log_2(\ell+1)\rceil$; out-of-range fetch returns a
specified default. This is a coherent circuit lookup, not an oracle
or assumed quantum random-access memory. It uses at most $\kappa$
journal accesses and at most $\kappa$ mark-stack accesses per tick,
where $\kappa$ is an absolute constant independent of
$\ell,a,w,N,d,u,v,\beta,g,b,h,z,q$ and $T$; it is part of the
constant hidden in $O(\cdot)$ in (B7). The cell width is at most $w$,
capacities $N$ and $d$, and binary cursor widths
$u=\lceil\log_2(N+1)\rceil$, $v=\lceil\log_2(d+1)\rceil$.
Take $w$ large enough for every descriptor and mark tuple.

Let $g$ be the sum of lengths of the supplied elementary action circuits
that are coherently selected during one tick, including all wire-address
selection, controlled gates, keeps and inverses. Let $\beta$ bound the
sum of sizes, including all input and output ports, of the other Boolean
circuits \emph{and their supplied inverse
routing computations}. These circuits must implement the certificate,
including behavior on invalid codes and completion. A table depending
on the clock is charged at its actual additional size in $\beta$ or $g$.
Let $h$ be all other logical machine bits and let $z$ bound reusable clean
workspace for the supplied component circuits. No small value of
$g,\beta,h,z$ follows from the existence of an abstract completion.
\end{definition}

\begin{theorem}[Conditional implementation cost]\label{thm:tB-cost}
Assume a certificate satisfying Definition~\ref{def:tB-cert}, its
specified tick and completion, and all supplied clean-circuit and
access hypotheses of Definition~\ref{def:tB-cost}. The following bound
is conditional on those data; it supplies no general dynamic source
compiler, which remains open (Open problem~\ref{prob:tB-compiler}).
A scan implementation of one tick has
\[
 G_{\rm tick}=O\bigl(\ell(a+w)+N(u+w)+d(v+w)
                  +\beta+g+b+w+1\bigr).                    \tag{B7}
\]
A sufficient total qubit count, including the input $q=|Q|$, is
\[
 q+h+(N+d)w+a+u+v+b+
 O(z+a+u+v+w+b+1).                                        \tag{B8}
\]
Here $h$ includes $A$, direction flags and retained routing records;
workspace for Boolean evaluation must already be included in $z$.
$T$ ticks cost $O(TG_{\rm tick})$ gates with the same wires.
The circuit always consists of exactly $T$ copies of the tick circuit.
A component in an idle rule sector still passes through that circuit;
no count of padded ticks for an individual sector component is subtracted
from the gate total. For a general input vector different sector
components may be idle at different ticks, so no single nonidle count
need exist. Separately, the static embedding of Theorem~\ref{thm:tB-A}
has $n$ scheduled logical-factor ticks and $T-n$ padding ticks, compared
with the direct Tier~A list of the $n$ logical factors $A_c$.
That schedule's factors are not counted as unit-cost elementary gates;
applying (B7) to any implementation requires all its circuit hypotheses.
A quantum-position interpreter therefore has
fetch and routing overhead even with classical control.
\end{theorem}
\begin{proof}
 \step{tB-cost-fetch}{Instruction lookup costs $O(\ell(a+w))$.}
 \begin{proof}\pf\ For each classical address $i$, reversibly compute
 the equality flag $[p=i]$ using $O(a)$ elementary Boolean gates and
 clean scratch, XOR the fixed $w$-bit word controlled by that flag,
 and undo the flag. Only one address can match. If no address matches,
 one further $O(a)$ range test and an $O(w)$ controlled XOR install
 the specified default word, after which the flag is uncomputed.
 Since $\ell\geq1$, this is absorbed by $O(\ell(a+w))$. Apply
 Lemma~\ref{lem:synthesis}; reversing the circuit erases a fetched
 descriptor when its inputs have been restored. A descriptor cannot
 be erased using a modified address without that restoration.\end{proof}
 \step{tB-cost-memory}{The stated journal and stack terms suffice.}
 \begin{proof}\pf\ Scan all $N$ journal cells, compute each cursor
 equality flag in $O(u)$ gates, and perform the required controlled
 bitwise XORs or swaps in $O(w)$ gates. Reverse the flag computation.
 The analogous scan of $d$ mark cells costs $O(d(v+w))$.
 The absolute access bound $\kappa$ absorbs repetitions of these scans.
 Cursor increments, decrements and comparisons use linear-size
 reversible binary arithmetic with retained carries uncomputed.
 The supplied microcode guarantees the required blank cells or retains
 their old values; a scan does not implement irreversible deletion.\end{proof}
 \step{tB-cost-rest}{The other terms and qubit count suffice.}
 \begin{proof}\pf\ Coherently control the supplied gate lists and
 evaluate the remaining Boolean circuits, charging $g$ and $\beta$.
 For the clock, evaluate the successor with a linear-size ripple
 circuit into a blank $b$-bit register, swap registers, and erase the
 old value by evaluating the predecessor into it. Boolean evaluation
 uses Lemma~\ref{lem:synthesis}; this costs $O(b+1)$ gates and clean
 workspace. Store the named cells and cursors, the $h$ other
 logical bits, a fetched word and a reusable set of equality/carry
 scratch bits. The clean component-workspace bound is $z$.
 This gives (B7)--(B8).\end{proof}
 \qedstep
 \begin{proof}\pf\ Repeat the same circuit $T$ times and reuse its
 zero-output workspace. For a superposition, there is
 one circuit consisting of $T$ ticks, not a sum of execution times over basis
 inputs.\end{proof}
\end{proof}

\begin{proposition}[Nesting costs in this semantics]\label{prop:tB-nesting}
Assume elementary signed journal entries, the certificate and explicit
tick-cost hypotheses of Theorem~\ref{thm:tB-cost}, and a fixed horizon
$T$. The bound contains no factor depending on nesting depth except
through its parameters $T,\ell,a,N,u,d,v,\beta,g,b,w,h,z$;
in particular it does not bound $T$ by a polynomial in the source
description. This conditional count supplies no general dynamic source
rewind construction, which remains open.
For the static family $P_0=a$ and
$P_{j+1}=\mathsf{mark}\{P_j\}\mathsf{rewind}\ K$, if the prescribed
expansion executes every elementary body gate and its adjoint and makes
no cancellations, its number $n_j$ of payload gate applications is
\[
 n_{j+1}=2n_j+k,\qquad n_j=2^jn_0+(2^j-1)k,                \tag{B9}
\]
where $n_0$ and $k$ are the lengths of the atom and keep circuits.
Any interpreter executing this list with at most one payload gate per
tick has $T\geq n_j$. This is a tight count for that expansion, not a
lower bound for all equivalent implementations. With $N=O(T)$ journal
capacity, (B7) gives the total scan bound
\[
 O\bigl(T^2(u+w)+
 T(\ell(a+w)+d(v+w)+\beta+g+b+w+1)\bigr),
\]
so no automatic linear-in-$T$ bound is obtained. Keeping a journal
does not itself remove (B9).
\end{proposition}
\begin{proof}
 \step{tB-nesting-global}{The total circuit has $T$ copies of its tick.}
 \begin{proof}\pf\ This is its definition and
 Theorem~\ref{thm:tB-cost}. Nesting affects the parameters and $T$,
 not an additional multiplication of this count.\end{proof}
 \step{tB-nesting-count}{Equation (B9) holds.}
 \begin{proof}\pf\ A marked expansion prints its body once, its
 $k$ keep gates once, and the adjoint body once. This gives the
 recurrence. Induction on $j$ gives the displayed solution.
 The circuit itself attains that count; serial execution cannot use
 fewer ticks than the number of gates it is required to execute.\end{proof}
 \qedstep
 \begin{proof}\pf\ Combine the steps and substitute $N=O(T)$ in
 (B7). No optimal-synthesis or general delayed-erasure theorem follows.
\end{proof}
\end{proof}

\subsection{Small exact instances and outstanding obligations}

\begin{example}[Finite matrices for independent checking]\label{ex:tB-table}
In each example below $p$ is a binary encoded cyclic instruction
counter, initially zero. It also encodes the fixed mode, journal
occupancy and fixed single-mark control; no unlisted auxiliary wire is
used. For a gate list $(G_0,\ldots,G_{L_{\Gamma}-1})$, notation
$[G_0,\ldots,G_{L_{\Gamma}-1}]$ means the complete unitary
$\sum_{p=0}^{L_{\Gamma}-1}\ket{p+1\bmod L_{\Gamma}}\bra p\otimes G_p$.
A fixed mark identifier needs no separate storage. These compressed
examples test algebraic and timing claims, not a variable-size stack
compiler.
\[
\begin{array}{c|c|c|c|l}
\text{case}&Q&M&T&\text{tick and endpoint from zero machine state}\\\hline
\mathrm{I}&q&p\ (1\text{ bit})&2&[X,X]:\ \psi\ket0\mapsto\psi\ket0\\
\mathrm{II}&q&p\ (2\text{ bits})&3&[H,X,H,\id]:\ \psi\ket0\mapsto Z\psi\ket3\\
\mathrm{III}&qr&p\ (2\text{ bits})&3&\text{(B5)--(B6); unequal completion ticks}\\
\mathrm{IV}&q&p\ (3\text{ bits})&7&[H,X,H,X,H,X,H,\id]:\
 \psi\ket0\mapsto -X\psi\ket7
\end{array}
\]
Case I uses two qubits, II three, and III--IV four. Case II is
$H^\dagger XH=Z$; case IV is an outer rewind with inner completed
operation $HXH=Z$ and outer keep $X$.
The minus sign in IV is part of the specified matrix and must not be
removed if the operation is later controlled.
\end{example}
\begin{proof}
 \step{tB-table-unit}{Each displayed tick is unitary.}
 \begin{proof}\pf\ Each counter input has a unique counter output,
 and its attached data gate is unitary. Case III was proved in
 Example~\ref{ex:tB-delay}.\end{proof}
 \step{tB-table-products}{The endpoints are as stated.}
 \begin{proof}\pf\ Apply the attached gates in counter order.
 $X^2=\id$, $HXH=Z$, and $ZXZ=-X$ by the two-by-two matrices.
 The counter ends at the indicated value.\end{proof}
 \qedstep
 \begin{proof}\pf\ The steps establish the table for every data
 input, including inputs entangled with an untouched reference.\end{proof}
\end{proof}

\begin{problem}[A compiler and a source simulation theorem]
\label{prob:tB-compiler}
For general bounded source programs with repeated dynamic mark visits,
construct the explicit coordinate encodings, local rules and efficient
completion in Definition~\ref{def:tB-cert}, and prove the reachable-domain
and disjoint-range obligations. This is not proved here. In particular,
after a keep changes a previously read input, newly executing the source
code can choose different instructions or fail to reach the saved visit.
It is not thereby the inverse of the recorded computation. A successful
translation must recover the exact old operation including its control,
or specify a different total completion and state the resulting semantic
change. Inspections of nonzero amplitudes, discarding a clock and taking
only a maximal deadline are not reversible Boolean operations on a
computational-basis state without additional records and a proof.
No equivalence with the executor's scheduling, rollback or budget rules
is established by this section.
\end{problem}

\begin{problem}[Synchronization, resources and unrestricted execution]
A source-level sufficient condition, proved for dynamic marks and arbitrary
intervening gates, which guarantees the factorization in
Theorem~\ref{thm:tB-clean}, remains open. Equal padded durations do not
clear the journal or mark stack. Extending the bounded construction to
an unknown unbounded runtime is not accomplished by the finite cyclic
clock. No efficient finite-alphabet synthesis of an arbitrary abstract
completion is asserted; exact synthesis of a prescribed arbitrary
completion need not exist over the finite alphabet, so a circuit claim
must supply its completion and circuit together, as in
Definition~\ref{def:tB-cert}. Finally, a general transformation
removing exponential repeated undo while retaining arbitrary quantum
keeps, with a proved time--space trade-off, is not supplied. The exact
count (B9) is not an impossibility result for such a transformation.
\end{problem}

\notedsection{Open problems}{sec:open-problems}
  {Some of these open problems concern \earlierforms.}
The bounded dynamic results are conditional but proved: a supplied certificate gives a unitary retained interpreter run and exact undo under matching encodings (Theorem~\ref{thm:tB-sound}); Tier~A is included as a static clock schedule, without constructing a dynamic certificate (Theorem~\ref{thm:tB-A}). The exact disposal criterion requires the entire interpreter register to factor in one fixed pure state for every data input from a fixed pure initial machine state (Theorem~\ref{thm:tB-clean}). Counterexample~\ref{ex:tB-delay} restores the target at unequal completion ticks but leaves a timing record whose disposal decoheres the data. Theorem~\ref{thm:tB-cost} charges instruction-table, journal and stack scans and all supplied action, Boolean and clock circuits at every tick; its $T$-tick total is conditional on the stated circuits and access bounds. These results supply neither a general source compiler nor a certificate for an unspecified implementation.

The open-point and restriction remarks~\ref{problem:remaining}, \ref{rem:r3-restrictions} and \ref{rem:sk} delimit the following questions.
\begin{enumerate}
\item Compiling general dynamic source mark/rewind into the certified bounded interpreter remains open (Open problem~\ref{prob:tB-compiler}). The missing construction must handle mutable mark stacks, data-dependent repeated visits, general recursive calls and an inverse stopped at a coherent mark identifier, and prove the local encodings, matching routing, disjoint ranges and reachable-domain obligations after every keep. A supplied unitary schedule or static finite menu does not establish those obligations. No source-level sufficient condition for the collective factorization of Theorem~\ref{thm:tB-clean} under arbitrary intervening gates is supplied; equal padded durations do not clear the journal or stack. The obstruction
is familiar from quantum Turing machines whose computations stop after
different numbers of steps: retained records can distinguish the stopping
times and prevent the intended interference when those records are
discarded. Tier~A fixes its schedule in the syntax; Tier~B supplies an
explicit bound and padding but must still satisfy the collective disposal
test. Counterexample~\ref{ex:tB-delay} exhibits the timing-record
obstruction, and Remark~\ref{rem:tB-halting} states the checked literature
context and its limitations.
\item Theorem~\ref{thm:resA-resources} settles resource counts for the prescribed unsimplified expansion: each enclosing mark doubles body work, and this dependence is tight. These lower counts apply only to that expansion and give no lower bound for optimal equivalent circuits. Theorem~\ref{thm:resA-delayed} proves delayed erasure for the stated local evaluators with preserved basis inputs, exact basis outputs and a product private factor on each basis input, scratch-free XOR copies, and no intervening changes to retained values. It does not give a transformation for arbitrary intervening quantum gates. Keeping an inner record and deleting its immediate inverse preserves the full retained operator exactly when the commutator condition of Theorem~\ref{thm:resA-deletion} holds; on a specified input subspace that theorem supplies the weaker exact test. The five-qubit Counterexample~\ref{ex:resA-quantum} shows that inner cleanliness alone is insufficient. General optimal time--space trade-offs and unrestricted delayed-erasure transformations remain open. The conditional multiple-deletion result, Corollary~\ref{cor:resA-flat}, gives no intermediate time--space continuum.
\item Source allocation inside a mark remains excluded. A reservation translation would require a proof that the reserved factors remain zero before their first use; equality of initialized embeddings does not provide an inverse allocation on arbitrary states. The bounded deferred-observation lemma does not admit observed measurement within marked reversible computations.
\item The four-qubit calculation uses explicit Boolean copying atoms. Its equality is not asserted for the full-snapshot translation of two source read instructions, whose additional configuration records require a separate calculation.
\item General actual-phase approximation certificates remain an assumption for additional matrices beyond the effective determinant-one one-qubit rotations with a fixed number of controls in Theorem~\ref{eqv:rotations}. That theorem charges the supplied matrix-entry procedure and uses exhaustive certified search; an efficient general angle evaluator or polynomial classical synthesis time is not proved here.
\item No theorem here infers a sufficient input-independent termination horizon from the instruction list, establishes an unbounded stopping simulation, or proves correctness of an unspecified software executor. Within the settled model, position restoration and auxiliary reset remain additional conditions rather than consequences of unitary soundness.
\end{enumerate}

\notedsection{Summary of results, restrictions and corrections}{sec:summary}
  {Part of this section summarizes results on \earlierforms.}
\subsection{Proved conclusions and pending results}
The settled results are the two bounded simulations (Theorems~\ref{eqv:forward} and \ref{eqv:reverse}), the conditional approximation theorem and its effective rotation construction (Theorems~\ref{eqv:approx} and \ref{eqv:rotations}), structured unitary soundness (Theorem~\ref{thm:sA-sound}), exact and recursive restoration tests (Theorems~\ref{thm:sA-clean} and \ref{thm:sA-recursive}), and the sufficient product-output restoration condition with its explicit limitations (Theorems~\ref{thm:sA-classical} and \ref{thm:sA-four}). Universal auxiliary restoration is false, as witnessed by the four-qubit circuit.
Counterexample~\ref{ex:sA-rec-limits} limits the reset tests themselves:
containing subspaces lose necessity, final cleanliness does not imply
cleanliness of inner marks, marginal restoration does not license
deletion, and a unitary reduced channel need not reset the auxiliary
register. Proposition~\ref{eqv:restrictions} and
Counterexamples~\ref{eqv:exact-obstruction} and
\ref{eqv:phase-obstruction} prove that the capacity, stopping, exact-angle
and phase restrictions cannot be suppressed;
Corollary~\ref{eqv:adaptive-circuits} packages the two simulations for
bounded adaptive device data. The resource and conditional bounded-interpreter results below are proved; compiling general dynamic source mark/rewind into a certified interpreter remains open. Open questions and exclusions are collected in Section~\ref{sec:open-problems}.

Primitive~10 has the controlled-reversible meaning of
Lemma~\ref{lem:read-product} when the subsequent operations preserve the
read register's computational-basis sectors. A gate changing that register
removes this guarantee; the full joint unitary and its retained records
then specify the semantics.

Theorem~\ref{thm:resA-resources} gives exact gate and serial-time counts for the prescribed expansion, sufficient qubit counts, and an explicit compiler bound. A mark doubles its body's work and adds its keep cost. Repeating this nesting gives a tight exponential count for that expansion. This is not a lower bound for optimal equivalent circuits: with an empty keep the expanded circuit can already denote identity. Repeated inverse uses reuse the same named records; extra gates alone do not require extra logical wires.

Theorem~\ref{thm:resA-delayed} proves a restricted alternative: retain all evaluator outputs and private wires until the final copy, then invert the evaluators in reverse order. For uniformly positive bounded evaluator costs and bounded copy costs, its serial time is linear in the number of levels, while the prescribed eager construction is exponential. The stated sufficient space allocations need not be strictly separated. More generally, keeping an inner record until the outer rewind and deleting its immediate inverse preserves the whole retained operator exactly when the inner forward operation commutes with the outer keep conjugated by the later computation (Theorem~\ref{thm:resA-deletion}). A weaker exact test applies on a specified input embedding. A clean inner computation alone does not suffice: the five-qubit Counterexample~\ref{ex:resA-quantum} changes the probability of $r=1$ from $1/2$ to zero, even after the machine register is traced. Certified repeated deletions give linear time in the once-counted body size only under Corollary~\ref{cor:resA-flat}'s tests and restrictions.

For Tier~B, Theorem~\ref{thm:tB-sound} proves that a bounded interpreter run is unitary on all retained logical registers when the supplied certificate verifies the stated rules, completions and reached subspaces. Under the whole-sector matching encodings, each undo with its label routing is the exact adjoint of its completed forward operation, including its retained records. A supplied clean implementation of the whole tick returns workspace to zero at every tick. Theorem~\ref{thm:tB-A} realizes Tier~A as a static specialization of the clock schedule with all its machine registers retained; it constructs no dynamic certificate and implies no machine reset.

Theorem~\ref{thm:tB-clean} gives an exact synchronization and disposal test. For all data inputs and a fixed pure initial machine state, tracing the interpreter register leaves a unitary data channel exactly when the entire final interpreter register factors in one fixed pure state independent of the input. This includes position, clock, journal, stack and routing records. The fixed final vector need not be the initial one: disposal and logical reset remain distinct. Equal durations alone do not establish the test. In the four-qubit Counterexample~\ref{ex:tB-delay}, exact undo restores the target after two or three ticks according to the data value. At the fixed three-tick endpoint the local timing record remains correlated with the data; disposing of it sends the control qubit's $\ket+$ state to the equal mixture of $\ket0$ and $\ket1$. This is a finite matrix example, not an exhibited dynamic-certificate instance.

Under the supplied-circuit and constant-access hypotheses of Definition~\ref{def:tB-cost}, Theorem~\ref{thm:tB-cost} gives the per-tick gate cost
\[
 O\bigl(\ell(a+w)+N(u+w)+d(v+w)+\beta+g+b+w+1\bigr).
\]
Here $\ell,N,d$ are instruction-table, journal and mark-stack capacities; $a,u,v$ are their address widths, $w$ the cell width, $\beta$ the supplied Boolean-circuit cost including inverse routing, $g$ the supplied action-circuit cost, and $b$ the clock width. The first three terms charge linear scans; there is no assumed memory-access oracle. Exactly $T$ copies of the tick circuit give the total gate cost; the sufficient wire allocation includes all retained records and clean workspace. Proposition~\ref{prop:tB-nesting} leaves nesting dependence in these parameters and the horizon: a journal alone does not remove the prescribed expansion's doubling, and capacity $N=O(T)$ gives a quadratic-in-$T$ scan term with its width factors retained.

The open step is compiling general dynamic source mark/rewind into an interpreter meeting the certificate, particularly data-dependent repeated mark visits and an inverse stopped at a coherent mark identifier (Open problem~\ref{prob:tB-compiler}). No such general compiler, source-level synchronization condition under arbitrary gates, unbounded-runtime extension or executor-equivalence theorem is proved.

\subsection{Every restriction of the admitted program class}
\begin{enumerate}
\item All instances and register lists are finite and fixed independently of the unknown quantum input. Uniform families are generated by one classical Turing machine. Polynomial uniformity additionally requires polynomial numerical bounds and polynomial generation time. Compilation counts charge the supplied descriptions, binary addresses and explicit output lists; numerical bounds in those estimates are supplied in unary. A binary horizon still costs its numerical number of execution steps.
\item Each device supplies the complete field layout, validity test, initial configuration, canonical total decoder, total update and total output circuits, including invalid-code outputs and fault behavior. All live-wire checks and flag-preservation rules hold on every configuration code. Unspecified arithmetic, addressing, overflow, division and call conventions are not identified with any particular source language. No efficiency of an arbitrary supplied Boolean function is presumed.
\item The initial configuration is valid and its live flags identify exactly the input wires. Allocations select never-used zero factors. Names are never reused, dropped names remain inaccessible, and no later operation acts on them. There is no device storage-release instruction. The output quantum list is fixed and live at final time for every reachable outcome list; the entire complementary retained register is traced.
\item Observed protocols for primitives 1--9 contain no marks and exclude coherent-read mode on every code. They use an input-independent horizon of at least one tick, terminal padding and a separate final timeout snapshot within that horizon. Each forward tick has distinct blank sample and configuration records. Measurements are computational-basis observations; classical output is a specified final tape or register. Equivalence concerns the final instrument on the same input and output lists, with arbitrary references; it preserves neither write times nor stopping times and describes no intermediate interaction.
\item A reversible Turing-machine transition is partial and injective under the stated tape convention. Finite total unitary use fixes its completion. Agreement with the original machine requires each forward transition to stay within the valid finite window and restricted transition domain. Efficient device simulation uses supplied circuits and claims no efficient arbitrary completion.
\item A Tier~A fixed-space marked term uses total unitary atoms, total Boolean keeps and a fixed retained register list. Marks and their pairings are static. Calls or loops containing marks need a supplied fixed finite syntactic expansion; general recursion and data-dependent numbers of mark visits are excluded. Bounded loops without marks may instead use scheduled ticks. Coherent tick decoders exclude measurement and allocation on every code. Continuation addresses must fit their fields.
\item No allocation, drop, observed measurement, external console observation or actual stop occurs inside a marked body. An unobserved console tape and bounded idle statuses may be updated. Logical records and implementation workspace are allocated before each enclosing run. Source allocation inside a mark needs a separate zero-input reservation proof and is not admitted automatically.
\item Atom occurrences have distinct retained record names; reversal reuses those same names in reverse order. Any factor needed by an enclosing inverse is retained. Nonblank records still have a total unitary action, but ordinary forward instruction semantics requires blankness. Continuation is a separate recorded update; position restoration and machine reset are additional properties.
\item Each keep uses a supplied total Boolean circuit on the whole specified machine list and a fixed list of distinct target wires disjoint from that machine list. Targets lie in the enclosing context; there is no computed-address exception. Keep entries may overlap targets. A target used by the body is allowed, with no automatic restoration guarantee.
\item Exact finite-alphabet implementation claims require supplied exact clean circuits for the atoms. Implementation workspace starts and ends at zero; no equality on its nonzero states is claimed. All physical qubits, including logical records, are retained before output partial trace. An ideal-angle atom has no automatic exact finite-alphabet implementation.
\item A controlled menu has a finite fixed collection of base structured terms, a common payload list, and a disjoint selector with sufficient capacity. Payloads contain no further menus, preserve the selector, and supply actual-phase exact finite-alphabet circuit lists for the implementation claim. Records are distinct across entries and occurrences and are retained, with the selector, for enclosing inverses. This does not define dynamic source mark stacks.
\item Restoration tests apply on the specified input subspace. Exact prefix images give necessity and sufficiency for actual inputs; containing subspaces give sufficiency alone. Blank-record reuse also requires an incoming zero condition. Logical reset points treat menu atoms as indivisible; boundaries within their entries or within physical implementations require the specified full circuit, its own initial subspace, all workspace, and the selector where present. Marginal restoration or a successful reset does not permit deletion of a register still required by an inverse. Final cleanliness does not imply cleanliness of every inner mark.
\item The sufficient product-output restoration theorem requires a total unitary forward operation independent of the keep target. For each basis input its forward data and machine outputs must be a product with a basis-valued machine output, with all keep controls included in that machine list. It need not preserve the input basis value. Removing the basis-output or unread-target requirement invalidates the universal claim; no necessity of these sufficient hypotheses for individual clean operations is asserted. The four-qubit witness uses explicit copying atoms rather than full-snapshot source-read instructions.
\item Approximation uses rational requests, counts every expanded factor occurrence including inverses and controls, bounds actual matrices with fixed uniform length constants, and reverses the same approximating circuit for inverse uses. General certificate-generation time is charged separately and need not be polynomial. The discharged case is determinant-one one-qubit matrices with supplied dyadic entry approximations of stated precision and cost, with a fixed number of controls. Arbitrary scalar phases, other additional matrices and growing control counts are outside that synthesis theorem. No matrix-entry procedure is derived merely from an angle description.
\item All resource gate counts use the prescribed unsimplified lists over the adjoint-closed alphabet $\mathcal G^{\pm}$ and serial execution. Sufficient wire allocations are not minimum-space lower bounds. The exponential lower counts apply only to this expansion; there is no lower bound here for optimal equivalent circuits, parallel depth or optimal time with a space restriction. These resource results do not imply logical-record reset. Zero-initialisation
of non-input wires, if charged at one tick per wire, adds
$N(P)-q_{\rm in}$ ticks, and preparing a nonzero initial configuration
is additional (Theorem~\ref{thm:resA-resources}).
\item The snapshot wire and scratch bounds use one shared configuration width, sample width and pair of device-circuit bounds across all ticks. They remain parametric in declared extra logical wires, arbitrary supplied-atom scratch and keep-circuit sizes. No efficiency bound for an arbitrary total completion is supplied. Payload record occurrences and explicit selector computations are charged once in the named terms; an externally supplied selector contributes its retained width, and its earlier computation cost is outside the displayed allocation. The compiler bound assumes an uncompressed prefix syntax encoding and already supplied delimited gate lists; constructing those lists is excluded.
\item The delayed-evaluator comparison requires total local evaluators, zero output and private inputs, a preserved basis input, an exact basis output and a product private factor for each basis input. That private factor may be a quantum state. All evaluator implementation wires are included in the declared private widths. XOR-copy circuits use no scratch and act only on the output and a disjoint target. Final targets never become evaluator inputs; earlier private wires remain untouched until their inverses, and no intervening operation changes retained values. Equality is of initialized retained isometries, up to auxiliary padding and factor permutation, and need not hold on nonzero auxiliary inputs. The linear versus exponential time comparison also requires the stated uniform cost bounds.
\item Deleting an inner inverse requires the exact commutator condition for equality on every retained input, or the stated weaker condition for the fixed input embedding. It does not itself reset a record or justify reuse of logical zero cells. The repeated-deletion corollary excludes menu atoms, requires a valid test at every deletion and requires that these deletions reduce the body to one visit per specified occurrence. It retains every logical wire and keeps the supplied exact clean local circuits. It provides no general transformation for arbitrary intervening quantum gates and no intermediate or optimal time--space trade-off. The corollary as stated has a single outermost mark.
\item Tier~B fixes finite capacities, field widths, program and an input-independent horizon $T$ before the unknown input; all wires are reserved initially. Its interpretation requires a supplied certificate for a linear initialized subspace and every reached sector component, including post-keep and partial-undo states. Rules have separately verified orthogonal input and output sectors. Capacity failure or an invalid mark identifier on the reached subspace defeats certification of successful source execution. No observation, allocation, drop or physical stop occurs during the run.
\item Every dynamic completed step has explicit forward and undo payload encodings. The matching identities (B2a), whole-sector routing realizations and consecutive-entry identifications hold on the entire retained payload, including post-keep vectors. All control changes, visit identifiers, popped records and earlier rewinds are retained or reversibly cleared with their exact action. Repeated visits require distinct live identifiers or a proved reversible reuse rule; merely naming mark and rewind proves none of these obligations. Constructing this interface for general dynamic source code remains open.
\item Tier~B circuit claims supply the specified completion and a clean implementation of the whole tick; abstract finite-dimensional extension gives neither an efficient circuit nor exact finite-alphabet synthesis of an arbitrary prescribed completion. The per-tick bound assumes the stated microcode, an absolute constant number of journal/stack accesses, Boolean-circuit sizes including ports and inverse routing, and supplied action, other-register and scratch costs. The run uses exactly $T$ tick circuits even when some sector components are idle. No polynomial source bound for $T$ follows. The static Tier~A embedding is a logical-factor schedule and supplies no dynamic certificate or unit-cost interpretation of a compound factor.
\item The Tier~B disposal equivalence tests every data input from a fixed pure initial interpreter state and requires collective pure factorization. A source interpretation additionally needs the specified zero clock and workspace factors and certified initial-subspace containment. Equal completion times or constant individual mixed marginals do not suffice. The unequal-duration example demonstrates failure for its displayed total unitary; it asserts neither that every unequal-duration run decoheres nor that the example supplies a general source translation.
\end{enumerate}

\subsection{Every correction needed for primitives 1--9}
\begin{enumerate}
\item (primitives 1--6) Integer arithmetic, comparison, memory access, conditional jumps, calls, returns, stopping and console output need complete classical functions and finite capacities. Growing uniform families replace one fixed finite device when comparison with unbounded computation is intended. Word arithmetic is the supplied bounded-width function; larger integers require charged multiword storage and work. Console output is a tape with specified overflow and final extraction.
\item (all primitives) The instruction names do not determine behaviour on
illegal codes. Each device must supply total decoder, update and output
circuits in the canonical layout, including a specified fault configuration
and the output on every invalid encoding (Definition~\ref{def:device}).
\item (primitive 5; all primitives) A horizon, padded terminal states and a scheduled timeout are extra data. The nine instruction names alone imply no termination theorem. Truncating an unbounded execution and adding timeout can change its instrument.
\item (primitive 9) Later quantum choices may depend on observed outcomes. Retained measurement copies, total classical updates and final environment trace define this adaptivity precisely. The bounded equivalence corollary (Corollary~\ref{eqv:adaptive-circuits}) takes an already supplied device presentation; it does not silently compile an undefined controller notation.
\item (primitive 7) Fresh allocation means a zero factor that has never been used. Drop is simulated by retaining an inaccessible factor and tracing it at output; storage reclamation and name reuse are not supplied by that simulation.
\item (primitive 8) Exact gate claims use the actual $H,T,\operatorname{CNOT}$ alphabet. $T^\dagger=T^7$ removes the inverse letter with constant overhead. Exact arbitrary-angle equivalence is false: the stated ideal rotation gives probability $1/3$, unavailable to the bounded finite-alphabet construction with basis initialization, deterministic classical data and unconditional outputs. Postselection and conditional renormalization are excluded from that obstruction.
\item (primitive 8, optional part) Optional rotations require an effective supply of matrix-entry approximations and a separate approximation theorem. Error is measured on actual matrices, including control-dependent phases; discarding a scalar phase before adding a control is invalid. The general theorem assumes certificates. They are supplied here only for the stated determinant-one one-qubit case and bounded controls, using an inverse-closed determinant-one set whose generated subgroup is dense. The word-length bound is polylogarithmic in inverse error, but the proved certificate search time is exponential in that length bound and includes entry-generation and full rational-request costs.
\end{enumerate}


\begin{thebibliography}{99}
\raggedright

\bibitem{Bennett1973}
C.~H.~Bennett,
\emph{Logical Reversibility of Computation},
IBM Journal of Research and Development \textbf{17}(6) (1973), 525--532.
\href{https://doi.org/10.1147/rd.176.0525}{doi:10.1147/rd.176.0525}.

\bibitem{Lamport2012}
L.~Lamport,
\emph{How to write a 21st century proof},
Journal of Fixed Point Theory and Applications \textbf{11} (2012), 43--63.
\href{https://doi.org/10.1007/s11784-012-0071-6}{doi:10.1007/s11784-012-0071-6}.

\bibitem{BoykinEtAl1999}
P.~O.~Boykin, T.~Mor, M.~Pulver, V.~Roychowdhury, and F.~Vatan,
\emph{On Universal and Fault-Tolerant Quantum Computing},
arXiv:quant-ph/9906054 (1999), preprint.
\href{https://arxiv.org/abs/quant-ph/9906054}{arXiv:quant-ph/9906054}.

\bibitem{DawsonNielsen2006}
C.~M.~Dawson and M.~A.~Nielsen,
\emph{The Solovay-Kitaev algorithm},
Quantum Information and Computation \textbf{6}(1) (2006), 81--95.
\href{https://doi.org/10.26421/QIC6.1-6}{doi:10.26421/QIC6.1-6}.

\bibitem{BernsteinVazirani1997}
E.~Bernstein and U.~Vazirani,
\emph{Quantum Complexity Theory},
SIAM Journal on Computing \textbf{26}(5) (1997), 1411--1473.
\href{https://doi.org/10.1137/S0097539796300921}{doi:10.1137/S0097539796300921}.

\bibitem{LindenPopescu1998}
N.~Linden and S.~Popescu,
\emph{The Halting Problem for Quantum Computers},
arXiv:quant-ph/9806054v2 (1998), preprint.
\href{https://arxiv.org/abs/quant-ph/9806054}{arXiv:quant-ph/9806054}.

\bibitem{Ozawa1998}
M.~Ozawa,
\emph{Quantum Nondemolition Monitoring of Universal Quantum Computers},
Physical Review Letters \textbf{80}(3) (1998), 631--634.
\href{https://doi.org/10.1103/PhysRevLett.80.631}{doi:10.1103/PhysRevLett.80.631}.

\bibitem{Shi2002}
Y.~Shi,
\emph{Remarks on universal quantum computer},
Physics Letters A \textbf{293}(5--6) (2002), 277--282.
\href{https://doi.org/10.1016/S0375-9601(02)00015-4}{doi:10.1016/S0375-9601(02)00015-4}.
\end{thebibliography}

\end{document}
